% Leçon 17 — Expression orale : protection de l’environnement — Chinois 1ère A — Cours signé OnBuch+
% Programme officiel MINESEC. Compiler avec : tectonic ZHA17-oral-protection-de-lenvironnement.tex
\newcommand{\DOCMATIERE}{Chinois}
\newcommand{\DOCNIVEAU}{1ère A}
\newcommand{\DOCMODULE}{Module 2 : Environnement et santé}
\newcommand{\DOCLECON}{ZHA17}
\newcommand{\DOCTITRE}{Leçon 17 — Expression orale : protection de l’environnement}
% OnBuch+ — préambule « cours signés » ENRICHI (compilé avec tectonic / XeLaTeX)
% Système visuel « Pop » d'OnBuch+ : fond crème (jamais blanc), bordures encre
% épaisses, ombres « dures » décalées, titres Archivo Black en capitales.
% Version MATHÉMATIQUES : graphes (pgfplots/tikz), boîtes pédagogiques
% (définition, propriété, méthode, attention, le savais-tu, exercice résolu,
% exercice, TP, bilan), page de garde et signature OnBuch+ ancrée en bas.
%
% Macros attendues dans main.tex AVANT \input{preamble} :
%   \DOCMATIERE  \DOCNIVEAU  \DOCMODULE  \DOCLECON (numéro)  \DOCTITRE
\documentclass[11pt]{article}

\usepackage{fontspec}
\usepackage[french]{babel} % français = langue principale ; le chinois via \zh{..} / textebox (xeCJK)
\usepackage{xeCJK}
\usepackage{pifont} % \ding{51} \ding{55} utilisés par les modèles
\usepackage[a4paper,margin=2cm,top=3.4cm,bottom=2.8cm,headheight=1.4cm,footskip=1.3cm]{geometry}
\usepackage{amsmath,amssymb,amsfonts}
\usepackage{mathtools} % \xmapsto, \coloneqq... souvent utilisés par les modèles
\usepackage{cancel} % simplifications barrées (bilans de forces)
% \abs{z} (module, valeur absolue) et \norm{u} (norme), utilisés par les modèles
\providecommand{\abs}{}\let\abs\relax\DeclarePairedDelimiter{\abs}{\lvert}{\rvert}
\providecommand{\norm}{}\let\norm\relax\DeclarePairedDelimiter{\norm}{\lVert}{\rVert}
\usepackage[table]{xcolor} % [table] : fournit aussi \rowcolors (alternance de lignes)
\usepackage{pagecolor}
\usepackage{tikz}
\usetikzlibrary{arrows.meta,positioning,calc,shapes.geometric,shapes.misc,shapes.arrows,decorations.pathmorphing,decorations.pathreplacing,decorations.markings,patterns,shadows,mindmap,trees,fit,backgrounds,matrix,intersections,angles,quotes}
\usepackage{pgfplots}
\usepackage[version=4]{mhchem} % équations nucléaires \ce{^{235}_{92}U}
\usepackage[european,siunitx]{circuitikz} % schémas électriques
\pgfplotsset{compat=1.18}
\usepgfplotslibrary{fillbetween,groupplots} % aires entre courbes (calcul intégral)
\usepackage{enumitem}
\usepackage{fancyhdr}
% [most] charge toutes les bibliothèques tcolorbox (listings, minted, raster,
% documentation...) et gonfle inutilement la mémoire TeX sur un document long
% (~300 boîtes) : on ne charge que ce qui est réellement utilisé.
\usepackage[skins,breakable]{tcolorbox}
\usepackage{graphicx}
\usepackage{colortbl}
\usepackage{booktabs}
\usepackage{tabularx}
\usepackage{diagbox} % case d'en-tête diagonale des tableaux à double entrée
% Colonnes extensibles centrée / gauche / droite (C, L, R), souvent utilisées par les modèles
\newcolumntype{C}{>{\centering\arraybackslash}X}
\newcolumntype{L}{>{\raggedright\arraybackslash}X}
\newcolumntype{R}{>{\raggedleft\arraybackslash}X}
\usepackage{array}
\usepackage{multicol}
\usepackage{multirow}
\usepackage{siunitx}
% parse-numbers=false : accepte aussi \SI{2,5 \times 10^{-3}}{...}
\sisetup{locale=FR,output-decimal-marker={,},group-minimum-digits=4,parse-numbers=false}
\DeclareSIUnit\ohm{\ensuremath{\symup{\Omega}}} % U+2126 absent de la police
\DeclareSIUnit\division{div} % réglages d'oscilloscope (ms/div, V/div)
\DeclareSIUnit\tr{tr} % tours (vitesse de rotation, ex. \SI{1500}{\tr\per\minute})
\DeclareSIUnit\molar{M} % concentration molaire (ex. \SI{3}{\molar}, \SI{1}{\micro\molar})
\DeclareSIUnit\an{an} % année (le modèle confond parfois avec \year, un compteur TeX) — ex. \SI{5}{\kilo\meter\per\an}
\DeclareSIUnit\FCFA{FCFA} % franc CFA (ex. \SI{500}{\FCFA\per\kilo\gram})
\DeclareSIUnit\degreeBrix{\textdegree Brix} % teneur en sucre d'un sirop/jus (réfractométrie)
\DeclareSIUnit\month{mois} % le modèle utilise parfois \month, un compteur TeX, comme unité de durée
\usepackage{qrcode}
% \degree : commande fréquemment utilisée par les modèles pour un angle ou une
% température en dehors de \SI{}{} (non fournie par les paquets chargés).
\providecommand{\degree}{^{\circ}}
% \qed (fin de démonstration), utilisé par les modèles sans amsthm
\providecommand{\qed}{\hfill$\square$}
% \barre : double barre des valeurs interdites dans un tableau de variations (tkz-tab)
\providecommand{\barre}{\ensuremath{\|}}
% environnement proof (amsthm non chargé) utilisé par les modèles
\ifdefined\proof\else\newenvironment{proof}[1][Démonstration]{\par\noindent\textit{#1.}\ }{\qed\par}\fi
\usepackage{lastpage}
\usepackage{hyperref}
\hypersetup{hidelinks}

% ============================================================
% Palette « Pop » OnBuch+ — mêmes hex que la classe P (plus_ui.dart)
% ============================================================
\definecolor{poporange}{HTML}{F59321}
\definecolor{popdark}{HTML}{E07A0C}
\definecolor{popcream}{HTML}{FFF6E8}
\definecolor{popink}{HTML}{1C1714}
\definecolor{poppurple}{HTML}{7A5AE0}
\definecolor{popgreen}{HTML}{2E9E6A}
\definecolor{poppink}{HTML}{E0457B}
\definecolor{popblue}{HTML}{2F7DE1}
\definecolor{popgold}{HTML}{C98A12}
\definecolor{popmuted}{HTML}{8A7F76}
\definecolor{popline}{HTML}{EADFD2}
\definecolor{popgray}{HTML}{8A7F76}
\colorlet{popgrey}{popgray}
% Alias tolérants (le modèle nomme parfois les couleurs différemment)
\colorlet{popwhite}{white}
\colorlet{popblack}{popink}
\colorlet{popred}{poppink}
\colorlet{popyellow}{popgold}
% Teintes claires (fonds de tableaux/encadrés) — le modèle nomme parfois
% les couleurs avec un suffixe L (ex. popblueL) sans les avoir définies.
\colorlet{poporangeL}{poporange!20}
\colorlet{popdarkL}{popdark!20}
\colorlet{poppurpleL}{poppurple!20}
\colorlet{popgreenL}{popgreen!20}
\colorlet{poppinkL}{poppink!20}
\colorlet{popblueL}{popblue!20}
\colorlet{popgoldL}{popgold!20}
\colorlet{popredL}{popred!20}
\colorlet{popgrayL}{popgray!20}
% Teintes claires pour fonds de tableaux / zones
\colorlet{popblueL}{popblue!10!white}
\colorlet{poporangeL}{poporange!14!white}
\colorlet{popgreenL}{popgreen!12!white}
\colorlet{poppurpleL}{poppurple!12!white}
\colorlet{poppinkL}{poppink!12!white}
\colorlet{popgoldL}{popgold!14!white}

\pagecolor{popcream}

% ============================================================
% Typographie
% ============================================================
\defaultfontfeatures{Ligatures=TeX}
\setmainfont{PlusJakartaSans-Medium.ttf}[Path=../../../fonts/,BoldFont=PlusJakartaSans-Bold.ttf]
% Symboles, API et emojis absents de la police principale : repli sur DejaVu Sans (caractères actifs)
\newfontface\symfont{DejaVuSans.ttf}[Path=../../../fonts/]
\makeatletter
\newcommand\symactive[1]{\catcode#1=13 \begingroup\lccode`\~=#1 \lowercase{\endgroup\def~}{{\symfont\char#1}}}
\newcommand\symblank[1]{\catcode#1=13 \begingroup\lccode`\~=#1 \lowercase{\endgroup\def~}{}}
\newcommand\symhyph[1]{\catcode#1=13 \begingroup\lccode`\~=#1 \lowercase{\endgroup\def~}{\mbox{-}}}
\newcount\symc
\newcommand\symrange[2]{\symc=#1 \@whilenum\symc<#2 \do{\expandafter\symactive\expandafter{\the\symc}\advance\symc by 1 }}
\newcommand\symblankrange[2]{\symc=#1 \@whilenum\symc<#2 \do{\expandafter\symblank\expandafter{\the\symc}\advance\symc by 1 }}
\makeatother
\symrange{"0250}{"02FF}\symactive{"2713}\symactive{"2714}\symactive{"2717}\symactive{"2718}\symactive{"2610}\symactive{"2611}\symactive{"2612}
\symrange{"2600}{"27BF}
\catcode"2705=13 \begingroup\lccode`\~="2705 \lowercase{\endgroup\def~}{{\symfont\char"2713}}\symblank{"26BD}\symblank{"26A1}
\symrange{"2460}{"257F}\symactive{"223C}\symactive{"2248}\symactive{"2260}
\symactive{"01F8}\symactive{"01F9}\symactive{"1E3E}\symactive{"1E3F}
\symhyph{"2011}\symhyph{"2E1A}\symblankrange{"1F300}{"1FAFF}\symblank{"FE0F}
% Caractères chinois : WenQuanYi Zen Hei (sous-ensemble GB2312 + pinyin), gras/italique simulés
\setCJKmainfont{WenQuanYiZenHei-subset.ttf}[Path=../../../fonts/,AutoFakeBold=3,AutoFakeSlant=0.2]
\xeCJKsetup{CJKspace=false}
% Pinyin : voyelles à caron (ǎ ǐ ǒ ǔ ǖ ǘ ǚ ǜ) absentes de la police principale -> caractères actifs
\newfontface\pyfont{WenQuanYiZenHei-subset.ttf}[Path=../../../fonts/]
\newcommand\pyactive[1]{\catcode#1=13 \begingroup\lccode`\~=#1 \lowercase{\endgroup\def~}{{\pyfont\char#1}}}
\pyactive{"01CD}\pyactive{"01CE}\pyactive{"01CF}\pyactive{"01D0}\pyactive{"01D1}\pyactive{"01D2}\pyactive{"01D3}\pyactive{"01D4}\pyactive{"01D5}\pyactive{"01D6}\pyactive{"01D7}\pyactive{"01D8}\pyactive{"01D9}\pyactive{"01DA}\pyactive{"01DB}\pyactive{"01DC}
% Maths en sans-serif (Fira Math) avec chiffres et lettres droites de Plus
% Jakarta, pour que les formules se fondent dans le texte.
\usepackage[math-style=ISO,bold-style=ISO]{unicode-math}
\setmathfont{FiraMath-Regular.otf}[Path=../../../fonts/]
\setmathfont{PlusJakartaSans-Medium.ttf}[Path=../../../fonts/,range={up/num,up/{Latin,latin},\mathup}]
\usepackage{icomma} % virgule décimale sans espace en mode math
% Caractères mathématiques Unicode absents des polices de texte (Plus Jakarta,
% Archivo Black) : rendus par leur équivalent mathématique, en texte comme en
% formule — sinon ils s'affichent en carré vide (ex. « Arithmétique dans ℤ »).
\usepackage{newunicodechar}
\newunicodechar{ℝ}{\ensuremath{\mathbb{R}}}
\newunicodechar{ℤ}{\ensuremath{\mathbb{Z}}}
\newunicodechar{ℂ}{\ensuremath{\mathbb{C}}}
\newunicodechar{ℕ}{\ensuremath{\mathbb{N}}}
\newunicodechar{ℚ}{\ensuremath{\mathbb{Q}}}
\newunicodechar{𝔻}{\ensuremath{\mathbb{D}}}
\newunicodechar{α}{\ensuremath{\alpha}}
\newunicodechar{β}{\ensuremath{\beta}}
\newunicodechar{γ}{\ensuremath{\gamma}}
\newunicodechar{δ}{\ensuremath{\delta}}
\newunicodechar{ε}{\ensuremath{\varepsilon}}
\newunicodechar{θ}{\ensuremath{\theta}}
\newunicodechar{λ}{\ensuremath{\lambda}}
\newunicodechar{μ}{\ensuremath{\mu}}
\newunicodechar{σ}{\ensuremath{\sigma}}
\newunicodechar{τ}{\ensuremath{\tau}}
\newunicodechar{φ}{\ensuremath{\varphi}}
\newunicodechar{ψ}{\ensuremath{\psi}}
\newunicodechar{ω}{\ensuremath{\omega}}
\newunicodechar{Σ}{\ensuremath{\Sigma}}
% majuscules produites par \MakeUppercase dans les titres (α-aminés -> Α)
\newunicodechar{Α}{\ensuremath{\alpha}}
\newunicodechar{Β}{\ensuremath{\beta}}
\newunicodechar{Γ}{\ensuremath{\Gamma}}
\newunicodechar{Δ}{\ensuremath{\Delta}}
\newunicodechar{Ε}{\ensuremath{\varepsilon}}
\newunicodechar{Θ}{\ensuremath{\Theta}}
\newunicodechar{Λ}{\ensuremath{\Lambda}}
\newunicodechar{Μ}{\ensuremath{\mu}}
\newunicodechar{Τ}{\ensuremath{\tau}}
\newunicodechar{Φ}{\ensuremath{\Phi}}
\newunicodechar{Ψ}{\ensuremath{\Psi}}
\newunicodechar{Ω}{\ensuremath{\Omega}}
\newunicodechar{↦}{\ensuremath{\mapsto}}
\newunicodechar{∈}{\ensuremath{\in}}
\newunicodechar{∉}{\ensuremath{\notin}}
\newunicodechar{∧}{\ensuremath{\wedge}}
\newunicodechar{⇔}{\ensuremath{\Leftrightarrow}}
\newunicodechar{⟺}{\ensuremath{\Longleftrightarrow}}
\newunicodechar{⇒}{\ensuremath{\Rightarrow}}
\newunicodechar{⟹}{\ensuremath{\Longrightarrow}}
\newunicodechar{∘}{\ensuremath{\circ}}
\newunicodechar{∪}{\ensuremath{\cup}}
\newunicodechar{∩}{\ensuremath{\cap}}
\newunicodechar{≡}{\ensuremath{\equiv}}
\newunicodechar{⊂}{\ensuremath{\subset}}
\newunicodechar{⊥}{\ensuremath{\perp}}
\newunicodechar{∃}{\ensuremath{\exists}}
\newunicodechar{∀}{\ensuremath{\forall}}
\newunicodechar{∛}{\ensuremath{\sqrt[3]{\,}}}
\newunicodechar{∜}{\ensuremath{\sqrt[4]{\,}}}
\newunicodechar{⃗}{\ensuremath{{}^{\to}}}
\newunicodechar{ⁿ}{\ifmmode^{n}\else\textsuperscript{\ensuremath{n}}\fi}
\newunicodechar{ˣ}{\ifmmode^{x}\else\textsuperscript{\ensuremath{x}}\fi}
\newunicodechar{ᵇ}{\ifmmode^{b}\else\textsuperscript{\ensuremath{b}}\fi}
\newunicodechar{ʳ}{\ifmmode^{r}\else\textsuperscript{\ensuremath{r}}\fi}
\newunicodechar{ᵘ}{\ifmmode^{u}\else\textsuperscript{\ensuremath{u}}\fi}
\newunicodechar{ʸ}{\ifmmode^{y}\else\textsuperscript{\ensuremath{y}}\fi}
\newunicodechar{ᵃ}{\ifmmode^{a}\else\textsuperscript{\ensuremath{a}}\fi}
\newunicodechar{ᵉ}{\ifmmode^{e}\else\textsuperscript{\ensuremath{e}}\fi}
\newunicodechar{ᵏ}{\ifmmode^{k}\else\textsuperscript{\ensuremath{k}}\fi}
\newunicodechar{ᵖ}{\ifmmode^{p}\else\textsuperscript{\ensuremath{p}}\fi}
\newunicodechar{ᴺ}{\ifmmode^{N}\else\textsuperscript{\ensuremath{N}}\fi}
\newunicodechar{ᶜ}{\ifmmode^{c}\else\textsuperscript{\ensuremath{c}}\fi}
\newunicodechar{ʰ}{\ifmmode^{h}\else\textsuperscript{\ensuremath{h}}\fi}
\newunicodechar{ᵠ}{\ifmmode^{q}\else\textsuperscript{\ensuremath{q}}\fi}
\newunicodechar{ₙ}{\ifmmode_{n}\else\textsubscript{\ensuremath{n}}\fi}
\newunicodechar{ₐ}{\ifmmode_{a}\else\textsubscript{\ensuremath{a}}\fi}
\newunicodechar{ᵢ}{\ifmmode_{i}\else\textsubscript{\ensuremath{i}}\fi}
\newunicodechar{ᵣ}{\ifmmode_{r}\else\textsubscript{\ensuremath{r}}\fi}
\newunicodechar{ₖ}{\ifmmode_{k}\else\textsubscript{\ensuremath{k}}\fi}
\newunicodechar{ᵦ}{\ifmmode_{\beta}\else\textsubscript{\ensuremath{\beta}}\fi}
\newunicodechar{ₓ}{\ifmmode_{x}\else\textsubscript{\ensuremath{x}}\fi}
\newfontfamily\popdisplay{ArchivoBlack-Regular.ttf}[Path=../../../fonts/]
\newfontfamily\poplogo{SpaceGrotesk-Bold.ttf}[Path=../../../fonts/]
\newfontfamily\popsemi{PlusJakartaSans-SemiBold.ttf}[Path=../../../fonts/]
\newfontfamily\popxbold{PlusJakartaSans-ExtraBold.ttf}[Path=../../../fonts/]
\newfontfamily\popxboldls{PlusJakartaSans-ExtraBold.ttf}[Path=../../../fonts/,LetterSpace=3.0]

\setlist[enumerate]{leftmargin=1.7em,itemsep=5pt,topsep=4pt}
\setlist[itemize]{leftmargin=1.7em,itemsep=4pt,topsep=4pt,label={\color{poporange}\textbullet}}
\setlength{\parindent}{0pt}
\setlength{\parskip}{5pt}
\renewcommand{\arraystretch}{1.25}

% pgfplots : style Pop par défaut
\pgfplotsset{
  every axis/.append style={axis line style={popink,line width=1pt},tick style={popink},
    grid style={popline,line width=0.5pt},label style={font=\small},tick label style={font=\footnotesize}},
  popcurve/.style={poporange,line width=1.8pt,no markers,smooth},
}
% Même style disponible dans un \draw TikZ simple (hors pgfplots)
\tikzset{popcurve/.style={poporange,line width=1.8pt}}
\tikzset{popgrid/.style={popline,very thin}}
\colorlet{popcurve}{poporange} % le nom du style est parfois utilisé comme couleur

% ============================================================
% Pill / squiggle
% ============================================================
\newcommand{\poppill}[3]{%
  \tikz[baseline=-0.55ex]\node[draw=popink,line width=1.2pt,rounded corners=2.6pt,%
    fill=#2,inner xsep=6.5pt,inner ysep=3pt]{%
    \popxboldls\fontsize{7}{7}\selectfont\color{#3}\MakeUppercase{#1}};%
}
\newcommand{\popsquiggle}[1]{%
  \tikz[baseline=-0.4ex]{
    \draw[#1,line width=2.6pt,line cap=round]
      plot [smooth] coordinates {(0,0.09) (0.34,0.01) (0.68,0.21) (1.02,0.01) (1.36,0.10)};
  }%
}
% Mot-clé surligné (marqueur orange sous le texte)
\newcommand{\cle}[1]{{\popxbold\color{popdark}#1}}

% ============================================================
% En-tête / pied
% ============================================================
\pagestyle{fancy}
\fancyhf{}
\renewcommand{\headrulewidth}{0pt}
\renewcommand{\footrulewidth}{0pt}
\lhead{\raisebox{-0.15ex}{\poplogo\fontsize{13}{13}\selectfont\color{popink}OnBuch\color{poporange}+}}
\rhead{\poppill{\DOCMATIERE\ \textbullet\ \DOCNIVEAU\ \textbullet\ Leçon \DOCLECON}{white}{popink}}
\lfoot{\popxbold\fontsize{7.6}{7.6}\selectfont\color{popmuted}\MakeUppercase{Cours signé OnBuch+}\ \textbullet\ {\popsemi\DOCTITRE}}
\rfoot{\tikz[baseline=-0.6ex]\node[rounded corners=8.5pt,fill=popink,minimum height=17pt,minimum width=17pt,inner xsep=5pt,inner sep=0]{\popxbold\fontsize{8}{8}\selectfont\color{popcream}\thepage\,/\,\pageref*{LastPage}};}

% ============================================================
% Titres
% ============================================================
\newcommand{\chaptitre}[2]{%
  \begin{tcolorbox}[enhanced,colback=poporange,colframe=popink,boxrule=2.6pt,arc=11pt,
      drop shadow=popink,left=15pt,right=15pt,top=12pt,bottom=11pt,before skip=3pt,after skip=18pt]
    \poppill{#2}{popink}{popcream}\par\vspace{9pt}
    {\raggedright\hyphenpenalty=10000\exhyphenpenalty=10000\popdisplay\fontsize{22}{24}\selectfont\color{popcream}\MakeUppercase{#1}\par}
  \end{tcolorbox}
}
% Section numérotée : pastille orange avec le numéro + titre Archivo
\newcounter{popsec}
\newcommand{\coursec}[1]{%
  \stepcounter{popsec}\setcounter{popsub}{0}%
  \par\vspace{16pt}\noindent
  \begin{minipage}{\linewidth}
  \tikz[baseline=-0.7ex]\node[draw=popink,line width=1.6pt,fill=poporange,rounded corners=4pt,minimum size=22pt,inner sep=0]{\popdisplay\fontsize{12}{12}\selectfont\color{popcream}\thepopsec};%
  \hspace{8pt}{\popdisplay\fontsize{15}{17}\selectfont\color{popink}\MakeUppercase{#1}}\par
  \vspace{4pt}\noindent{\color{poporange}\rule{36pt}{3.6pt}}
  \end{minipage}\par\nopagebreak\vspace{8pt}
}
\newcounter{popsub}
\newcommand{\courssub}[1]{%
  \stepcounter{popsub}%
  \par\vspace{9pt}\noindent{\popxbold\fontsize{11.5}{13}\selectfont\color{popink}{\color{poporange}\thepopsec.\thepopsub}\hspace{6pt}#1}\par\nopagebreak\vspace{4pt}
}

% ============================================================
% Boîtes pédagogiques — même recette : fond blanc, bordure épaisse colorée,
% ombre dure encre, badge plein. Titre optionnel [#1].
% ============================================================
\tcbset{popbase/.style={breakable,enhanced,colback=white,boxrule=2.3pt,arc=9pt,
  shadow={3pt}{-3pt}{0pt}{popink},left=14pt,right=14pt,top=11pt,bottom=11pt,before skip=11pt,after skip=15pt,
  fontupper=\popsemi\fontsize{10.6}{14.5}\selectfont\color{popink},
  fontlower=\popsemi\fontsize{10.6}{14.5}\selectfont\color{popink}}}
% Coche dessinée (le glyphe ✓ n'existe pas dans les polices du document)
\newcommand{\popcheck}[1]{\tikz[baseline=-0.5ex]\draw[#1,line width=1.6pt,line cap=round,line join=round](-0.13,0.0)--(-0.04,-0.09)--(0.13,0.1);}
\providecommand{\coche}{\popcheck{popgreen}}
\providecommand{\resultat}[1]{\par\smallskip\noindent\fcolorbox{popgreen}{popgreenL}{\parbox{\dimexpr\linewidth-2\fboxsep-2\fboxrule\relax}{\textbf{Résultat.} #1}}\par\smallskip}
\newcommand{\popboxhead}[3]{\poppill{#1}{#2}{white}\ifx\relax#3\relax\else\hspace{6pt}{\popxbold\fontsize{10.6}{12}\selectfont\color{popink}#3}\fi\par\vspace{7pt}}

\newtcolorbox{aretenir}[1][]{popbase,colframe=popgreen,before upper={\popboxhead{À retenir}{popgreen}{#1}}}
% Texte en chinois (document de lecture, dialogue, extrait) : cadre
% d'étude. Un titre optionnel (source, auteur) s'affiche dans l'en-tête.
\newtcolorbox{textebox}[1][]{popbase,colframe=popblue,before upper={\popboxhead{文本 · Texte}{popblue}{#1}},after upper={}}
% Marques ✓ / ✗ (forme correcte / fautive) souvent utilisées par les modèles
\providecommand{\cross}{\textcolor{poppink}{\ensuremath{\times}}}
\providecommand{\xmark}{\cross}\providecommand{\crossmark}{\cross}\providecommand{\cmark}{\ensuremath{\checkmark}}
% Ligne de réponse / justification à compléter (inventée par les modèles)
\providecommand{\justification}[1]{\par\noindent\textit{Justification~:} \dotfill\par}
% \textipa (package tipa non chargé) : la police ne couvre pas l'API -> simple passage
\providecommand{\textipa}[1]{#1}
% Soulignement (ulem non chargé) : \uline, \ul
\providecommand{\uline}[1]{\underline{#1}}\providecommand{\ul}[1]{\underline{#1}}
% \glossaire{\item ...} inventée par les modèles : liste de glossaire
\providecommand{\glossaire}[1]{\begin{itemize}#1\end{itemize}}
% \alert (beamer) inventée par les modèles : texte mis en évidence
\providecommand{\alert}[1]{\textbf{\textcolor{poppink}{#1}}}
% variantes ASCII d'ordinaux écrites en macro
\providecommand{\iere}{\textsuperscript{re}}\providecommand{\ieme}{\textsuperscript{e}}\providecommand{\ere}{\textsuperscript{re}}\providecommand{\eme}{\textsuperscript{e}}\providecommand{\er}{\textsuperscript{er}}\providecommand{\ie}{\textsuperscript{e}}
% Ordinaux français écrits en macro par les modèles : 1\ière, 3\ième
\newcommand{\ière}{\textsuperscript{re}}\newcommand{\ième}{\textsuperscript{e}}
% \exemple (inventée par les modèles) : étiquette « Exemple : »
\providecommand{\exemple}{\textbf{Exemple~:}\ }
% \square utilisable aussi hors mode math (cases à cocher)
\AtBeginDocument{\let\squareMath\square\renewcommand{\square}{\ensuremath{\squareMath}}}
% Mot ou phrase en chinois à l'intérieur d'un paragraphe français
\newcommand{\zh}[1]{\ifmmode\text{#1}\else{#1}\fi}
\let\eng\zh \let\esp\zh \let\all\zh % alias tolérés
% flèches d'intonation inventées par les modèles
\providecommand{\textsearrow}{}\renewcommand{\textsearrow}{\ensuremath{\searrow}}\providecommand{\textnearrow}{}\renewcommand{\textnearrow}{\ensuremath{\nearrow}}
% \fr{..} : mot français (inventé par les modèles)
\providecommand{\fr}[1]{\textit{#1}}
% zhbox : encadré de texte chinois inventé par les modèles
\newenvironment{zhbox}{\begin{quote}}{\end{quote}}
% \vs : « versus » inventé par les modèles
\providecommand{\vs}{\textit{vs}\ }
% \blank : case à compléter inventée par les modèles
\providecommand{\blank}[1][]{\underline{\hspace{1.6em}}}
\newtcolorbox{exemplebox}[1][]{popbase,colframe=poporange,before upper={\popboxhead{Exemple}{poporange}{#1}}}
\newtcolorbox{definition}[1][]{popbase,colframe=popblue,before upper={\popboxhead{Définition}{popblue}{#1}}}
\newtcolorbox{propriete}[1][]{popbase,colframe=poppurple,before upper={\popboxhead{Propriété}{poppurple}{#1}}}
\newtcolorbox{methode}[1][]{popbase,colframe=popgold,before upper={\popboxhead{Méthode}{popgold}{#1}}}
\newtcolorbox{attention}[1][]{popbase,colframe=poppink,before upper={\popboxhead{Attention}{poppink}{#1}}}
\newtcolorbox{experience}[1][]{popbase,colframe=popblue,colback=popblueL,before upper={\popboxhead{Expérience}{popblue}{#1}}}
% « Le savais-tu ? » : carte violette pleine (contexte, culture, Cameroun)
\newtcolorbox{savaistu}[1][]{popbase,colback=poppurple,colframe=popink,
  fontupper=\popsemi\fontsize{10.4}{14.2}\selectfont\color{white},
  before upper={\poppill{Le savais-tu\,?}{white}{poppurple}\ifx\relax#1\relax\else\hspace{6pt}{\popxbold\color{white}#1}\fi\par\vspace{7pt}}}
% Exercice résolu : énoncé en haut, \tcblower puis la solution sur fond crème
\newcounter{popexo}\newcounter{popexr}
\newtcolorbox{exoresolu}[1][]{popbase,colframe=poporange,colbacklower=popcream,
  segmentation style={popink,line width=1.2pt,dashed},
  before upper={\stepcounter{popexr}\popboxhead{Exercice résolu \thepopexr}{poporange}{#1}},
  before lower={\poppill{Solution}{popink}{popcream}\par\vspace{6pt}}}
% Exercice (non résolu en place) — la correction est dans la partie Corrigés
\newtcolorbox{exercice}[1][]{popbase,colframe=popink,
  before upper={\stepcounter{popexo}\popboxhead{Exercice \thepopexo}{popink}{#1}}}
% Corrigé
\newtcolorbox{corrige}[1][]{popbase,colframe=popgreen,colback=popgreenL,
  before upper={\popboxhead{Corrigé}{popgreen}{#1}}}
% Objectifs / prérequis / bilan
\newtcolorbox{objectifs}{popbase,colframe=popink,colback=poporangeL,
  before upper={\popboxhead{Objectifs de la leçon}{popink}{}}}
\newtcolorbox{prerequis}{popbase,colframe=popink,colback=popblueL,
  before upper={\popboxhead{Prérequis}{popblue}{}}}
\newtcolorbox{bilan}{popbase,colframe=popink,colback=white,boxrule=2.8pt,
  before upper={\popboxhead{Fiche bilan}{poporange}{L'essentiel à connaître pour l'examen}}}
% Figure encadrée avec légende
\newtcolorbox{popfigure}[1][]{enhanced,breakable=false,colback=white,colframe=popline,boxrule=1.4pt,arc=8pt,
  left=10pt,right=10pt,top=10pt,bottom=8pt,before skip=10pt,after skip=12pt,halign=center,
  lower separated=false,halign lower=center,
  fontlower=\popsemi\fontsize{9}{11.5}\selectfont\color{popmuted},
  #1}
\newcommand{\legende}[1]{\par\vspace{6pt}{\popsemi\fontsize{9}{11.5}\selectfont\color{popmuted}#1}}

% ============================================================
% Signature OnBuch+
% ============================================================
\newcommand{\poplogoinline}[1]{{\poplogo\fontsize{#1}{#1}\selectfont\color{popink}OnBuch\color{poporange}+}}
\newcommand{\signatureonbuch}{%
  \begin{tcolorbox}[enhanced,colback=popink,colframe=popink,boxrule=2.2pt,arc=10pt,
      drop shadow=poporange,left=14pt,right=14pt,top=10pt,bottom=10pt,before skip=0pt,after skip=0pt]
    \tikz[baseline=-0.7ex]\node[circle,fill=popgreen,draw=popcream,line width=1.3pt,minimum size=22pt,inner sep=0]{\popcheck{white}};%
    \hspace{9pt}\begin{minipage}[c]{0.62\linewidth}
      {\popxbold\fontsize{12}{13}\selectfont\color{popcream}Signé\ \textbullet\ {\poplogo OnBuch\color{poporange}+}}\par\vspace{2pt}
      {\raggedright\popsemi\fontsize{8}{10.5}\selectfont\color{popline}\MakeUppercase{Équipe pédagogique OnBuch+}\\\MakeUppercase{Conforme au programme officiel MINESEC}\par}
    \end{minipage}\hfill
    {\poplogo\fontsize{22}{22}\selectfont\color{popcream}OnBuch\color{poporange}+}
  \end{tcolorbox}%
}

% ============================================================
% Page de garde
% #1 titre · #2 matière · #3 niveau · #4 module · #5 n° leçon · #6 durée ·
% #7 description / objectifs (texte ou itemize)
% ============================================================
\newcommand{\pagegarde}[7]{%
  \thispagestyle{empty}
  \newgeometry{a4paper,margin=1.7cm,top=1.6cm,bottom=1.4cm}%
  \begin{tikzpicture}[remember picture,overlay]
    \fill[poporange!18] ([xshift=-3.2cm,yshift=-3.6cm]current page.north east) circle (4.6cm);
    \fill[poppurple!14] ([xshift=2.4cm,yshift=7.6cm]current page.south west) circle (3.8cm);
  \end{tikzpicture}%
  {\poplogo\fontsize{30}{30}\selectfont\color{popink}OnBuch\color{poporange}+}\hfill
  \raisebox{0.6ex}{\poppill{Cours signé \textbullet\ Premium}{popink}{popcream}}\par
  \vspace{1pt}\popsquiggle{poppurple}\par
  \vspace{8pt}
  \begin{tcolorbox}[enhanced,colback=poporange,colframe=popink,boxrule=2.8pt,arc=13pt,
      drop shadow=popink,left=18pt,right=18pt,top=12pt,bottom=13pt,after skip=10pt]
    \poppill{#2 \textbullet\ #3}{popink}{popcream}\hspace{5pt}\poppill{#4}{white}{popink}\par\vspace{8pt}
    {\popdisplay\fontsize{14}{14}\selectfont\color{popink}LEÇON #5}\par\vspace{4pt}
    {\raggedright\hyphenpenalty=10000\exhyphenpenalty=10000\popdisplay\fontsize{25}{27}\selectfont\color{popcream}\MakeUppercase{#1}\par}
    \vspace{8pt}
    {\popxbold\fontsize{9.5}{9.5}\selectfont\color{popink}\MakeUppercase{Durée conseillée : #6}}
  \end{tcolorbox}
  \begin{tcolorbox}[enhanced,colback=white,colframe=popink,boxrule=2.2pt,arc=10pt,
      drop shadow=popink,left=16pt,right=16pt,top=10pt,bottom=10pt,after skip=8pt]
    \poppill{Dans cette leçon}{poppurple}{white}\par\vspace{5pt}
    \popsemi\fontsize{9.3}{12.6}\selectfont\color{popink}#7
  \end{tcolorbox}
  \noindent\begin{minipage}[t]{0.487\linewidth}
    \begin{tcolorbox}[enhanced,colback=popgreen,colframe=popink,boxrule=2.2pt,arc=10pt,
        drop shadow=popink,left=11pt,right=11pt,top=10pt,bottom=10pt,height=3.1cm,valign=center]
      \tcbox[colback=white,colframe=popink,boxrule=1.3pt,arc=5pt,boxsep=0pt,nobeforeafter,
        left=3pt,right=3pt,top=3pt,bottom=3pt,valign=center]{\qrcode[height=1.9cm]{https://play.google.com/store/apps/details?id=com.luvvix.onbuch&hl=fr}}%
      \hspace{8pt}\begin{minipage}[c]{0.5\linewidth}
        {\popdisplay\fontsize{11}{12.5}\selectfont\color{white}\MakeUppercase{Télécharge OnBuch}}\par\vspace{4pt}
        {\raggedright\popsemi\fontsize{8.4}{11}\selectfont\color{white}Gratuit \textbullet\ Google Play\\Tuteur IA, annales, concours, résultats\par}
      \end{minipage}
    \end{tcolorbox}
  \end{minipage}\hfill
  \begin{minipage}[t]{0.487\linewidth}
    \begin{tcolorbox}[enhanced,colback=poppurple,colframe=popink,boxrule=2.2pt,arc=10pt,
        drop shadow=popink,left=11pt,right=11pt,top=10pt,bottom=10pt,height=3.1cm,valign=center]
      \tikz[baseline=-0.7ex]\node[circle,fill=white,draw=popink,line width=1.3pt,minimum size=1.6cm,inner sep=0]{\popdisplay\fontsize{24}{24}\selectfont\color{poppurple}+};%
      \hspace{8pt}\begin{minipage}[c]{0.6\linewidth}
        {\popdisplay\fontsize{11}{12.5}\selectfont\color{white}\MakeUppercase{Passe à OnBuch+}}\par\vspace{4pt}
        {\raggedright\popsemi\fontsize{8.4}{11}\selectfont\color{white}Tous les cours signés, Tuteur IA illimité, fiches PDF — onglet OnBuch+ de l'app\par}
      \end{minipage}
    \end{tcolorbox}
  \end{minipage}
  \vspace{8pt}
  \popcontenu
  \vfill
  \signatureonbuch
  \restoregeometry
  \newpage
}

% Rangée « Ce cours contient » (page de garde)
\newcommand{\poptile}[3]{% #1 couleur #2 gros texte #3 légende
  \begin{tcolorbox}[enhanced,colback=#1,colframe=popink,boxrule=2pt,arc=9pt,shadow={2.5pt}{-2.5pt}{0pt}{popink},
      left=6pt,right=6pt,top=9pt,bottom=9pt,halign=center,valign=center,height=1.75cm,width=\linewidth,nobeforeafter]
    {\popdisplay\fontsize{12.5}{13}\selectfont\color{white}\MakeUppercase{#2}}\par\vspace{3pt}
    {\centering\popsemi\fontsize{7.8}{9.5}\selectfont\color{white}#3\par}
  \end{tcolorbox}}
\newcommand{\popcontenu}{%
  \par\noindent{\popxboldls\fontsize{7.5}{7.5}\selectfont\color{popmuted}CE COURS SIGNÉ CONTIENT}\par\vspace{7pt}
  \noindent\begin{minipage}[t]{0.235\linewidth}\poptile{popblue}{Cours}{complet, illustré, pas à pas}\end{minipage}\hfill
  \begin{minipage}[t]{0.235\linewidth}\poptile{poporange}{Méthodes}{et exercices résolus}\end{minipage}\hfill
  \begin{minipage}[t]{0.235\linewidth}\poptile{poppink}{Exercices}{type examen + corrigés}\end{minipage}\hfill
  \begin{minipage}[t]{0.235\linewidth}\poptile{popgreen}{Fiche bilan}{l'essentiel à retenir}\end{minipage}\par}

% Sommaire
\newtcolorbox{sommaire}{enhanced,colback=white,colframe=popink,boxrule=2.2pt,arc=10pt,
  drop shadow=popink,left=18pt,right=18pt,top=13pt,bottom=13pt,before skip=6pt,after skip=20pt,
  before upper={\poppill{Sommaire}{poppurple}{white}\par\vspace{9pt}\popsemi\fontsize{10.6}{14.5}\selectfont\color{popink}}}

% Signature de fin de cours — poussée tout en bas de la dernière page
\newcommand{\findecours}{%
  \par\vfill\nopagebreak
  \begin{center}\popsquiggle{poporange}\end{center}\vspace{4pt}
  \signatureonbuch
}
\AtBeginDocument{\def\llbracket{\mathopen{[\mkern-3.5mu[}}\def\rrbracket{\mathclose{]\mkern-3.5mu]}}} % intervalles d'entiers (glyphe ⟦ absent de la police)
% Glyphes absents de Fira Math : redéfinis à partir de glyphes disponibles
\AtBeginDocument{%
  \def\setminus{\mathbin{\backslash}}\let\smallsetminus\setminus
  \def\vdots{\mathord{\vbox{\baselineskip3.2pt\lineskiplimit0pt\kern4pt\hbox{.}\hbox{.}\hbox{.}}}}%
  \def\ddots{\mathinner{\mkern1mu\raise6.5pt\hbox{.}\mkern2mu\raise3.6pt\hbox{.}\mkern2mu\raise0.7pt\hbox{.}\mkern1mu}}%
  \def\triangle{\mathord{\scalebox{0.9}{$\symup{\Delta}$}}}%
  \def\nabla{\mathord{\raisebox{\depth}{\scalebox{1}[-1]{$\symup{\Delta}$}}}}%
  \def\checkmark{\popcheck{popdark}}%
  \def\star{\mathbin{\ast}}\def\bigstar{\mathord{\ast}}%
  \def\diamond{\mathbin{\scalebox{0.6}{$\Diamond$}}}%
  \def\bigcup{\mathop{\mathchoice{\raisebox{-0.3em}{\scalebox{1.7}{$\cup$}}}{\scalebox{1.25}{$\cup$}}{\cup}{\cup}}\displaylimits}%
  \def\bigcap{\mathop{\mathchoice{\raisebox{-0.3em}{\scalebox{1.7}{$\cap$}}}{\scalebox{1.25}{$\cap$}}{\cap}{\cap}}\displaylimits}%
  \def\lesseqqgtr{\mathrel{\lessgtr}}\def\bigoplus{\mathop{\oplus}\displaylimits}%
}
\providecommand{\ssi}{si et seulement si} % abréviation fréquente
\AtBeginDocument{\providecommand{\wideparen}{\overparen}} % arc de cercle

%% FIGFIX-BEGIN v5 — correctifs globaux des figures (docs/onbuch-generation/tools/figfix)
\usepackage{adjustbox}\usepackage{etoolbox}\tcbuselibrary{raster}
% toute figure TikZ/pgfplots plus large que la ligne est réduite à la largeur disponible
\BeforeBeginEnvironment{tikzpicture}{\begin{adjustbox}{max width=\linewidth}}
\AfterEndEnvironment{tikzpicture}{\end{adjustbox}}
% légendes pgfplots lisibles par-dessus les courbes
\pgfplotsset{every axis/.append style={legend style={fill=white,fill opacity=0.92,text opacity=1,draw=black!15,inner sep=3pt}}}
% étiquettes d'arêtes (syntaxe edge["..."]) avec fond blanc
\tikzset{every edge quotes/.append style={fill=white,inner sep=1.2pt}}
%% FIGFIX-END
\begin{document}

\pagegarde{Leçon 17 — Expression orale : protection de l’environnement}{Chinois}{1ère A}{Module 2 : Environnement et santé}{ZHA17}{2 h}{Cette leçon d'expression orale apprend à présenter un exposé simple et à participer à un débat guidé en chinois sur le thème de la protection de l'environnement. Elle mobilise les formules fonctionnelles pour présenter, donner son avis et relancer la parole, avec un travail spécifique sur la prononciation des tons.
\vspace{4pt}
\begin{multicols}{2}\small
\begin{itemize}
\item À la fin de cette leçon, je sais présenter oralement un exposé simple en chinois sur la protection de l'environnement
\item je sais utiliser les formules fonctionnelles pour donner mon avis (我觉得..., 我认为...)
\item je sais exprimer l'accord et le désaccord poliment (我同意 / 我不同意)
\item je sais relancer la parole dans un débat (你呢？你怎么看？)
\item je sais prononcer correctement les tons des mots clés du thème environnement
\item je sais utiliser le vocabulaire essentiel : 环境, 保护, 污染, 垃圾, 树, 水, 空气
\end{itemize}\end{multicols}}

\begin{sommaire}
\begin{multicols}{2}
\begin{enumerate}
\item Vocabulaire clé : l'environnement
\item Dialogue modèle 1 : donner son avis
\item Formules fonctionnelles : présenter, opiner, relancer
\item Dialogue modèle 2 : mini-débat
\item Travail phonétique : les tons du thème
\item Jeux de rôle et débat guidé
\item Activité d'intégration
\item Méthodes et exercices résolus
\item Exercices
\item Corrigés des exercices
\item Fiche bilan
\end{enumerate}
\end{multicols}
\end{sommaire}

\begin{objectifs}
\begin{itemize}
\item À la fin de cette leçon, je sais présenter oralement un exposé simple en chinois sur la protection de l'environnement
\item je sais utiliser les formules fonctionnelles pour donner mon avis (我觉得..., 我认为...)
\item je sais exprimer l'accord et le désaccord poliment (我同意 / 我不同意)
\item je sais relancer la parole dans un débat (你呢？你怎么看？)
\item je sais prononcer correctement les tons des mots clés du thème environnement
\item je sais utiliser le vocabulaire essentiel : 环境, 保护, 污染, 垃圾, 树, 水, 空气
\item je sais construire des phrases d'opinion avec 应该 et 要
\item je sais participer à un jeu de rôle et à un débat guidé en binôme ou en groupe
\end{itemize}
\end{objectifs}

\begin{prerequis}
\begin{itemize}
\item Vocabulaire de la nature et de la santé vu dans le Module 2 (水, 空气, 病, 身体)
\item Structure de phrase sujet-verbe-objet et verbes modaux 要, 应该, 可以
\item Pronoms personnels et questions simples avec 吗, 呢, 怎么
\item Bases de la prononciation des quatre tons et du ton neutre
\item Formules de salutation et de présentation vues en Seconde
\end{itemize}
\end{prerequis}

\newpage

\coursec{Vocabulaire clé : l'environnement}

\courssub{Situation d'accroche et problématique}

À Yaoundé comme à Douala, les pluies de septembre ont encore inondé les quartiers à cause des caniveaux bouchés par les plastiques. Pendant ce temps, en Chine, des millions de lycéens participent chaque année à des campagnes de nettoyage et de plantation d'arbres. Comment dire en chinois : « Il faut protéger l'environnement » ? Comment donner ton avis, être d'accord ou pas d'accord avec ton camarade ? Aujourd'hui, tu vas apprendre à débattre en chinois sur un sujet qui te concerne directement : la protection de l'environnement.

\courssub{Observation : un texte court pour entrer dans le lexique}

\begin{textebox}[Texte d'entrée : protéger notre cadre de vie]
我们要保护环境。\\
Wǒmen yào bǎohù huánjìng.\\
Nous devons protéger l'environnement.\\
\\
空气和水很重要。\\
Kōngqì hé shuǐ hěn zhòngyào.\\
L'air et l'eau sont très importants.\\
\\
不要扔垃圾，请节约用水。\\
Búyào rēng lājī, qǐng jiéyuē yòng shuǐ.\\
Ne jetez pas de déchets, économisez l'eau s'il vous plaît.
\end{textebox}

\noindent
\textbf{Observation guidée :}
\begin{itemize}
    \item Repère les mots qui reviennent : \zh{环境}、\zh{保护}、\zh{空气}、\zh{水}、\zh{垃圾}、\zh{节约}。
    \item Note la structure \zh{要 + verbe} pour exprimer l'obligation / la nécessité (« devoir », « il faut »).
    \item Remarque l'ordre : \zh{不要 + verbe} pour l'interdiction (« ne pas »).
\end{itemize}

\courssub{Lexique de base : noms, verbes et adjectifs du thème}

\begin{definition}[Vocabulaire clé — Environnement et protection]
Le tableau ci-dessous regroupe les \textbf{20 items lexicaux} indispensables pour cette leçon. Ils couvrent les \textbf{noms} (éléments naturels, polluants), les \textbf{verbes d'action} (protéger, polluer, jeter, planter, économiser) et les \textbf{adjectifs d'état} (propre, sale). Apprends-les par \textbf{paires sémantiques} (contraires ou collocations fréquentes) pour faciliter la mémorisation et l'oral.
\end{definition}

\begin{center}
\begin{tabularx}{\linewidth}{|X|X|X|X|}
\hline
\rowcolor{popblueL} \textbf{Caractères} & \textbf{Pinyin} & \textbf{Nature} & \textbf{Traduction} \\
\hline
环境 & huánjìng & nom & environnement \\
\hline
保护 & bǎohù & verbe & protéger \\
\hline
保护环境 & bǎohù huánjìng & locution verbale & protéger l'environnement \\
\hline
污染 & wūrǎn & nom / verbe & pollution / polluer \\
\hline
空气污染 & kōngqì wūrǎn & nom composé & pollution de l'air \\
\hline
水污染 & shuǐ wūrǎn & nom composé & pollution de l'eau \\
\hline
垃圾 & lājī & nom & déchets, ordures \\
\hline
扔垃圾 & rēng lājī & locution verbale & jeter des déchets \\
\hline
树 & shù & nom & arbre \\
\hline
种树 & zhòng shù & locution verbale & planter des arbres \\
\hline
水 & shuǐ & nom & eau \\
\hline
空气 & kōngqì & nom & air \\
\hline
干净 & gānjìng & adjectif & propre \\
\hline
脏 & zāng & adjectif & sale \\
\hline
塑料 & sùliào & nom & plastique \\
\hline
塑料袋 & sùliào dài & nom composé & sac plastique \\
\hline
节约 & jiéyuē & verbe & économiser \\
\hline
节约用水 & jiéyuē yòng shuǐ & locution verbale & économiser l'eau \\
\hline
问题 & wèntí & nom & problème \\
\hline
环境问题 & huánjìng wèntí & nom composé & problème d'environnement \\
\hline
\end{tabularx}
\end{center}

\courssub{Analyse morphologique et collocations fréquentes}

\begin{propriete}[Construction des mots composés : nom + verbe / verbe + nom]
En chinois, les mots composés se forment souvent par \textbf{association de deux morphèmes} :
\begin{itemize}
    \item \textbf{Nom + Verbe} : \zh{空气污染} (air + polluer = pollution de l'air), \zh{水污染} (eau + polluer).
    \item \textbf{Verbe + Nom} : \zh{保护环境} (protéger + environnement), \zh{种树} (planter + arbre), \zh{扔垃圾} (jeter + déchet), \zh{节约用水} (économiser + utiliser + eau).
    \item \textbf{Nom + Nom} : \zh{塑料袋} (plastique + sac), \zh{环境问题} (environnement + problème).
\end{itemize}
\emph{Repère le noyau (tête) du composé : c'est souvent le second élément pour les noms (污染, 袋, 问题) et le premier pour les actions verbales (保护, 种, 扔, 节约).}
\end{propriete}

\begin{exemplebox}[Collocations types à mémoriser par cœur]
\begin{itemize}
    \item \zh{严重污染} (yánzhòng wūrǎn) — pollution grave
    \item \zh{减少垃圾} (jiǎnshǎo lājī) — réduire les déchets
    \item \zh{保护树木} (bǎohù shùmù) — protéger les arbres
    \item \zh{空气干净} (kōngqì gānjìng) — l'air est propre
    \item \zh{水很脏} (shuǐ hěn zāng) — l'eau est très sale
    \item \zh{少用塑料袋} (shǎo yòng sùliào dài) — utiliser moins de sacs plastiques
    \item \zh{解决环境问题} (jiějué huánjìng wèntí) — résoudre les problèmes d'environnement
\end{itemize}
\end{exemplebox}

\courssub{Point de grammaire lexicale : \zh{污染} — nom et verbe}

\begin{definition}[污染 wūrǎn : double nature grammaticale]
\zh{污染} signifie à la fois \textbf{« pollution » (nom)} et \textbf{« polluer » (verbe)}, selon sa place dans la phrase.
\begin{itemize}
    \item \textbf{Nom} : sujet ou objet. Ex. \zh{空气污染很严重。} (La pollution de l'air est grave.)
    \item \textbf{Verbe} : prédicat, souvent suivi d'un objet. Ex. \zh{工厂污染了河水。} (L'usine a pollué l'eau de la rivière.)
\end{itemize}
\emph{En français, on distingue « pollution » (nom) et « polluer » (verbe). En chinois, le même mot couvre les deux emplois. Le contexte (particules, place, compléments) tranche.}
\end{definition}

\begin{exemplebox}[污染 : nom vs verbe]
\begin{tabularx}{\linewidth}{|X|X|X|}
\hline
\rowcolor{popgreenL} \textbf{Phrase chinoise} & \textbf{Pinyin} & \textbf{Traduction} \\
\hline
\zh{空气污染很严重。} & Kōngqì wūrǎn hěn yánzhòng. & La pollution de l'air est très grave. (nom) \\
\hline
\zh{谁污染了环境？} & Shéi wūrǎn le huánjìng ? & Qui a pollué l'environnement ? (verbe + \zh{了}) \\
\hline
\zh{不要污染水源。} & Búyào wūrǎn shuǐyuán. & Ne polluez pas la source d'eau. (verbe) \\
\hline
\zh{解决污染问题。} & Jiějué wūrǎn wèntí. & Résoudre le problème de pollution. (nom en attribut) \\
\hline
\end{tabularx}
\end{exemplebox}

\coursec{Dialogue modèle 1 : Donner son avis}

\courssub{Mise en situation : deux élèves dans la cour du lycée}

\begin{textebox}[Dialogue 1 : L'état de l'école]
A: 你看，操场上有很多垃圾。\\
Nǐ kàn, cāochǎng shàng yǒu hěn duō lājī.\\
Regarde, il y a beaucoup de déchets sur le terrain de sport.\\
\\
B: 是啊，环境很脏。我认为我们要保护环境。\\
Shì a, huánjìng hěn zāng. Wǒ rènwéi wǒmen yào bǎohù huánjìng.\\
Oui, l'environnement est sale. Je pense qu'il faut protéger l'environnement.\\
\\
A: 我同意。少用塑料袋，多种树。\\
Wǒ tóngyì. Shǎo yòng sùliào dài, duō zhòng shù.\\
Je suis d'accord. Utilisons moins de sacs plastiques, plantons plus d'arbres.\\
\\
B: 好主意！我们一起行动吧。\\
Hǎo zhǔyi! Wǒmen yìqǐ xíngdòng ba.\\
Bonne idée ! Agissons ensemble.
\end{textebox}

\noindent
\textbf{Observation guidée :}
\begin{itemize}
    \item \zh{我认为} (wǒ rènwéi) — « Je pense que / À mon avis », introduit un avis personnel.
    \item \zh{我同意} (wǒ tóngyì) — « Je suis d'accord ». \zh{我不同意} (wǒ bù tóngyì) — « Je ne suis pas d'accord ».
    \item \zh{好主意} (hǎo zhǔyi) — « Bonne idée ! » (réaction positive).
    \item \zh{一起行动吧} (yìqǐ xíngdòng ba) — « Agissons ensemble » (\zh{吧} suggère une proposition douce).
\end{itemize}

\coursec{Formules fonctionnelles : présenter, opiner, relancer}

\begin{propriete}[Boîte à outils pour l'oral : structures clés]
Ces formules sont les \textbf{briques de base} pour tout exposé ou débat en classe. Mémorise-les par blocs fonctionnels.
\end{propriete}

\begin{exemplebox}[1. Présenter un sujet / Introduire son exposé]
\begin{tabularx}{\linewidth}{|X|X|X|}
\hline
\rowcolor{popblueL} \textbf{Chinois} & \textbf{Pinyin} & \textbf{Français} \\
\hline
\zh{今天我要谈谈……} & Jīntiān wǒ yào tántan…… & Aujourd'hui, je vais parler de… \\
\hline
\zh{我的题目是……} & Wǒ de tímù shì…… & Mon sujet est… \\
\hline
\zh{关于环境保护，我想说……} & Guānyú huánjìng bǎohù, wǒ xiǎng shuō…… & Concernant la protection de l'environnement, je veux dire… \\
\hline
\zh{首先，……} & Shǒuxiān,…… & Premièrement,… \\
\hline
\zh{其次，……} & Qícì,…… & Deuxièmement,… \\
\hline
\zh{最后，……} & Zuìhòu,…… & Enfin,… \\
\hline
\end{tabularx}
\end{exemplebox}

\begin{exemplebox}[2. Donner son avis / Exprimer accord ou désaccord]
\begin{tabularx}{\linewidth}{|X|X|X|}
\hline
\rowcolor{poporangeL} \textbf{Chinois} & \textbf{Pinyin} & \textbf{Français} \\
\hline
\zh{我认为……} & Wǒ rènwéi…… & Je pense que… / À mon avis… \\
\hline
\zh{我觉得……} & Wǒ juéde…… & Je trouve que… / J'ai l'impression que… \\
\hline
\zh{按我说，……} & Àn wǒ shuō,…… & Selon moi,… / Si tu veux mon avis,… \\
\hline
\zh{我同意。/ 我赞成。} & Wǒ tóngyì. / Wǒ zànchéng. & Je suis d'accord. / J'approuve. \\
\hline
\zh{我不同意。/ 我反对。} & Wǒ bù tóngyì. / Wǒ fǎnduì. & Je ne suis pas d'accord. / Je m'oppose. \\
\hline
\zh{这有道理。} & Zhè yǒu dàoli. & C'est logique / Il y a du vrai. \\
\hline
\zh{恐怕不太对吧？} & Kǒngpà bù tài duì ba? & Je crains que ce ne soit pas tout à fait exact ? \\
\hline
\end{tabularx}
\end{exemplebox}

\begin{exemplebox}[3. Relancer / Interagir / Conclure]
\begin{tabularx}{\linewidth}{|X|X|X|}
\hline
\rowcolor{popgreenL} \textbf{Chinois} & \textbf{Pinyin} & \textbf{Français} \\
\hline
\zh{你呢？} & Nǐ ne? & Et toi ? \\
\hline
\zh{你怎么看？} & Nǐ zěnme kàn? & Qu'en penses-tu ? / Comment tu vois les choses ? \\
\hline
\zh{还有什么建议？} & Hái yǒu shénme jiànyi? & D'autres suggestions ? \\
\hline
\zh{我们一起……吧。} & Wǒmen yìqǐ…… ba. & Et si nous… ensemble ? \\
\hline
\zh{总之，……} & Zǒngzhī,…… & En résumé,… / Bref,… \\
\hline
\zh{谢谢大家。} & Xièxie dàjiā. & Merci à tous. \\
\hline
\end{tabularx}
\end{exemplebox}

\coursec{Dialogue modèle 2 : Mini-débat}

\courssub{Sujet : « Faut-il interdire les sacs plastiques à l'école ? »}

\begin{textebox}[Dialogue 2 : Débat guidé (4 élèves)]
Modérateur: 今天我们讨论：学校要不要禁止塑料袋？\\
Jīntiān wǒmen tǎolùn: xuéxiào yào bù yào jìnzhǐ sùliào dài?\\
Aujourd'hui nous discutons : l'école doit-elle interdire les sacs plastiques ?\\
\\
Élève A (Pour): 我赞成禁止。塑料袋污染环境，很难分解。\\
Wǒ zànchéng jìnzhǐ. Sùliào dài wūrǎn huánjìng, hěn nán fēnjiě.\\
Je suis pour l'interdiction. Les sacs plastiques polluent l'environnement, ils se décomposent difficilement.\\
\\
Élève B (Contre): 我不同意。塑料袋很方便，学生买饭要用。\\
Wǒ bù tóngyì. Sùliào dài hěn fāngbiàn, xuéshēng mǎi fàn yào yòng.\\
Je ne suis pas d'accord. Les sacs plastiques sont pratiques, les élèves en ont besoin pour acheter à manger.\\
\\
Élève C (Nuance): 我觉得可以少用，自带饭盒。\\
Wǒ juéde kěyǐ shǎo yòng, zìdài fànhé.\\
Je trouve qu'on peut en utiliser moins, en apportant sa propre boîte à repas.\\
\\
Modérateur: 好，我们一起少用塑料袋，保护环境。\\
Hǎo, wǒmen yìqǐ shǎo yòng sùliào dài, bǎohù huánjìng.\\
Bien, ensemble, utilisons moins de sacs plastiques, protégeons l'environnement.
\end{textebox}

\noindent
\textbf{Analyse du débat :}
\begin{itemize}
    \item \zh{讨论} (tǎolùn) — discuter, débattre.
    \item \zh{禁止} (jìnzhǐ) — interdire (verbe).
    \item \zh{分解} (fēnjiě) — se décomposer (biodégradabilité).
    \item \zh{方便} (fāngbiàn) — pratique, commode.
    \item \zh{自带} (zìdài) — apporter soi-même.
    \item Structure \zh{要不要} (yào bù yào) — « faut-il / est-ce qu'il faut » (question alternative).
\end{itemize}

\coursec{Travail phonétique : les tons du thème}

\begin{attention}[垃圾 lājī : attention au deuxième ton !]
\zh{垃圾} se prononce \textbf{lājī} (1er ton + 1er ton). Les francophones ont tendance à dire « lāji » avec un \textbf{ton neutre fautif} sur la deuxième syllabe (\zh{ji} au lieu de \zh{jī}). En chinois standard, \zh{垃圾} garde ses deux tons pleins. Entraîne-toi : \zh{扔垃圾} (rēng lājī), \zh{很多垃圾} (hěn duō lājī).
\end{attention}

\begin{attention}[Le sandhi de \zh{不} (bù) : règle d'or à l'oral]
\zh{不} (bù, 4e ton) change de ton devant un autre 4e ton : il devient \textbf{2e ton (bú)}.
\begin{itemize}
    \item \zh{不要} → \textbf{bú yào} (interdiction)
    \item \zh{不干净} → \textbf{bù gānjìng} (1er ton après \zh{不} → \zh{不} reste 4e ton)
    \item \zh{不脏} → \textbf{bù zāng} (1er ton après \zh{不} → \zh{不} reste 4e ton)
    \item \zh{不同意} → \textbf{bú tóngyì} (2e ton \zh{同} → \zh{不} devient 2e ton)
\end{itemize}
\emph{Erreur fréquente : dire \zh{bù yào} ou \zh{bù tóngyì}. Corrige-toi systématiquement à l'oral.}
\end{attention}

\begin{propriete}[Tons à surveiller pour cette leçon — Tableau récapitulatif]
\begin{center}
\begin{tabularx}{\linewidth}{|X|X|X|}
\hline
\rowcolor{poporangeL} \textbf{Mot} & \textbf{Pinyin} & \textbf{Difficulté / Note} \\
\hline
环境 & huánjìng & 2e + 4e (montant puis descendant) \\
\hline
保护 & bǎohù & 3e + 4e (3e ton bas + 4e ton chute) \\
\hline
污染 & wūrǎn & 1e + 3e (haut plat + bas) \\
\hline
垃圾 & lājī & 1e + 1e (deux tons hauts plats, pas de ton neutre) \\
\hline
种树 & zhòng shù & 4e + 4e (deux chutes ; \zh{种} = planter, 4e ton) \\
\hline
干净 & gānjìng & 1e + 4e (haut plat + chute) \\
\hline
脏 & zāng & 1e ton (haut plat) \\
\hline
塑料 & sùliào & 4e + 4e (deux chutes) \\
\hline
节约 & jiéyuē & 2e + 1e (montant + haut plat) \\
\hline
问题 & wèntí & 4e + 2e (chute + montant) \\
\hline
禁止 & jìnzhǐ & 4e + 3e (chute + bas) \\
\hline
分解 & fēnjiě & 1e + 3e (haut plat + bas) \\
\hline
方便 & fāngbiàn & 1e + 4e (haut plat + chute) \\
\hline
\end{tabularx}
\end{center}
\end{propriete}

\courssub{Exercice phonétique guidé : lecture en chaîne}
Lis à haute voix la liste ci-dessous en respectant les tons et le sandhi de \zh{不}. Le professeur corrige chaque élève.
\begin{enumerate}
    \item \zh{保护环境} (bǎohù huánjìng)
    \item \zh{空气污染} (kōngqì wūrǎn)
    \item \zh{不要扔垃圾} (bú yào rēng lājī)
    \item \zh{种树} (zhòng shù)
    \item \zh{节约用水} (jiéyuē yòng shuǐ)
    \item \zh{少用塑料袋} (shǎo yòng sùliào dài)
    \item \zh{我不同意} (wǒ bú tóngyì)
    \item \zh{环境问题很严重} (huánjìng wèntí hěn yánzhòng)
\end{enumerate}

\coursec{Jeux de rôle et débat guidé}

\begin{experience}[Activité 1 : Jeu de rôle « Au marché / À la cantine » (par paires)]
\textbf{Consigne :} L'élève A est un client, l'élève B est un vendeur / caissier. Le client refuse un sac plastique et propose une alternative.
\begin{itemize}
    \item \textbf{A:} \zh{你好，买两个包子。} (Bonjour, deux brioches vapeur.)
    \item \textbf{B:} \zh{要塑料袋吗？} (Vous voulez un sac plastique ?)
    \item \textbf{A:} \zh{不要塑料袋，我自带盒子。/ 请少用塑料袋，保护环境。} (Pas de sac, j'ai ma boîte. / Moins de sacs svp, protégeons l'env.)
    \item \textbf{B:} \zh{好，好主意！} (Bien, bonne idée !)
\end{itemize}
\textbf{Variante :} Inverser les rôles. Ajouter : \zh{节约用水}, \zh{不扔垃圾}.
\end{experience}

\begin{experience}[Activité 2 : Débat structuré « Notre école verte » (groupes de 4-5)]
\textbf{Sujet :} \zh{如何让我们的学校更绿色？} (Comment rendre notre école plus verte ?)
\textbf{Rôles :} 1 Modérateur (distribue la parole, résume), 2-3 Débateurs, 1 Secrétaire (note les arguments clés en chinois).
\textbf{Étapes (15 min) :}
\begin{enumerate}
    \item \textbf{Préparation (3 min) :} Chaque élève prépare 2 arguments avec le vocabulaire de la leçon (\zh{种树}, \zh{节约用水}, \zh{少用塑料袋}, \zh{不扔垃圾}, \zh{保护环境}).
    \item \textbf{Débat (8 min) :} Le modérateur lance : \zh{今天我们讨论：如何让学校更绿色？请说说你的看法。} Les débateurs utilisent les formules : \zh{我认为……}, \zh{我同意/不同意}, \zh{你呢？}, \zh{还有建议？}.
    \item \textbf{Synthèse orale (2 min) :} Le secrétaire présente les 3 meilleures idées du groupe en chinois.
    \item \textbf{Restitution (2 min) :} Un groupe volontaire rejoue son débat devant la classe.
\end{enumerate}
\textbf{Critères d'évaluation (grille partagée) :} Prononciation des tons (5 pts), Utilisation du vocabulaire cible (5 pts), Structures \zh{要/不要}, \zh{我认为}, \zh{同意} (5 pts), Fluidité / Interaction (5 pts).
\end{experience}

\coursec{Élément socioculturel : la Journée de la plantation d'arbres en Chine}

\begin{savaistu}[植树节 — Journée nationale de la plantation d'arbres]
En Chine, le \textbf{12 mars} est la \zh{植树节} (Zhíshùjié), Journée nationale de la plantation d'arbres, instaurée en 1979. Ce jour-là, des millions de Chinois — écoliers, lycéens, employés, retraités — sortent planter des arbres dans les parcs, sur les collines, le long des routes. C'est une obligation civique : chaque citoyen âgé de 11 à 60 ans doit planter \textbf{3 à 5 arbres par an} (ou payer une taxe équivalente). Depuis 1982, plus de \textbf{78 milliards d'arbres} ont été plantés par la population. Au Cameroun, la \zh{植树节} n'existe pas officiellement, mais des opérations « \emph{Green Cameroon} » ou « \emph{Un élève, un arbre} » existent dans certains lycées. \emph{Comparaison :} la Chine a institutionnalisé le geste écologique en devoir citoyen annuel ; le Cameroun mise davantage sur des campagnes ponctuelles portées par des ONG ou le Ministère de l'Environnement (MINEPDED).
\end{savaistu}

\coursec{Carte mentale lexicale : organisation visuelle du vocabulaire}

\begin{popfigure}
\begin{center}\tcbox[colback=popink,colframe=popink,coltext=white,arc=9pt,boxsep=4pt,left=6pt,right=6pt]{\bfseries\small 环境 huánjìng}\end{center}
\vspace{-2pt}
\begin{tcbraster}[raster columns=2,raster equal height=rows,raster column skip=6pt,raster row skip=6pt,enhanced,blankest]
\begin{tcolorbox}[enhanced,colback=white,colframe=poporange,boxrule=0.9pt,arc=3pt,title={污染 wūrǎn},fonttitle=\bfseries\scriptsize,coltitle=white,colbacktitle=poporange,left=4pt,right=4pt,top=2pt,bottom=2pt,boxsep=1pt,toptitle=1.5pt,bottomtitle=1.5pt,halign=left]
\begin{itemize}[nosep,leftmargin=*,itemsep=1pt,font=\scriptsize]
\item 空气污染 kōngqì wūrǎn
\item 水污染 shuǐ wūrǎn
\item 严重污染 yánzhòng wūrǎn
\end{itemize}
\end{tcolorbox}
\begin{tcolorbox}[enhanced,colback=white,colframe=popgreen,boxrule=0.9pt,arc=3pt,title={垃圾 lājī},fonttitle=\bfseries\scriptsize,coltitle=white,colbacktitle=popgreen,left=4pt,right=4pt,top=2pt,bottom=2pt,boxsep=1pt,toptitle=1.5pt,bottomtitle=1.5pt,halign=left]
\begin{itemize}[nosep,leftmargin=*,itemsep=1pt,font=\scriptsize]
\item 扔垃圾 rēng lājī
\item 减少垃圾 jiǎnshǎo lājī
\item 塑料袋 sùliào dài
\end{itemize}
\end{tcolorbox}
\begin{tcolorbox}[enhanced,colback=white,colframe=poppurple,boxrule=0.9pt,arc=3pt,title={自然 zìrán},fonttitle=\bfseries\scriptsize,coltitle=white,colbacktitle=poppurple,left=4pt,right=4pt,top=2pt,bottom=2pt,boxsep=1pt,toptitle=1.5pt,bottomtitle=1.5pt,halign=left]
\begin{itemize}[nosep,leftmargin=*,itemsep=1pt,font=\scriptsize]
\item 树 / 种树 shù / zhòng shù
\item 水 shuǐ
\item 空气 kōngqì
\item 干净 / 脏 gānjìng / zāng
\end{itemize}
\end{tcolorbox}
\begin{tcolorbox}[enhanced,colback=white,colframe=poppink,boxrule=0.9pt,arc=3pt,title={行动 xíngdòng},fonttitle=\bfseries\scriptsize,coltitle=white,colbacktitle=poppink,left=4pt,right=4pt,top=2pt,bottom=2pt,boxsep=1pt,toptitle=1.5pt,bottomtitle=1.5pt,halign=left]
\begin{itemize}[nosep,leftmargin=*,itemsep=1pt,font=\scriptsize]
\item 保护环境 bǎohù huánjìng
\item 节约用水 jiéyuē yòng shuǐ
\item 解决问题 jiějué wèntí
\end{itemize}
\end{tcolorbox}
\end{tcbraster}

\legende{Carte mentale du vocabulaire : 4 branches thématiques autour de 环境}
\end{popfigure}

\coursec{Exercice résolu : mise en pratique immédiate}

\begin{exoresolu}[Compléter et traduire]
\textbf{Énoncé :} Complète les phrases avec le mot approprié du tableau de vocabulaire, puis donne le pinyin et la traduction française.
\begin{enumerate}
    \item \zh{工厂 \_\_\_\_ 了河水。} (pollué)
    \item \zh{请大家 \_\_\_\_ 用水。} (économiser)
    \item \zh{这个城市的 \_\_\_\_ 很严重。} (pollution)
    \item \zh{我们要 \_\_\_\_ 树。} (planter)
    \item \zh{不要 \_\_\_\_ 垃圾。} (jeter)
\end{enumerate}

\tcblower
\textbf{Solution détaillée :}
\begin{enumerate}
    \item \zh{工厂 \textbf{污染} 了河水。} — Gōngchǎng \textbf{wūrǎn} le héshuǐ. — L'usine a pollué l'eau de la rivière. (\zh{污染} verbe + \zh{了} aspect accompli)
    \item \zh{请大家 \textbf{节约} 用水。} — Qǐng dàjiā \textbf{jiéyuē} yòng shuǐ. — Je prie tout le monde d'économiser l'eau. (\zh{节约} verbe + objet \zh{用水})
    \item \zh{这个城市的 \textbf{污染} 很严重。} — Zhège chéngshì de \textbf{wūrǎn} hěn yánzhòng. — La pollution de cette ville est grave. (\zh{污染} nom, sujet de \zh{很严重})
    \item \zh{我们要 \textbf{种} 树。} — Wǒmen yào \textbf{zhòng} shù. — Nous devons planter des arbres. (\zh{种树} locution verbale, \zh{种} seul suffit devant \zh{树})
    \item \zh{不要 \textbf{扔} 垃圾。} — Búyào \textbf{rēng} lājī. — Ne jetez pas de déchets. (\zh{扔垃圾} locution verbale, \zh{扔} seul suffit devant \zh{垃圾})
\end{enumerate}
\end{exoresolu}

\coursec{Exercices d'entraînement (à faire en classe ou à la maison)}

\begin{exercice}[Entraînement 1 : Associer mot chinois et définition française]
Relie chaque terme chinois à sa définition.
\begin{enumerate}
    \item \zh{环境} \quad a. Problème, question
    \item \zh{污染} \quad b. Propre
    \item \zh{垃圾} \quad c. Environnement
    \item \zh{干净} \quad d. Pollution / polluer
    \item \zh{问题} \quad e. Déchets, ordures
\end{enumerate}
\end{exercice}

\begin{exercice}[Entraînement 2 : Traduire en chinois (caractères + pinyin)]
\begin{enumerate}
    \item L'air est pollué.
    \item Il faut économiser l'eau.
    \item Ne jetez pas les sacs plastiques.
    \item Planter des arbres protège l'environnement.
    \item Les problèmes d'environnement sont graves.
\end{enumerate}
\end{exercice}

\begin{exercice}[Entraînement 3 : Production orale guidée — « Mon quartier »]
Prépare une courte présentation orale (4-5 phrases) sur l'environnement de ton quartier / ton école à Yaoundé, Douala, Buea, etc. Utilise au moins \textbf{6 mots} du vocabulaire de cette leçon. Structure suggérée :
\begin{enumerate}
    \item Présenter le lieu : \zh{我住在……} (J'habite à…)
    \item Décrire l'air / l'eau : \zh{空气…… / 水……}
    \item Citer un problème : \zh{有……问题} / \zh{很脏}
    \item Proposer une action : \zh{我们要…… / 请……}
    \item Conclure : \zh{保护环境，人人有责。}
\end{enumerate}
\emph{Exemple :} \zh{我住在雅温得。这里空气不干净，有很多垃圾。水很脏。我们要种树，少用塑料袋。保护环境，人人有责。}
\end{exercice}


\begin{corrige}[Entraînement 1]
1–c, 2–d, 3–e, 4–b, 5–a.
\end{corrige}

\begin{corrige}[Entraînement 2]
\begin{enumerate}
    \item \zh{空气污染了。} (Kōngqì wūrǎn le.) ou \zh{空气被污染了。}
    \item \zh{要节约用水。} (Yào jiéyuē yòng shuǐ.)
    \item \zh{不要扔塑料袋。} (Búyào rēng sùliào dài.)
    \item \zh{种树保护环境。} (Zhòng shù bǎohù huánjìng.)
    \item \zh{环境问题很严重。} (Huánjìng wèntí hěn yánzhòng.)
\end{enumerate}
\end{corrige}

\begin{corrige}[Entraînement 3]
Production libre. Vérification par le professeur : tons corrects (surtout \zh{lājī}, \zh{zhòng shù}, \zh{jiéyuē}, \zh{bú yào}), ordre des mots (sujet + temps/lieu + verbe + objet), utilisation de \zh{要/不要} pour obligation/interdiction.
\end{corrige}

\coursec{Dialogue modèle 1 : donner son avis}

Dans la section précédente, nous avons découvert le vocabulaire clé de l'environnement : \zh{环境} (huánjìng, environnement), \zh{垃圾} (lājī, déchets), \zh{保护} (bǎohù, protéger), \zh{种树} (zhòng shù, planter des arbres). Posséder du vocabulaire ne suffit pas : il faut maintenant apprendre à l'\emph{utiliser} dans une conversation réelle. C'est exactement l'objectif de cette section : à travers un dialogue entre deux élèves camerounais, Lili et Dawei, tu vas apprendre à \cle{donner ton avis} en chinois, à \cle{approuver} ton interlocuteur et à \cle{proposer des solutions}. Lis d'abord le dialogue à voix haute, sans chercher à tout analyser : laisse-toi guider par la situation.

\courssub{Le dialogue : Lili et Dawei à Yaoundé}

\begin{textebox}[Dialogue : 在校园里 (zài xiàoyuán lǐ) — Dans la cour de l'école]
A (Lili) : 你看，这里的垃圾太多了！\\
Nǐ kàn, zhèlǐ de lājī tài duō le !\\
« Regarde, il y a trop de déchets ici ! »\\[0.4em]
B (Dawei) : 是啊，我觉得我们应该保护环境。\\
Shì a, wǒ juéde wǒmen yīnggāi bǎohù huánjìng.\\
« Oui, je pense que nous devrions protéger l'environnement. »\\[0.4em]
A (Lili) : 你说得对。我们可以做什么？\\
Nǐ shuō de duì. Wǒmen kěyǐ zuò shénme ?\\
« Tu as raison. Que pouvons-nous faire ? »\\[0.4em]
B (Dawei) : 我们不要扔垃圾，还要多种树。\\
Wǒmen bú yào rēng lājī, hái yào duō zhòng shù.\\
« Ne jetons pas de déchets, et plantons aussi plus d'arbres. »
\end{textebox}

\begin{methode}[Travail oral en binôme]
\begin{enumerate}
\item Lis le dialogue à voix haute avec ton voisin, chacun un rôle, en insistant sur les tons : \zh{垃圾} se prononce \emph{lājī} (premier ton haut et plat, puis premier ton), et \zh{保护} se prononce \emph{bǎohù} (troisième ton qui descend puis remonte, puis quatrième ton qui tombe).
\item Échangez les rôles et relisez en regardant votre partenaire, comme dans une vraie conversation.
\item Cache la traduction française et vérifie que tu comprends chaque réplique rien qu'en lisant les caractères.
\item Termine en récitant le dialogue sans regarder le pinyin.
\end{enumerate}
\end{methode}

\courssub{Observation : repérer les formules d'opinion}

Avant d'énoncer la règle, observons. Dans le dialogue, relève les expressions qui servent à donner un avis ou à réagir à l'avis de l'autre :

\begin{center}
\begin{tabularx}{\linewidth}{|X|X|X|}\hline
\rowcolor{popblueL} Formule & Pinyin & Fonction \\ \hline
\zh{你看} & nǐ kàn & attirer l'attention : « regarde » \\ \hline
\zh{我觉得} & wǒ juéde & donner son avis : « je pense, je trouve » \\ \hline
\zh{应该} & yīnggāi & exprimer le devoir : « devoir, il faudrait » \\ \hline
\zh{你说得对} & nǐ shuō de duì & approuver : « tu as raison » \\ \hline
\zh{可以} & kěyǐ & exprimer la possibilité : « pouvoir » \\ \hline
\zh{不要} & bú yào & interdire / conseiller de ne pas : « ne pas » \\ \hline
\zh{还要} & hái yào & ajouter une action : « et aussi, il faut encore » \\ \hline
\end{tabularx}
\end{center}

Tu remarques que le dialogue suit une progression très naturelle, que tu pourras reproduire dans n'importe quel débat : \textbf{constat} (\zh{垃圾太多了}) $\rightarrow$ \textbf{avis} (\zh{我觉得...}) $\rightarrow$ \textbf{accord} (\zh{你说得对}) $\rightarrow$ \textbf{question} (\zh{我们可以做什么？}) $\rightarrow$ \textbf{proposition} (\zh{我们不要...还要...}). Retiens cette « colonne vertébrale » du débat : c'est elle qui te permettra de parler sans notes.

\courssub{La structure d'opinion : 我觉得 / 我认为 + phrase complète}

\begin{propriete}[Donner son avis avec 觉得 et 认为]
Pour exprimer une opinion, le chinois place le verbe d'opinion \zh{觉得} (juéde, « penser, trouver », plus oral) ou \zh{认为} (rènwéi, « considérer, estimer », plus soutenu) \textbf{directement devant la phrase complète}, sans aucun mot de liaison :\par
\textbf{Sujet + 觉得 / 认为 + phrase complète (sujet + verbe + complément)}\par
Contrairement au français, il n'y a \textbf{pas de « que »} entre le verbe d'opinion et la phrase qui suit :\par
\zh{我觉得 + 我们应该保护环境。} = « Je pense \emph{que} nous devrions protéger l'environnement. »\par
Le mot \zh{觉得} convient aux conversations quotidiennes ; \zh{认为} convient mieux aux exposés et aux débats un peu formels. Les deux se construisent exactement de la même façon.
\end{propriete}

\begin{popfigure}
\begin{tikzpicture}[node distance=0.5cm and 0.7cm,
box/.style={rounded corners=4pt, draw=popblue, fill=popblueL, align=center, inner sep=6pt, font=\small},
box2/.style={rounded corners=4pt, draw=poporange, fill=poporangeL, align=center, inner sep=6pt, font=\small},
box3/.style={rounded corners=4pt, draw=popgreen, fill=popgreenL, align=center, inner sep=6pt, font=\small}]
\node[box, align=center] (s){\zh{我}\\wǒ\\(sujet)};
\node[box2, right=of s, align=center] (v){\zh{觉得 / 认为}\\juéde / rènwéi\\(verbe d'opinion)};
\node[box3, right=of v, align=center] (p){\zh{我们应该保护环境}\\wǒmen yīnggāi bǎohù huánjìng\\(phrase complète)};
\draw[-{Stealth}, thick, popdark] (s) -- (v);
\draw[-{Stealth}, thick, popdark] (v) -- (p);
\node[below=0.6cm of v, align=center, font=\small\itshape, text=popmuted] {« Je pense que nous devrions protéger l'environnement. »};
\end{tikzpicture}
\legende{Schéma de la structure d'opinion : sujet + verbe d'opinion + phrase complète, sans « que ».}
\end{popfigure}

\begin{attention}[Le piège du « 是 » et du « que » français]
Les francophones traduisent souvent « je pense \emph{que}... » par \zh{我觉得是...}. C'est une erreur : \zh{是} ne joue \textbf{pas} le rôle de « que ».\par
Dis : \zh{我觉得我们应该保护环境。} (Wǒ juéde wǒmen yīnggāi bǎohù huánjìng.)\par
Ne dis pas : \zh{我觉得是我们应该保护环境。}\par
Autre piège : ne laisse pas la phrase d'opinion sans sujet. En chinois, après \zh{觉得}, la phrase doit être complète : \zh{我觉得水污染很严重。} (Wǒ juéde shuǐ wūrǎn hěn yánzhòng. « Je pense que la pollution de l'eau est grave. »), et non \zh{觉得很严重} si l'on ne sait pas de quoi on parle.
\end{attention}

\begin{exemplebox}[Exemples traduits : donner son avis]
\zh{我认为种树很重要。}\\
\emph{Wǒ rènwéi zhòng shù hěn zhòngyào.}\\
« Je pense que planter des arbres est très important. »\\[0.4em]
\zh{我觉得雅温得的空气不太好。}\\
\emph{Wǒ juéde Yǎwēndé de kōngqì bú tài hǎo.}\\
« Je trouve que l'air de Yaoundé n'est pas très bon. »\\[0.4em]
\zh{我认为我们应该少用塑料袋。}\\
\emph{Wǒ rènwéi wǒmen yīnggāi shǎo yòng sùliào dài.}\\
« J'estime que nous devrions moins utiliser de sacs en plastique. »\\[0.4em]
\zh{他觉得保护环境是每个人的责任。}\\
\emph{Tā juéde bǎohù huánjìng shì měi ge rén de zérèn.}\\
« Il pense que protéger l'environnement est la responsabilité de chacun. »
\end{exemplebox}

Observe le dernier exemple : ici \zh{是} est correct, car il se trouve \emph{à l'intérieur} de la phrase complète qui suit \zh{觉得} (« protéger l'environnement \emph{est} la responsabilité de chacun »). Ce qui est interdit, c'est de placer \zh{是} immédiatement après \zh{觉得} pour traduire « que ».

\courssub{Approuver et réagir : 你说得对 et ses variantes}

Dans un dialogue, donner son avis ne suffit pas : il faut aussi \textbf{réagir} à celui de l'autre. Le dialogue nous offre la formule la plus utile : \zh{你说得对} (nǐ shuō de duì), littéralement « ce que tu dis est juste ». Observe sa construction : \zh{说} (dire) + \zh{得} (particule de manière) + \zh{对} (juste, correct). La particule \zh{得} introduit ici l'évaluation de l'action : « tu parles de manière juste ».

\begin{center}
\begin{tabularx}{\linewidth}{|X|X|X|}\hline
\rowcolor{popgreenL} Formule de réaction & Pinyin & Traduction et emploi \\ \hline
\zh{你说得对。} & nǐ shuō de duì & « Tu as raison » : accord poli et direct \\ \hline
\zh{是啊。} & shì a & « Oui / en effet » : accord simple, très oral \\ \hline
\zh{我同意。} & wǒ tóngyì & « Je suis d'accord » : accord plus formel \\ \hline
\zh{我也觉得。} & wǒ yě juéde & « Moi aussi je pense ainsi » \\ \hline
\zh{我不太同意。} & wǒ bú tài tóngyì & « Je ne suis pas tout à fait d'accord » : désaccord poli \\ \hline
\end{tabularx}
\end{center}

Note la politesse chinoise du désaccord : on dit rarement \zh{我不同意} (je ne suis pas d'accord) de façon sèche ; on adoucit avec \zh{不太} (pas tout à fait). Au Cameroun comme en Chine, adoucir son désaccord permet de garder une discussion amicale.

\courssub{Proposer des solutions : 应该、可以、不要、还要}

La fin du dialogue montre comment passer de l'opinion à l'action. Quatre petits mots-outils font tout le travail :

\begin{itemize}
\item \zh{应该} (yīnggāi, devoir) : pour dire ce qu'il faudrait faire. \zh{我们应该保护环境。} (Wǒmen yīnggāi bǎohù huánjìng.) « Nous devrions protéger l'environnement. »
\item \zh{可以} (kěyǐ, pouvoir) : pour demander ou indiquer une possibilité. \zh{我们可以做什么？} (Wǒmen kěyǐ zuò shénme ?) « Que pouvons-nous faire ? »
\item \zh{不要} (bú yào, ne pas) : pour conseiller d'éviter une action. \zh{我们不要扔垃圾。} (Wǒmen bú yào rēng lājī.) « Ne jetons pas de déchets. » Attention à la règle de sandhi : \zh{不} passe au deuxième ton devant \zh{要} (quatrième ton) : \emph{bú yào}.
\item \zh{还要} (hái yào, il faut encore / et aussi) : pour ajouter une deuxième action. \zh{还要多种树。} (Hái yào duō zhòng shù.) « Et plantons aussi plus d'arbres. »
\end{itemize}

\begin{methode}[Construire ta propre phrase d'avis en trois étapes]
\begin{enumerate}
\item Commence toujours par \zh{我觉得} ou \zh{我认为}.
\item Ajoute ensuite ta phrase complète, avec son sujet et son verbe : \zh{我觉得 + 水污染很严重。} (Wǒ juéde shuǐ wūrǎn hěn yánzhòng. « Je pense que la pollution de l'eau est grave. »).
\item Si tu proposes une solution, utilise \zh{应该} ou \zh{不要} : \zh{我认为我们应该多种树。} (Wǒ rènwéi wǒmen yīnggāi duō zhòng shù. « Je pense que nous devrions planter plus d'arbres. »).
\end{enumerate}
\end{methode}

\begin{exoresolu}[Traduire et compléter]
\textbf{Énoncé.} 1) Traduis en chinois : « Je pense que nous devrions protéger les rivières du Cameroun. » 2) Complète avec la bonne formule (\zh{我觉得}, \zh{你说得对}, \zh{不要}) : A : 水污染太严重了！B : \_\_\_\_\_\_，我们应该少扔垃圾。
\tcblower
\textbf{Solution détaillée.}\par
1) On applique la structure : sujet + verbe d'opinion + phrase complète, sans « que ». Le sujet est \zh{我}, le verbe d'opinion \zh{觉得} ou \zh{认为}, puis la phrase complète \zh{我们应该保护喀麦隆的河流}. On obtient : \zh{我觉得我们应该保护喀麦隆的河流。} (Wǒ juéde wǒmen yīnggāi bǎohù Kāmàilóng de héliú.) On n'ajoute surtout pas \zh{是} après \zh{觉得}.\par
2) A fait un constat (« la pollution de l'eau est très grave ») ; B approuve puis propose une solution. La formule d'accord attendue est donc \zh{你说得对} : \zh{你说得对，我们应该少扔垃圾。} (Nǐ shuō de duì, wǒmen yīnggāi shǎo rēng lājī.) « Tu as raison, nous devrions jeter moins de déchets. » \zh{我觉得} conviendrait grammaticalement mais ne marquerait pas l'accord avec A ; \zh{不要} ne peut pas ouvrir la réplique seul.
\end{exoresolu}

\begin{exercice}[À toi de parler]
\begin{enumerate}
\item Traduis en chinois : « Je trouve que l'air de Douala n'est pas très bon. »
\item Traduis en chinois : « Tu as raison, nous ne devrions pas jeter de déchets dans la rivière. »
\item Avec ton voisin, improvise un mini-dialogue de quatre répliques en suivant la progression : constat $\rightarrow$ avis $\rightarrow$ accord $\rightarrow$ proposition. Utilise au moins \zh{我觉得}, \zh{你说得对} et \zh{应该}.
\end{enumerate}
\end{exercice}

\begin{corrige}[Exercice : À toi de parler]
1) \zh{我觉得杜阿拉的空气不太好。} (Wǒ juéde Dù'ālā de kōngqì bú tài hǎo.) Structure : sujet + \zh{觉得} + phrase complète, sans « que ».\par
2) \zh{你说得对，我们不应该往河里扔垃圾。} (Nǐ shuō de duì, wǒmen bù yīnggāi wǎng hé lǐ rēng lājī.) On peut aussi dire plus simplement : \zh{你说得对，我们不要扔垃圾。}\par
3) Exemple de production attendue : A : \zh{校园里的垃圾太多了！} B : \zh{我觉得我们应该保护环境。} A : \zh{你说得对。} B : \zh{我们应该每天打扫校园。} (Měitiān dǎsǎo xiàoyuán : nettoyer la cour chaque jour.)
\end{corrige}

\begin{aretenir}[L'essentiel de la section]
\begin{itemize}
\item Pour donner ton avis : \textbf{sujet + 觉得 / 认为 + phrase complète}, sans « que » et sans \zh{是} juste après le verbe d'opinion.
\item Pour approuver : \zh{你说得对。} (tu as raison), \zh{是啊。} (en effet), \zh{我同意。} (je suis d'accord) ; pour un désaccord poli : \zh{我不太同意。}
\item Pour proposer : \zh{应该} (il faudrait), \zh{可以} (on peut), \zh{不要} (ne... pas), \zh{还要} (et aussi).
\item Structure-type d'un échange : constat $\rightarrow$ avis $\rightarrow$ accord $\rightarrow$ question $\rightarrow$ proposition.
\item Soigne les tons : \zh{垃圾} (lājī), \zh{保护} (bǎohù), \zh{不要} (bú yào).
\end{itemize}
\end{aretenir}

Tu maîtrises maintenant la mécanique de l'opinion. Dans la section suivante, nous élargirons ton répertoire avec toutes les \textbf{formules fonctionnelles} pour présenter un exposé, nuancer ton point de vue et relancer la parole dans un débat.

\coursec{Formules fonctionnelles : présenter, opiner, relancer}

Dans la section précédente, nous avons écouté et étudié le \emph{Dialogue modèle 1}, où deux élèves donnaient leur avis sur la protection de l'environnement. Vous avez remarqué que certaines phrases revenaient tout le temps : \zh{我觉得…}, \zh{我同意}, \zh{你呢？}… Ces phrases ne sont pas du vocabulaire ordinaire : ce sont des \cle{formules fonctionnelles}, c'est-à-dire des « outils de communication » tout prêts, que l'on réutilise dans n'importe quel exposé ou débat. Dans cette section, nous allons les classer, les expliquer une à une, puis apprendre à les enchaîner pour construire un vrai mini-exposé en chinois.

\courssub{Observer d'abord : une phrase d'ouverture}

\begin{textebox}[Phrase d'ouverture d'un exposé]
今天我要谈一谈环境保护。\\
Jīntiān wǒ yào tán yi tán huánjìng bǎohù.\\
« Aujourd'hui, je vais parler de la protection de l'environnement. »\\[0.5em]
环境保护很重要，我们应该保护我们的环境。\\
Huánjìng bǎohù hěn zhòngyào, wǒmen yīnggāi bǎohù wǒmen de huánjìng.\\
« La protection de l'environnement est très importante, nous devons protéger notre environnement. »
\end{textebox}

Observons la première phrase. Elle suit un plan très régulier :

\begin{center}
\zh{今天} (aujourd'hui) + \zh{我} (je) + \zh{要} (vouloir / aller) + \zh{谈一谈} (parler un peu de) + sujet.
\end{center}

Le mot \zh{要} (yào) exprime ici l'intention proche, comme le futur proche français « je vais parler ». La forme \zh{谈一谈} (tán yi tán) est un verbe redoublé avec \zh{一} au milieu : ce redoublement adoucit le verbe et signifie « parler un peu, dire deux mots de ». C'est une marque de modestie très appréciée en chinois : l'orateur ne prétend pas tout dire, il « en parle un peu ». On trouve la même structure avec \zh{说一说} (shuō yi shuō, « en dire un mot ») : \zh{我想说一说环境保护。} (Wǒ xiǎng shuō yi shuō huánjìng bǎohù. — « Je voudrais dire quelques mots sur la protection de l'environnement. »)

\begin{definition}[Formule fonctionnelle]
Une \cle{formule fonctionnelle} est une expression figée ou semi-figée qui remplit une fonction de communication précise : ouvrir un exposé, donner son avis, réagir, relancer, conclure. Elle s'apprend par cœur comme un tout, puis on remplace seulement la partie variable (le sujet, l'opinion).
\end{definition}

\courssub{Présenter un sujet}

\begin{propriete}[Présenter le sujet d'un exposé]
Structure : \zh{今天我要谈一谈 / 我想说一说} + \cle{sujet} + \zh{。}\\
Le sujet se place \emph{après} le verbe, sans préposition : contrairement au français « parler \emph{de} quelque chose », le chinois met directement le nom après \zh{谈} ou \zh{说}.\\
Exemples : \zh{今天我要谈一谈塑料污染。} (Jīntiān wǒ yào tán yi tán sùliào wūrǎn. — « Aujourd'hui, je vais parler de la pollution plastique. ») ; \zh{我想说一说雅温得的垃圾问题。} (Wǒ xiǎng shuō yi shuō Yǎwēndé de lājī wèntí. — « Je voudrais parler du problème des ordures à Yaoundé. »)
\end{propriete}

\begin{exemplebox}[Varier les sujets de présentation]
\zh{今天我要谈一谈水污染。} Jīntiān wǒ yào tán yi tán shuǐ wūrǎn. « Aujourd'hui, je vais parler de la pollution de l'eau. »\\
\zh{我想说一说砍树的问题。} Wǒ xiǎng shuō yi shuō kǎn shù de wèntí. « Je voudrais parler du problème de la coupe des arbres. »\\
\zh{今天我要谈一谈怎么保护我们的学校。} Jīntiān wǒ yào tán yi tán zěnme bǎohù wǒmen de xuéxiào. « Aujourd'hui, je vais parler de comment protéger notre école. »
\end{exemplebox}

\courssub{Donner son avis : 我觉得， 我认为， 在我看来}

Trois formules permettent d'exprimer une opinion, de la plus courante à la plus formelle :

\begin{center}
\begin{tabularx}{\linewidth}{|X|X|X|X|}
\hline
\rowcolor{popblueL} Formule & Pinyin & Registre & Équivalent français \\
\hline
我觉得… & wǒ juéde… & courant, oral & je trouve que…, je pense que… \\
\hline
我认为… & wǒ rènwéi… & soutenu, écrit et oral & j'estime que…, je considère que… \\
\hline
在我看来… & zài wǒ kànlái… & soutenu, ouvre une argumentation & à mon avis…, selon moi… \\
\hline
\end{tabularx}
\end{center}

\begin{propriete}[Structure de l'avis]
\zh{我觉得 / 我认为} + \cle{phrase complète} (sujet + verbe).\\
\zh{在我看来} se place en tête de phrase, suivi d'une virgule : \zh{在我看来，} + phrase.\\
Exemple : \zh{我觉得塑料污染是一个大问题。} (Wǒ juéde sùliào wūrǎn shì yí ge dà wèntí. — « Je trouve que la pollution plastique est un grand problème. ») ; \zh{在我看来，每个人都应该保护环境。} (Zài wǒ kànlái, měi ge rén dōu yīnggāi bǎohù huánjìng. — « À mon avis, chacun doit protéger l'environnement. »)
\end{propriete}

\begin{attention}[Ne pas traduire mot à mot « je pense que oui »]
Après \zh{觉得} et \zh{认为}, il faut une \emph{phrase complète}, jamais un mot seul. En français on dit « Je pense que oui » ; en chinois, on répète l'idée : \zh{我觉得是这样。} (Wǒ juéde shì zhèyàng. — « Je pense que c'est ainsi. ») et non pas \zh{觉得是}. Autre piège : \zh{觉得} exprime une opinion ou une impression, pas une sensation physique ; pour « j'ai chaud », on dit \zh{我很热} (Wǒ hěn rè), tandis que \zh{我觉得很热} signifie « je trouve qu'il fait chaud » (une impression, un jugement).
\end{attention}

\courssub{Être d'accord et ne pas être d'accord}

\begin{exemplebox}[Exprimer l'accord]
\zh{我同意。} Wǒ tóngyì. « Je suis d'accord. »\\
\zh{我同意你的看法。} Wǒ tóngyì nǐ de kànfǎ. « Je suis d'accord avec ton point de vue. »\\
\zh{你说得对。} Nǐ shuō de duì. « Tu as raison (ce que tu dis est juste). »\\
\zh{我也是这么想的。} Wǒ yě shì zhème xiǎng de. « Je le pense aussi (c'est aussi ce que je pense). »
\end{exemplebox}

\begin{attention}[同意 se construit directement]
\zh{同意} (tóngyì) est un verbe qui prend directement son complément : \zh{我同意你} (je suis d'accord avec toi), \zh{我同意你的看法} (je suis d'accord avec ton avis). N'ajoutez \emph{jamais} \zh{是} : on ne dit pas « 我同意是 ». De même, \zh{你说得对} utilise \zh{得} (et non \zh{的}) : c'est la structure verbe + 得 + adjectif, « tu parles de façon juste ».
\end{attention}

Pour le désaccord, le chinois préfère, comme le français poli, adoucir la contradiction. On commence souvent par \zh{可是} (kěshì, « mais ») ou on reformule avec \zh{我觉得} :

\begin{exemplebox}[Exprimer le désaccord poliment]
\zh{我不同意。} Wǒ bù tóngyì. « Je ne suis pas d'accord. »\\
\zh{我觉得不是这样。} Wǒ juéde bú shì zhèyàng. « Je ne pense pas que ce soit ainsi. »\\
\zh{可是塑料袋很方便。} Kěshì sùliàodài hěn fāngbiàn. « Mais les sacs plastiques sont très pratiques. »\\
\zh{我不同意，我觉得塑料袋很方便。} Wǒ bù tóngyì, wǒ juéde sùliàodài hěn fāngbiàn. « Je ne suis pas d'accord, je trouve les sacs plastiques pratiques. »
\end{exemplebox}

Remarquez la négation : \zh{不同意} prend \zh{不} (bù), qui passe au deuxième ton devant un quatrième ton : \emph{bú tóngyì}. De même \zh{不是} se prononce \emph{bú shì}.

\courssub{Relancer la parole : 呢 et les questions d'opinion}

Un débat n'est pas un monologue : après avoir parlé, il faut \cle{relancer} son camarade. Le chinois possède un outil très économique pour cela : la particule \zh{呢}.

\begin{propriete}[La particule 呢 pour relancer]
\zh{呢} (ne, ton neutre) en fin de phrase reprend le contexte sans répéter la question : \zh{你呢？} (Nǐ ne ?) = « Et toi ? ». Après avoir dit \zh{我觉得应该少用塑料袋} (Wǒ juéde yīnggāi shǎo yòng sùliàodài. — « Je pense qu'on devrait utiliser moins de sacs plastiques. »), il suffit d'ajouter \zh{你呢？} pour demander « et toi, qu'en penses-tu ? ».\\
Formes plus explicites : \zh{你怎么看？} (Nǐ zěnme kàn ? — « Qu'en penses-tu ? ») ; \zh{你怎么看这个问题？} (Nǐ zěnme kàn zhège wèntí ? — « Comment vois-tu ce problème ? ») ; \zh{你同意吗？} (Nǐ tóngyì ma ? — « Es-tu d'accord ? »)
\end{propriete}

\begin{attention}[呢 n'est pas 吗]
\zh{吗} (ma) transforme une phrase affirmative en question fermée (oui/non) : \zh{你同意吗？} — « Es-tu d'accord ? ». \zh{呢} (ne) ne forme pas une question nouvelle : il \emph{rappelle} la question précédente en changeant de personne. On ne peut donc pas commencer un dialogue par \zh{你呢？} ; il faut qu'une information ou une question ait déjà été donnée. Ordre des mots : dans \zh{你怎么看这个问题？}, le mot interrogatif \zh{怎么} (comment) se place \emph{avant} le verbe, et l'ordre de la phrase ne change pas — pas d'inversion comme en français.
\end{attention}

\courssub{Conclure un exposé}

Deux formules suffisent pour terminer proprement :

\begin{exemplebox}[Conclure]
\zh{所以，我们应该少用塑料袋，多种树。} Suǒyǐ, wǒmen yīnggāi shǎo yòng sùliàodài, duō zhòng shù. « Donc, nous devrions utiliser moins de sacs plastiques et planter plus d'arbres. »\\
\zh{谢谢大家！} Xièxie dàjiā ! « Merci à tous ! »
\end{exemplebox}

\zh{所以} (suǒyǐ, « donc, c'est pourquoi ») introduit la conclusion logique ; \zh{应该} (yīnggāi, « devoir » au sens du conseil) introduit la solution ; \zh{谢谢大家} est la formule de politesse finale obligatoire après un exposé en classe.

\begin{methode}[Structure d'un mini-exposé en 4 phrases]
\begin{enumerate}
\item \textbf{Présenter le sujet} : \zh{今天我要谈一谈…} (Jīntiān wǒ yào tán yi tán…).
\item \textbf{Décrire le problème} : une phrase simple avec \zh{有} ou un adjectif : \zh{我们的城市有很多垃圾。} (Wǒmen de chéngshì yǒu hěn duō lājī. — « Notre ville a beaucoup d'ordures. »)
\item \textbf{Donner son avis} avec \zh{我觉得} : \zh{我觉得这是一个大问题。} (Wǒ juéde zhè shì yí ge dà wèntí. — « Je trouve que c'est un grand problème. »)
\item \textbf{Proposer une solution} avec \zh{应该} et conclure : \zh{所以，我们应该保护环境。谢谢大家！} (Suǒyǐ, wǒmen yīnggāi bǎohù huánjìng. Xièxie dàjiā ! — « Donc, nous devons protéger l'environnement. Merci à tous ! »)
\end{enumerate}
\end{methode}

\begin{popfigure}
\begin{tikzpicture}[
  node distance=0.55cm,
  box/.style={draw=popblue, fill=popblueL, rounded corners=3pt, align=center, minimum width=4.6cm, minimum height=1.05cm, font=\small},
  arr/.style={-{Stealth[length=2.5mm]}, thick, popdark}
]
\node[box, align=center] (n1){\textbf{Présenter}\\今天我要谈一谈…};
\node[box, below=of n1, align=center] (n2){\textbf{Donner son avis}\\我觉得… / 我认为…};
\node[box, below=of n2, align=center] (n3){\textbf{Réagir}\\我同意 / 我不同意};
\node[box, below=of n3, align=center] (n4){\textbf{Relancer}\\你呢？ / 你怎么看？};
\node[box, below=of n4, align=center] (n5){\textbf{Conclure}\\所以，我们应该… 谢谢大家！};
\draw[arr] (n1) -- (n2);
\draw[arr] (n2) -- (n3);
\draw[arr] (n3) -- (n4);
\draw[arr] (n4.west) .. controls +(-1.4,0.4) and +(-1.4,-0.4) .. (n2.west);
\draw[arr] (n4) -- (n5);
\end{tikzpicture}
\legende{Organigramme du débat : la flèche de retour montre que l'on peut relancer plusieurs fois avant de conclure.}
\end{popfigure}

\courssub{Tableau récapitulatif des formules}

\begin{center}
\begin{tabularx}{\linewidth}{|X|X|X|X|}
\hline
\rowcolor{popblueL} Fonction & Formule & Pinyin & Traduction \\
\hline
Présenter & 今天我要谈一谈… & Jīntiān wǒ yào tán yi tán… & Aujourd'hui, je vais parler de… \\
\hline
Présenter & 我想说一说… & Wǒ xiǎng shuō yi shuō… & Je voudrais dire un mot de… \\
\hline
Opiner & 我觉得… & wǒ juéde… & je trouve que… \\
\hline
Opiner & 我认为… & wǒ rènwéi… & j'estime que… \\
\hline
Opiner & 在我看来，… & zài wǒ kànlái, … & à mon avis, … \\
\hline
Accord & 我同意（你的看法）。 & Wǒ tóngyì (nǐ de kànfǎ). & Je suis d'accord (avec ton avis). \\
\hline
Accord & 你说得对。 & Nǐ shuō de duì. & Tu as raison. \\
\hline
Accord & 我也是这么想的。 & Wǒ yě shì zhème xiǎng de. & Je le pense aussi. \\
\hline
Désaccord & 我不同意。 & Wǒ bù tóngyì. & Je ne suis pas d'accord. \\
\hline
Désaccord & 我觉得不是这样。 & Wǒ juéde bú shì zhèyàng. & Je ne pense pas que ce soit ainsi. \\
\hline
Relancer & 你呢？ & Nǐ ne ? & Et toi ? \\
\hline
Relancer & 你怎么看（这个问题）？ & Nǐ zěnme kàn (zhège wèntí) ? & Qu'en penses-tu ? \\
\hline
Relancer & 你同意吗？ & Nǐ tóngyì ma ? & Es-tu d'accord ? \\
\hline
Conclure & 所以，我们应该… & Suǒyǐ, wǒmen yīnggāi… & Donc, nous devrions… \\
\hline
Conclure & 谢谢大家！ & Xièxie dàjiā ! & Merci à tous ! \\
\hline
\end{tabularx}
\end{center}

\begin{exoresolu}[Compléter un échange avec la bonne formule]
Énoncé — Complétez ce mini-débat entre Awa et Lin avec les formules étudiées : \zh{我觉得} / \zh{你呢} / \zh{我不同意} / \zh{所以}.\\
Awa : 今天我要谈一谈塑料袋的问题。\_\_\_\_\_\_ 塑料袋污染环境，我们应该少用。\\
Awa : \_\_\_\_\_\_ ？\\
Lin : \_\_\_\_\_\_ ，我觉得塑料袋很方便。\\
Awa : 好吧！\_\_\_\_\_\_ ，我们应该少用，也应该回收。
\tcblower
Solution détaillée — 1) Awa donne son avis : \zh{我觉得塑料袋污染环境…} (Wǒ juéde sùliàodài wūrǎn huánjìng. — « Je trouve que les sacs plastiques polluent l'environnement. »). 2) Elle relance Lin : \zh{你呢？} (Nǐ ne ?). 3) Lin n'est pas d'accord : \zh{我不同意，我觉得塑料袋很方便。} (Wǒ bù tóngyì, wǒ juéde sùliàodài hěn fāngbiàn.). 4) Awa conclut : \zh{所以，我们应该少用，也应该回收。} (Suǒyǐ, wǒmen yīnggāi shǎo yòng, yě yīnggāi huíshōu. — « Donc, nous devrions en utiliser moins et aussi recycler. »). Vérifiez la logique : avis → relance → réaction → conclusion, exactement comme dans l'organigramme.
\end{exoresolu}

\begin{exercice}[À vous de parler]
Préparez oralement un mini-exposé de 4 phrases sur le thème « 保护我们的学校 » (bǎohù wǒmen de xuéxiào, protéger notre école) en suivant la méthode : présentation, problème, avis avec \zh{我觉得}, solution avec \zh{应该}. Terminez par \zh{谢谢大家！}, puis relancez votre voisin avec \zh{你呢？}
\end{exercice}

\begin{corrige}[Exercice « À vous de parler »]
Modèle de production attendue (les réponses peuvent varier) :\\
\zh{今天我要谈一谈保护我们的学校。我们的学校有很多垃圾。我觉得这不漂亮，也不健康。所以，我们应该每天打扫，少用塑料袋。谢谢大家！你呢？}\\
Jīntiān wǒ yào tán yi tán bǎohù wǒmen de xuéxiào. Wǒmen de xuéxiào yǒu hěn duō lājī. Wǒ juéde zhè bù piàoliang, yě bù jiànkāng. Suǒyǐ, wǒmen yīnggāi měitiān dǎsǎo, shǎo yòng sùliàodài. Xièxie dàjiā ! Nǐ ne ?\\
« Aujourd'hui, je vais parler de la protection de notre école. Il y a beaucoup d'ordures dans notre école. Je trouve que ce n'est ni beau ni sain. Donc, nous devrions nettoyer chaque jour et utiliser moins de sacs plastiques. Merci à tous ! Et toi ? »\\
Critères de réussite : 4 phrases au minimum, une formule de chaque fonction, tons corrects sur \zh{我觉得} (wǒ juéde : 3-2) et \zh{应该} (yīnggāi : 1-1).
\end{corrige}

\begin{aretenir}[L'essentiel de la section]
\begin{itemize}
\item Présenter : \zh{今天我要谈一谈…} ; opiner : \zh{我觉得 / 我认为 / 在我看来}.
\item Accord : \zh{我同意（你的看法）}, \zh{你说得对} ; désaccord poli : \zh{我不同意}, \zh{我觉得不是这样}, \zh{可是…}.
\item Relancer : \zh{你呢？} (particule \zh{呢} qui rappelle la question), \zh{你怎么看？}, \zh{你同意吗？}
\item Conclure : \zh{所以，我们应该…} + \zh{谢谢大家！}
\item Pièges : pas de \zh{是} après \zh{同意} ; \zh{得} dans \zh{说得对} ; \zh{不} devient \emph{bú} devant un 4\textsuperscript{e} ton.
\end{itemize}
\end{aretenir}

\coursec{Vocabulaire clé : l'environnement}

Cette leçon mobilise un lexique précis pour décrire la nature, les menaces qui pèsent sur elle et les actions de protection. Maîtriser ces mots est la condition \textbf{sine qua non} pour réussir l'exposé et le débat oral. Nous classons le vocabulaire par champs sémantiques pour en faciliter la mémorisation.

\begin{textebox}[Lexique thématique : Environnement et protection]
\textbf{1. Le milieu naturel (大自然)}\\
\zh{环境} (huánjìng) \textbf{environnement} \quad \zh{自然} (zìrán) \textbf{nature} \quad \zh{空气} (kōngqì) \textbf{air} \quad \zh{水} (shuǐ) \textbf{eau} \quad \zh{树} (shù) \textbf{arbre} \quad \zh{森林} (sēnlín) \textbf{forêt} \quad \zh{河} (hé) \textbf{rivière} \quad \zh{天空} (tiānkōng) \textbf{ciel}\\[6pt]
\textbf{2. Problèmes écologiques (环境问题)}\\
\zh{污染} (wūrǎn) \textbf{pollution (verbe/nom)} \quad \zh{白色污染} (báisè wūrǎn) \textbf{pollution blanche (sacs plastiques)} \quad \zh{垃圾} (lājī) \textbf{déchets, ordures} \quad \zh{废气} (fèiqì) \textbf{gaz d'échappement} \quad \zh{噪音} (zàoyīn) \textbf{bruit, nuisance sonore} \quad \zh{砍伐} (kǎnfá) \textbf{abattre (arbres)} \quad \zh{浪费} (làngfèi) \textbf{gaspiller} \quad \zh{全球变暖} (quánqiú biǎnnuǎn) \textbf{réchauffement climatique}\\[6pt]
\textbf{3. Actions de protection (保护行动)}\\
\zh{保护} (bǎohù) \textbf{protéger} \quad \zh{节约} (jiéyuē) \textbf{économiser} \quad \zh{回收} (huíshōu) \textbf{recycler, récupérer} \quad \zh{分类} (fēnlèi) \textbf{trier (déchets)} \quad \zh{种树} (zhòng shù) \textbf{planter des arbres} \quad \zh{减少} (jiǎnshǎo) \textbf{réduire} \quad \zh{推广} (tuīguǎng) \textbf{promouvoir} \quad \zh{自带} (zìdài) \textbf{apporter soi-même}\\[6pt]
\textbf{4. Alternatives écologiques (替代品)}\\
\zh{塑料袋} (sùliàodài) \textbf{sac plastique} \quad \zh{纸袋} (zhǐdài) \textbf{sac en papier} \quad \zh{布袋} (bùdài) \textbf{sac en tissu} \quad \zh{一次性筷子} (yīcìxìng kuàizi) \textbf{baguettes jetables} \quad \zh{竹筷子} (zhú kuàizi) \textbf{baguettes en bambou} \quad \zh{金属筷子} (jīnshǔ kuàizi) \textbf{baguettes en métal} \quad \zh{公交车} (gōngjiāochē) \textbf{bus} \quad \zh{自行车} (zìxíngchē) \textbf{vélo}
\end{textebox}

\begin{center}
\begin{tabularx}{\linewidth}{|X|X|X|X|}
\hline
\rowcolor{popblueL} \textbf{Caractères} & \textbf{Pinyin} & \textbf{Nature} & \textbf{Traduction} \\
\hline
环境 & huánjìng & nom & environnement \\
\hline
污染 & wūrǎn & verbe/nom & polluer / pollution \\
\hline
保护 & bǎohù & verbe & protéger \\
\hline
塑料袋 & sùliàodài & nom & sac en plastique \\
\hline
纸袋 & zhǐdài & nom & sac en papier \\
\hline
布袋 & bùdài & nom & sac en tissu \\
\hline
一次性 & yīcìxìng & adj. & à usage unique, jetable \\
\hline
节约 & jiéyuē & verbe & économiser \\
\hline
回收 & huíshōu & verbe & recycler \\
\hline
分类 & fēnlèi & verbe & trier / classer \\
\hline
推广 & tuīguǎng & verbe & promouvoir, vulgariser \\
\hline
白色污染 & báisè wūrǎn & nom & pollution blanche \\
\hline
可降解 & kě jiàngjiě & adj. & biodégradable \\
\hline
方便 & fāngbiàn & adj. & pratique, commode \\
\hline
便宜 & piányi & adj. & bon marché \\
\hline
\end{tabularx}
\end{center}

\begin{aretenir}[Écriture des caractères clés : Radicaux et ordre des traits]
\begin{itemize}
    \item \zh{环} (huán) : Clé \zh{王} (jade/roi) + \zh{亢}. Ordre : haut → bas (王), puis structure enveloppante. Sens : « anneau, environnement ».
    \item \zh{污} (wū) : Clé \zh{氵} (trois points d'eau). Phonétique \zh{于}. L'eau sale → pollution.
    \item \zh{护} (hù) : Clé \zh{扌} (main). Phonétique \zh{戶}. La main qui protège la porte/maison.
    \item \zh{塑} (sù) : Clé \zh{土} (terre). Phonétique \zh{溯}. Terre modelée → plastique.
    \item \zh{纸} (zhǐ) : Clé \zh{纟} (fil/soie). Phonétique \zh{氏}. Fibres végétales → papier.
    \item \zh{布} (bù) : Clé \zh{巾} (étoffe). Tissu, toile.
    \item \zh{竹} (zhú) : Clé \zh{竹} (bambou). Pictogramme de pousses de bambou.
\end{itemize}
\cle{Rappel ordre des traits :} Horizontal avant vertical ; Gauche avant droite ; Haut avant bas ; Extérieur avant intérieur ; Traits centrales avant symétriques ; Point final à droite.
\end{aretenir}

\begin{attention}[Classificateurs (量词) indispensables pour ce thème]
En chinois, tout nom comptable nécessite un classificateur entre le nombre/démonstratif et le nom.
\begin{center}
\begin{tabularx}{\linewidth}{|X|X|X|}
\hline
\rowcolor{popblueL} \textbf{Classificateur} & \textbf{Noms concernés (thème)} & \textbf{Exemple} \\
\hline
\zh{个} (gè) & universel (problème, idée, argument) & \zh{一个论点} (un argument) \\
\hline
\zh{个} / \zh{棵} (kē) & \zh{树} (arbre), \zh{竹子} & \zh{一棵树}, \zh{种两棵树} \\
\hline
\zh{个} / \zh{条} (tiáo) & \zh{河} (rivière), \zh{路} (route), \zh{筷子} (baguettes - long) & \zh{一条河}, \zh{一双筷子} (\zh{双} pour la paire) \\
\hline
\zh{个} / \zh{只} (zhī) & \zh{袋子} (sac - conteneur), \zh{动物} & \zh{一个袋子}, \zh{一只布袋} \\
\hline
\zh{张} (zhāng) & \zh{纸} (feuille de papier), \zh{桌子}, \zh{票} & \zh{一张纸}, \zh{几张纸袋} (si plat) \\
\hline
\zh{辆} (liàng) & \zh{车} (voiture, bus, vélo) & \zh{一辆公交车}, \zh{两辆自行车} \\
\hline
\end{tabularx}
\end{center}
\cle{Erreur fréquente :} Dire \zh{*一个筷子} (incorrect). On dit \zh{一双筷子} (une paire) ou \zh{一把筷子} (une poignée/quantité). Pour un seul bâton : \zh{一根筷子} (\zh{根} pour objets longs et fins).
\end{attention}

\coursec{Dialogue modèle 1 : donner son avis}

Avant d'aborder le débat contradictoire, il faut savoir présenter un point de vue simple et structuré. Ce dialogue met en scène deux lycéens, \zh{李明} (Lǐ Míng) et \zh{王芳} (Wáng Fāng), membres du club « Terre Verte » de leur lycée à Yaoundé. Ils préparent une affiche pour la \zh{世界环境日} (Journée mondiale de l'environnement, 5 juin).

\begin{textebox}[Dialogue : Préparer une affiche pour la Journée de l'environnement]
\textbf{A :} 李明，世界环境日快到了。我们做什么？\\
Lǐ Míng, Shìjiè Huánjìng Rì kuài dào le. Wǒmen zuò shénme?\\
« Li Ming, la Journée de l'environnement approche. Qu'est-ce qu'on fait ? »\\[4pt]
\textbf{B :} 我认为我们应该做一个海报。\\
Wǒ rènwéi wǒmen yīnggāi zuò yī gè hǎibào.\\
« Je pense qu'on devrait faire une affiche. »\\[4pt]
\textbf{A :} 好主意！海报上写什么？\\
Hǎo zhǔyi! Hǎibào shàng xiě shénme?\\
« Bonne idée ! Qu'est-ce qu'on écrit sur l'affiche ? »\\[4pt]
\textbf{B :} 写“保护环境，从我做起”。\\
Xiě “Bǎohù huánjìng, cóng wǒ zuòqǐ”.\\
« Écrire "Protéger l'environnement, commencer par moi-même". »\\[4pt]
\textbf{A :} 我也觉得很好。我们还可以写“少用塑料袋，多用布袋”。\\
Wǒ yě juéde hěn hǎo. Wǒmen hái kěyǐ xiě “Shǎo yòng sùliàodài, duō yòng bùdài”.\\
« Moi aussi je trouve ça très bien. On peut aussi écrire "Moins de sacs plastiques, plus de sacs en tissu". »\\[4pt]
\textbf{B :} 同意。那我们现在开始画吧。\\
Tóngyì. Nà wǒmen xiànzài kāishǐ huà ba.\\
« D'accord. Alors on commence à dessiner maintenant. »
\end{textebox}

\begin{center}
\begin{tabularx}{\linewidth}{|X|X|X|X|}
\hline
\rowcolor{popblueL} \textbf{Caractères} & \textbf{Pinyin} & \textbf{Nature} & \textbf{Traduction} \\
\hline
世界环境日 & Shìjiè Huánjìng Rì & n. propre & Journée mondiale de l'environnement \\
\hline
海报 & hǎibào & nom & affiche, poster \\
\hline
从…做起 & cóng… zuòqǐ & locution & commencer par…, s'initier par… \\
\hline
少 & shǎo & adv/verbe & peu, réduire, moins \\
\hline
多 & duō & adv/verbe & beaucoup, augmenter, plus \\
\hline
同意 & tóngyì & verbe & être d'accord \\
\hline
开始 & kāishǐ & verbe & commencer \\
\hline
画 & huà & verbe & dessiner, peindre \\
\hline
吧 & ba & particule & suggestion, adoucissement \\
\hline
\end{tabularx}
\end{center}

\begin{propriete}[Structures pour exprimer son avis : \zh{我认为} vs \zh{我觉得}]
\begin{itemize}
    \item \zh{我认为} (Wǒ rènwéi) : « Je pense que / J'estime que ». \textbf{Registre neutre à formel}. Basé sur un raisonnement, une réflexion. Utilisé pour ouvrir un exposé ou un débat.
    \item \zh{我觉得} (Wǒ juéde) : « Je trouve que / J'ai l'impression que ». \textbf{Registre oral, subjectif}. Basé sur un ressenti, une impression personnelle.
\end{itemize}
\textbf{Structure commune :} \zh{主语 + 认为/觉得 + [Proposition complète]}
\begin{center}
\begin{tabularx}{\linewidth}{|X|X|X|}
\hline
\rowcolor{popblueL} \textbf{Structure} & \textbf{Exemple} & \textbf{Traduction} \\
\hline
\zh{我认为 + Suj + 应该 + V} & \zh{我认为我们应该保护环境。} & Je pense que nous devrions protéger l'environnement. \\
\hline
\zh{我觉得 + Adj / SV} & \zh{我觉得这个主意很好。} & Je trouve que cette idée est très bien. \\
\hline
\zh{我觉得 + Suj + V + 得 + Adv} & \zh{我觉得他说得很有道理。} & Je trouve que ce qu'il dit est très sensé. \\
\hline
\end{tabularx}
\end{center}
\end{propriete}

\begin{propriete}[Structure \zh{少 / 多 + Verbe} : « Moins / Plus de… »]
Pour formuler des recommandations concrètes (éco-gestes), on place les adverbes \zh{少} (peu/moins) et \zh{多} (beaucoup/plus) \textbf{devant le verbe d'action}.
\begin{itemize}
    \item \zh{少用塑料袋} (Shǎo yòng sùliàodài) — Utiliser moins de sacs plastiques.
    \item \zh{多用布袋} (Duō yòng bùdài) — Utiliser plus de sacs en tissu.
    \item \zh{少开车，多坐公交} (Shǎo kāichē, duō zuò gōngjiāo) — Moins conduire, plus prendre le bus.
    \item \zh{少说多做} (Shǎo shuō duō zuò) — Moins parler, plus agir (idiome).
\end{itemize}
\cle{Attention :} Ne pas confondre \zh{少} (adverbe « moins ») et \zh{小} (adjectif « petit »). \zh{*小用} est faux.
\end{propriete}

\coursec{Formules fonctionnelles : présenter, opiner, relancer}

Pour réussir l'expression orale (exposé + débat), vous devez maîtriser ces « briques » communicatives. Elles sont votre \textbf{boîte à outils} pour l'examen oral.

\begin{definition}[Boîte à outils orale : 4 fonctions clés]
\textbf{1. PRÉSENTER un exposé (Ouvrir / Structurer / Clore)}
\begin{center}
\begin{tabularx}{\linewidth}{|X|X|X|}
\hline
\rowcolor{popblueL} \textbf{Fonction} & \textbf{Formule chinoise + Pinyin} & \textbf{Traduction} \\
\hline
Ouvrir & \zh{大家好，今天我要谈谈…} (Dàjiā hǎo, jīntiān wǒ yào tántan…) & Bonjour à tous, aujourd'hui je vais parler de… \\
\hline
Ouvrir & \zh{我的题目是…} (Wǒ de tímù shì…) & Mon sujet est… \\
\hline
Structurer & \zh{首先…，其次…，最后…} (Shǒuxiān…, qícì…, zuìhòu…) & D'abord…, ensuite…, enfin… \\
\hline
Donner ex. & \zh{比如说…} (Bǐrú shuō…) & Par exemple… \\
\hline
Clore & \zh{总之…} (Zǒngzhī…) & En résumé / En un mot… \\
\hline
Clore & \zh{谢谢大家。} (Xièxie dàjiā.) & Merci à tous. \\
\hline
\end{tabularx}
\end{center}

\textbf{2. DONNER SON AVIS / PROPOSER}
\begin{center}
\begin{tabularx}{\linewidth}{|X|X|X|}
\hline
\rowcolor{popblueL} \textbf{Nuance} & \textbf{Formule} & \textbf{Traduction} \\
\hline
Opinion raisonnée & \zh{我认为 / 据我看…} (Wǒ rènwéi / Jù wǒ kàn…) & Je pense que / À mon avis… \\
\hline
Sentiment & \zh{我觉得…} (Wǒ juéde…) & Je trouve que / J'ai l'impression que… \\
\hline
Suggestion forte & \zh{我建议…} (Wǒ jiànyì…) & Je suggère / propose que… \\
\hline
Obligation morale & \zh{我们应该 / 得…} (Wǒmen yīnggāi / děi…) & Nous devons / Il faut que nous… \\
\hline
\end{tabularx}
\end{center}

\textbf{3. EXPRIMER ACCORD / DÉSACCORD}
\begin{center}
\begin{tabularx}{\linewidth}{|X|X|X|}
\hline
\rowcolor{popblueL} \textbf{Position} & \textbf{Formule} & \textbf{Traduction} \\
\hline
Accord total & \zh{我同意。/ 我赞成。} (Wǒ tóngyì. / Wǒ zànchéng.) & Je suis d'accord. / J'approuve. \\
\hline
Accord partiel & \zh{你说得有道理，但是…} (Nǐ shuō de yǒu dàolǐ, dànshì…) & Tu as raison, mais… \\
\hline
Désaccord poli & \zh{我有不同的看法。} (Wǒ yǒu bùtóng de kànfǎ.) & J'ai un avis différent. \\
\hline
Désaccord direct & \zh{我不同意。/ 我不赞成。} (Wǒ bù tóngyì. / Wǒ bù zànchéng.) & Je ne suis pas d'accord. / Je n'approuve pas. \\
\hline
\end{tabularx}
\end{center}

\textbf{4. RELANCER / CONCÉDER / NUANCER (Le "liant" du débat)}
\begin{center}
\begin{tabularx}{\linewidth}{|X|X|X|}
\hline
\rowcolor{popblueL} \textbf{Fonction} & \textbf{Formule} & \textbf{Traduction} \\
\hline
Relancer & \zh{你怎么看？} (Nǐ zěnme kàn?) & Qu'en penses-tu ? / Quel est ton avis ? \\
\hline
Relancer & \zh{还有什么建议？} (Hái yǒu shénme jiànyì?) & D'autres suggestions ? \\
\hline
Concéder & \zh{你说得也有道理。} (Nǐ shuō de yě yǒu dàolǐ.) & Ce que tu dis a aussi du sens. \\
\hline
Concéder & \zh{有道理，不过…} (Yǒu dàolǐ, bùguò…) & C'est sensé, toutefois… \\
\hline
Transition & \zh{那 / 那么…} (Nà / Nàme…) & Alors / Dans ce cas… \\
\hline
Proposer alt. & \zh{我们可以…} (Wǒmen kěyǐ…) & Nous pouvons… \\
\hline
Valider & \zh{好主意！/ 就这么定！} (Hǎo zhǔyi! / Jiù zhème dìng!) & Bonne idée ! / C'est décidé comme ça ! \\
\hline
\end{tabularx}
\end{center}
\end{definition}

\begin{methode}[Méthode : Préparer un mini-exposé de 1 minute (3 étapes)]
\begin{enumerate}
    \item \textbf{Planifier (3 min) :} Choisir UN sujet précis (ex: \zh{垃圾分类}). Noter 3 arguments en mots-clés (caractères) : \zh{1. 减少污染 2. 变废为宝 3. 养成习惯}. Utiliser \zh{首先、其次、最后}.
    \item \textbf{Rédiger le script (5 min) :} Écrire des phrases complètes avec pinyin. Vérifier : \zh{我认为…} (ouverture) + \zh{首先…} + \zh{其次…} + \zh{最后…} + \zh{总之…} (clôture) + \zh{谢谢大家}.
    \item \textbf{S'entraîner à l'oral (5 min) :} Lire à voix haute 3 fois. \textbf{Focus tons :} \zh{环境} (2-4), \zh{污染} (1-3), \zh{应该} (1-1), \zh{垃圾} (1-1). S'enregistrer si possible.
\end{enumerate}
\end{methode}

\coursec{Dialogue modèle 2 : mini-débat}

Après avoir maîtrisé les formules fonctionnelles pour présenter un exposé, exprimer son opinion et relancer l'échange, nous passons maintenant à la mise en situation authentique : un mini-débat contradictoire. Dans la section précédente, nous avons appris à dire « \zh{我认为…} » (je pense que…) ou « \zh{我赞成} » (j'approuve). Ici, nous allons observer comment deux interlocuteurs gèrent un désaccord poli, argumentent et trouvent un compromis sur un thème concret : l'interdiction des sacs plastiques au marché. C'est une situation courante, tant à Yaoundé qu'à Pékin, qui mobilise du vocabulaire HSK 2-3 et des structures de liaison logiques indispensables à l'oral.

\courssub{Observation et analyse du dialogue modèle}

Écoutons et lisons attentivement le dialogue suivant entre deux camarades, \zh{小李} (Xiǎo Lǐ) et \zh{小王} (Xiǎo Wáng). Repérez les marqueurs d'opposition (\zh{可是}, \zh{但是}), la formule de concession polie (\zh{你说得也有道理}) et la proposition de solution alternative.

\begin{textebox}[Dialogue : Faut-il interdire les sacs plastiques au marché ?]
\textbf{A :} 我认为我们应该不用塑料袋。\\
Wǒ rènwéi wǒmen yīnggāi bú yòng sùliàodài.\\
« Je pense que nous ne devrions plus utiliser de sacs plastiques. »\\[4pt]
\textbf{B :} 我不同意。塑料袋很方便。\\
Wǒ bù tóngyì. Sùliàodài hěn fāngbiàn.\\
« Je ne suis pas d'accord. Les sacs plastiques sont pratiques. »\\[4pt]
\textbf{A :} 可是塑料袋污染环境。\\
Kěshì sùliàodài wūrǎn huánjìng.\\
« Mais les sacs plastiques polluent l'environnement. »\\[4pt]
\textbf{B :} 你说得也有道理。那我们可以用纸袋。\\
Nǐ shuō de yě yǒu dàolǐ. Nà wǒmen kěyǐ yòng zhǐdài.\\
« Ce que tu dis a aussi du sens. Alors nous pouvons utiliser des sacs en papier. »\\[4pt]
\textbf{A :} 好主意！\\
Hǎo zhǔyi!\\
« Bonne idée ! »
\end{textebox}

\begin{center}
\begin{tabularx}{\linewidth}{|X|X|X|X|}
\hline
\rowcolor{popblueL} \textbf{Caractères} & \textbf{Pinyin} & \textbf{Nature} & \textbf{Traduction} \\
\hline
塑料袋 & sùliàodài & nom & sac en plastique \\
\hline
方便 & fāngbiàn & adj. & pratique, commode \\
\hline
污染 & wūrǎn & verbe & polluer \\
\hline
环境 & huánjìng & nom & environnement \\
\hline
道理 & dàolǐ & nom & raison, sens, principe \\
\hline
纸袋 & zhǐdài & nom & sac en papier \\
\hline
主意 & zhǔyi & nom & idée, avis, décision \\
\hline
不同意 & bù tóngyì & verbe & ne pas être d'accord \\
\hline
应该 & yīnggāi & verbe modal & devoir, devrait \\
\hline
不用 & bú yòng & verbe modal (nég.) & ne pas utiliser / ne pas avoir besoin de \\
\hline
\end{tabularx}
\end{center}

\courssub{Point de grammaire 1 : Marquer l'opposition avec \zh{可是} et \zh{但是}}

\begin{propriete}[Les connecteurs d'opposition \zh{可是} (kěshì) et \zh{但是} (dànshì)]
En chinois, pour introduire une objection, une concession ou un contraste (l'équivalent de « mais », « cependant », « pourtant »), on place le connecteur en \textbf{début de proposition} (souvent après le sujet s'il y en a un, ou tout au début de la phrase).

\textbf{Structure :} [Sujet] + \zh{可是 / 但是} + [Proposition contrastive]

\textbf{Distinction d'emploi :}
\begin{itemize}
    \item \zh{可是} (kěshì) : registre \textbf{oral}, courant, plus doux. Idéal pour un dialogue entre camarades.
    \item \zh{但是} (dànshì) : registre \textbf{plus formel}, écrit ou soutenu. Marque une opposition plus forte.
\end{itemize}

\textbf{Exemples comparatifs :}
\begin{center}
\begin{tabularx}{\linewidth}{|X|X|X|}
\hline
\rowcolor{popblueL} \textbf{Structure} & \textbf{Exemple chinois + Pinyin} & \textbf{Traduction} \\
\hline
Sujet + \zh{可是} + Phrase & \zh{我想去，可是下雨了。}\newline Wǒ xiǎng qù, kěshì xiàyǔ le. & Je veux y aller, mais il pleut. \\
\hline
Sujet + \zh{但是} + Phrase & \zh{这个方法好，但是太贵。}\newline Zhège fāngfǎ hǎo, dànshì tài guì. & Cette méthode est bonne, mais trop chère. \\
\hline
\zh{可是} en tête & \zh{可是塑料袋污染环境。}\newline Kěshì sùliàodài wūrǎn huánjìng. & Mais les sacs plastiques polluent l'environnement. \\
\hline
\end{tabularx}
\end{center}

\cle{Attention :} En français, « mais » peut parfois se placer en milieu de phrase pour insister (\emph{Il est fatigué, mais il continue}). En chinois, \zh{可是/但是} reste \textbf{fermement ancré au début de la seconde proposition}. On ne dit pas \zh{*塑料袋可是污染环境}.
\end{propriete}

\begin{exemplebox}[Variantes courantes à connaître]
\begin{itemize}
    \item \zh{不过} (bùguò) : « toutefois, simplement » (plus léger, souvent pour nuancer). \zh{塑料袋方便，不过污染大。} (Les sacs sont pratiques, toutefois la pollution est grande.)
    \item \zh{然而} (rán'ér) : « cependant, pourtant » (très formel, écrit).
\end{itemize}
\end{exemplebox}

\courssub{Point de grammaire 2 : La concession polie \zh{你说得也有道理}}

C'est la structure clé pour \textbf{débattre sans heurter} son interlocuteur, valeur culturelle forte en Chine (保全面子, « sauver la face »).

\begin{methode}[Nuancer un désaccord : la stratégie en 3 temps]
Pour passer de l'opposition à la construction d'une solution commune, suivez ce schéma rhétorique standard :
\begin{enumerate}
    \item \textbf{Reconnaître la validité partielle de l'autre} : \zh{你说得也有道理} (Nǐ shuō de yě yǒu dàolǐ) — « Ce que tu dis a aussi du sens / Tu as raison sur ce point. »
    \item \textbf{Marquer la transition} : \zh{那} (nà) ou \zh{那么} (nàme) — « Alors / Dans ce cas. »
    \item \textbf{Proposer l'alternative ou le compromis} : \zh{我们可以用纸袋} (Wǒmen kěyǐ yòng zhǐdài) — « Nous pouvons utiliser des sacs en papier. »
\end{enumerate}

\textbf{Analyse grammaticale de \zh{你说得也有道理} :}
\begin{itemize}
    \item \zh{你} (Sujet) + \zh{说} (Verbe) + \zh{得} (Particule de complément de degré/manière) + \zh{也有道理} (Prédicat : « a aussi de la raison »).
    \item \zh{得} (de) introduit ici l'évaluation de la manière de parler. C'est la structure \textbf{Verbe + 得 + Adjectif/Verbe d'évaluation}.
    \item \zh{也} (yě) = « aussi / également », marque l'ouverture d'esprit.
\end{itemize}

\end{methode}

\begin{exemplebox}[Autres formules de concession pour varier]
\begin{center}
\begin{tabularx}{\linewidth}{|X|X|X|}
\hline
\rowcolor{popblueL} \textbf{Formule} & \textbf{Pinyin} & \textbf{Nuance} \\
\hline
\zh{你的观点很有道理。} & Nǐ de guāndiǎn hěn yǒu dàolǐ. & Ton point de vue est très fondé. (Plus formel) \\
\hline
\zh{你说得对，但是…} & Nǐ shuō de duì, dànshì… & Tu as raison, mais… (Concession + opposition directe) \\
\hline
\zh{有道理，不过…} & Yǒu dàolǐ, bùguò… & C'est sensé, toutefois… (Oral, rapide) \\
\hline
\end{tabularx}
\end{center}
\end{exemplebox}

\coursec{Travail phonétique : les tons du thème}

La prononciation exacte des tons est cruciale pour l'intelligibilité. En débat, une erreur de ton sur un mot-clé (ex: \zh{方便} vs \zh{放扁}) peut changer le sens de votre argument.

\begin{attention}[Prononciation de \zh{方便} (fāngbiàn) : 1er ton + 4e ton — \textbf{Erreur n°1 des francophones}]
\begin{itemize}
    \item \zh{方} (fāng) : \textbf{1er ton (haut et plat — \emph{ā})}. La voix reste haute, stable, ne monte pas, ne descend pas. Imaginez une note de musique tenue longuement.
    \item \zh{便} (biàn) : \textbf{4e ton (descendant fort — \emph{à})}. La voix chute brutalement du haut vers le bas, comme un « Non ! » sec et affirmé.
\end{itemize}
\textbf{Erreur fréquente :} Prononcer « fang-bian » sur un ton plat (ton 1-1), montant (ton 2) ou descendant-descendant (ton 3-4), ce qui rend le mot incompréhensible ou change le sens (ex: \zh{放扁} fàngbiǎn « aplatir »).
\textbf{Exercice oral (Drill) :} Répétez 5 fois en exagérant le contraste : \zh{方──便↘} (fāng──biàn↘). Associez un geste : main horizontale haute pour \zh{方}, main qui tombe verticalement pour \zh{便}.
\end{attention}

\begin{attention}[Autres pièges tonaux du thème]
\begin{center}
\begin{tabularx}{\linewidth}{|X|X|X|X|}
\hline
\rowcolor{popblueL} \textbf{Mot} & \textbf{Pinyin} & \textbf{Tons} & \textbf{Piège fréquent / Astuce} \\
\hline
\zh{环境} & huánjìng & 2 - 4 & \zh{环} monte (2), \zh{境} tombe (4). Ne pas dire *huánjǐng (3). \\
\hline
\zh{污染} & wūrǎn & 1 - 3 & \zh{污} haut plat (1), \zh{染} bas-remontant (3). Attention au 3e ton "demi-ton" en milieu de phrase. \\
\hline
\zh{塑料} & sùliào & 4 - 4 & Deux 4e tons : chute nette sur \zh{塑}, rechute nette sur \zh{料}. Éviter le "chant". \\
\hline
\zh{垃圾} & lājī & 1 - 1 (souvent 1-neutre) & \zh{圾} souvent ton neutre (qīngshēng) à l'oral : \zh{lājī} $\rightarrow$ \zh{lāji}. \\
\hline
\zh{回收} & huíshōu & 2 - 1 & \zh{回} monte (2), \zh{收} haut plat (1). \\
\hline
\zh{一次性} & yīcìxìng & 1 - 4 - 4 & \zh{一} (yī) ton 1 isolé/comptage. \zh{次} (cì) 4e ton. \zh{性} (xìng) 4e ton. \\
\hline
\end{tabularx}
\end{center}
\end{attention}

\courssub{Lexique argumentatif : Structurer le POUR et le CONTRE}

Pour débattre, il faut pouvoir lister des arguments des deux côtés. Le tableau ci-dessous synthétise le vocabulaire utile pour le sujet des sacs plastiques, réutilisable pour d'autres thèmes (eau, arbres, déchets).

\begin{popfigure}
\begin{tikzpicture}[
    node distance=0.8cm and 1.5cm,
    argbox/.style={rectangle, rounded corners=4pt, draw=popline, thick, fill=popcream, text width=5.5cm, align=center, minimum height=1.8cm, font=\small},
    titlebox/.style={rectangle, rounded corners=4pt, draw=popdark, thick, fill=popblue, text=white, text width=5.5cm, align=center, font=\bfseries\normalsize},
    arrowstyle/.style={->, >=Stealth, thick, popblue}
]
% Column titles
\node[titlebox] (pourtitle) {POUR interdire (论点)};
\node[titlebox] (contretitle) [right=of pourtitle] {CONTRE interdire (反方论点)};

% POUR arguments
\node[argbox] (p1) [below=of pourtitle] {\zh{保护环境}\newline bǎohù huánjìng\newline Protéger l'environnement};
\node[argbox] (p2) [below=of p1] {\zh{减少白色污染}\newline jiǎnshǎo báisè wūrǎn\newline Réduire la « pollution blanche »};
\node[argbox] (p3) [below=of p2] {\zh{保护动物健康}\newline bǎohù dòngwù jiànkāng\newline Protéger la santé des animaux};
\node[argbox] (p4) [below=of p3] {\zh{推广纸袋/布袋}\newline tuīguǎng zhǐdài/bùdài\newline Promouvoir sacs papier/tissu};

% CONTRE arguments
\node[argbox] (c1) [below=of contretitle] {\zh{很方便}\newline hěn fāngbiàn\newline Très pratique};
\node[argbox] (c2) [below=of c1] {\zh{很便宜}\newline hěn piányi\newline Très bon marché};
\node[argbox] (c3) [below=of c2] {\zh{防水}\newline fángshuǐ\newline Imperméable};
\node[argbox] (c4) [below=of c3] {\zh{习惯了}\newline xíguàn le\newline On a l'habitude};

% Arrows for visual flow
\draw[arrowstyle] (pourtitle.south) -- (p1.north);
\draw[arrowstyle] (p1.south) -- (p2.north);
\draw[arrowstyle] (p2.south) -- (p3.north);
\draw[arrowstyle] (p3.south) -- (p4.north);

\draw[arrowstyle] (contretitle.south) -- (c1.north);
\draw[arrowstyle] (c1.south) -- (c2.north);
\draw[arrowstyle] (c2.south) -- (c3.north);
\draw[arrowstyle] (c3.south) -- (c4.north);

% Central "VS" separator
\node[draw=poporange, thick, circle, fill=poporange!20, font=\bfseries, minimum size=1.2cm] at ($(pourtitle)!0.5!(contretitle) + (0,-4.5)$) {VS};
\end{tikzpicture}
\legende{Cartographie des arguments : POUR vs CONTRE l'interdiction des sacs plastiques.}
\end{popfigure}

\begin{center}
\begin{tabularx}{\linewidth}{|X|X|X|X|}
\hline
\rowcolor{popblueL} \textbf{Caractères} & \textbf{Pinyin} & \textbf{Nature} & \textbf{Traduction} \\
\hline
白色污染 & báisè wūrǎn & nom & pollution blanche (sacs plastiques) \\
\hline
推广 & tuīguǎng & verbe & promouvoir, vulgariser \\
\hline
布袋 & bùdài & nom & sac en tissu \\
\hline
便宜 & piányi & adj. & bon marché, pas cher \\
\hline
防水 & fángshuǐ & verbe/adj. & imperméable, étanche \\
\hline
习惯 & xíguàn & verbe & avoir l'habitude de \\
\hline
论点 & lùndiǎn & nom & argument (dans un débat) \\
\hline
反方 & fǎngfāng & nom & partie adverse (dans un débat) \\
\hline
\end{tabularx}
\end{center}

\courssub{Exercice résolu : Transposition du dialogue (Changement de sujet)}

\begin{exoresolu}[Adapter le mini-débat au thème : « Gaspillage de l'eau »]
\textbf{Consigne :} Reprenez la structure exacte du dialogue modèle (A: Proposition $\rightarrow$ B: Opposition + argument praticité $\rightarrow$ A: Objection environnementale $\rightarrow$ B: Concession + Solution alternative $\rightarrow$ A: Validation) pour le sujet : \zh{节约用水} (économiser l'eau). Utilisez \zh{可是} et \zh{你说得也有道理}.

\textbf{Solution détaillée :}

\textbf{A :} 我认为我们应该节约用水。\\
Wǒ rènwéi wǒmen yīnggāi jiéyuē yòngshuǐ.\\
« Je pense que nous devrions économiser l'eau. »

\textbf{B :} 我不同意。用水多很方便。\\
Wǒ bù tóngyì. Yòngshuǐ duō hěn fāngbiàn.\\
« Je ne suis pas d'accord. Utiliser beaucoup d'eau est pratique. »

\textbf{A :} 可是浪费水污染环境。\\
Kěshì làngfèi shuǐ wūrǎn huánjìng.\\
« Mais gaspiller l'eau pollue l'environnement. » \textit{(Alternative : \zh{可是水资源很少} — Mais les ressources en eau sont rares.)}

\textbf{B :} 你说得也有道理。那我们可以收集雨水使用。\\
Nǐ shuō de yě yǒu dàolǐ. Nà wǒmen kěyǐ shōují yǔshuǐ shǐyòng.\\
« Ce que tu dis a du sens. Alors nous pouvons collecter l'eau de pluie pour l'utiliser. »

\textbf{A :} 好主意！\\
Hǎo zhǔyi!\\
« Bonne idée ! »

\textbf{Analyse de la transposition :}
\begin{itemize}
    \item Remplacement du vocabulaire thème : \zh{塑料袋} $\rightarrow$ \zh{水/用水} ; \zh{污染环境} conservé (ou \zh{水资源少}) ; \zh{纸袋} $\rightarrow$ \zh{收集雨水}.
    \item Conservation des \textbf{structures grammaticales cibles} : \zh{我认为我们应该…}, \zh{我不同意}, \zh{可是…}, \zh{你说得也有道理}, \zh{那我们可以…}.
    \item Respect de l'enchaînement logique : Proposition $\rightarrow$ Objection praticienne $\rightarrow$ Argument éthique/écologique $\rightarrow$ Concession $\rightarrow$ Solution technique $\rightarrow$ Accord.
\end{itemize}
\end{exoresolu}

\coursec{Jeux de rôle et débat guidé}

\courssub{Entraînement à l'expression orale : Mini-débat en binôme}

\begin{experience}[Mini-débat en binôme : Trois sujets au choix]
\textbf{Durée :} 15 minutes (préparation 5 min, passage 5 min/binôme, feedback 5 min).

\textbf{Déroulement :}
\begin{enumerate}
    \item \textbf{Formez des binômes.} L'un joue le rôle \textbf{A} (propose une mesure écologique), l'autre le rôle \textbf{B} (objecte par praticité/coût, puis accepte la concession).
    \item \textbf{Choisissez UN sujet parmi les trois :}
    \begin{itemize}
        \item \textbf{Sujet 1 :} \zh{保护树木} (Protéger les arbres) — \textit{Arguments : couper bois vs planter / papier recyclé. Voc. clé : \zh{砍树} (kǎn shù), \zh{种树} (zhòng shù), \zh{循环纸} (xúnhuán zhǐ).}
        \item \textbf{Sujet 2 :} \zh{垃圾分类} (Tri des déchets) — \textit{Arguments : trop compliqué vs nécessaire pour recyclage. Voc. clé : \zh{麻烦} (máfan), \zh{必要} (bìyào), \zh{可回收} (kě huíshōu), \zh{厨余垃圾} (chúyú lājī).}
        \item \textbf{Sujet 3 :} \zh{少开车，多坐公交} (Moins de voiture, plus de bus) — \textit{Arguments : voiture pratique vs pollution/embouteillages. Voc. clé : \zh{堵车} (dǔchē), \zh{尾气} (wěiqì), \zh{地铁} (dìtiě), \zh{共享单车} (gòngxiǎng dānchē).}
    \end{itemize}
    \item \textbf{Préparez 5 arguments (POUR/CONTRE)} en vous aidant du tableau TikZ et du vocabulaire ci-dessus. Écrivez vos répliques en caractères + pinyin.
    \item \textbf{Jouez le dialogue devant la classe.} Critères d'évaluation :
    \begin{itemize}
        \item Respect de la structure en 5 tours de parole (A-B-A-B-A).
        \item Utilisation correcte de \zh{可是/但是} et \zh{你说得也有道理}.
        \item Prononciation distincte des tons (notamment \zh{方便}, \zh{污染}, \zh{环境}, \zh{垃圾}, \zh{堵车}).
        \item Fluidité et intonation naturelle (pas de lecture robotique ; regard vers l'interlocuteur).
    \end{itemize}
\end{enumerate}

\textbf{Aide-mémoire pour le rôle B (l'opposant constructif) :}
\begin{textebox}[Fiche rôle B - Mémo]
1. \zh{我不同意。} + Argument praticité/coût (\zh{很方便 / 很便宜 / 习惯了 / 麻烦})\\
2. Écouter l'argument A (\zh{可是…})\\
3. \zh{你说得也有道理。} (OBLIGATOIRE pour la politesse du débat)\\
4. \zh{那我们可以…} + Solution alternative (\zh{用纸袋 / 用布袋 / 收集雨水 / 坐公交 / 用循环纸 / 骑共享单车})\\
5. Sourire et valider : \zh{好主意！} / \zh{就这么定！}
\end{textebox}
\end{experience}

\courssub{Le saviez-vous ? : Législation comparée Cameroun – Chine}

\begin{savaistu}[La guerre aux sacs plastiques : 2008 (Chine) vs 2014 (Cameroun)]
Saviez-vous que la Chine a été pionnière en Asie sur l'interdiction des sacs plastiques gratuits ?
\begin{itemize}
    \item \textbf{Chine (1er juin 2008) :} Entrée en vigueur du \zh{限塑令} (Xiàn sù lìng — « Ordre de restriction du plastique »). Interdiction de la production, vente et utilisation de sacs plastiques d'épaisseur < 0,025 mm. Les supermarchés ne doivent plus les donner gratuitement. Résultat : réduction de \textbf{60 à 80\%} de la consommation de sacs dans les grandes surfaces les premières années.
    \item \textbf{Cameroun (avril 2014) :} Promulgation de la loi n°2014/006 interdisant la fabrication, l'importation et la commercialisation des emballages plastiques non biodégradables (épaisseur < 60 microns). Application progressive, avec des opérations \zh{清理环境} (qīnglǐ huánjìng — nettoyage de l'environnement) dans les marchés (Marché Central Yaoundé, Marché Mokolo, Marché Sandaga Douala).
\end{itemize}
\textbf{Point commun :} Les deux pays visent la réduction de la \zh{白色污染} (pollution blanche). \textbf{Différence :} La Chine a misé sur la taxation (payant) ; le Cameroun sur l'interdiction pure (sauf biodégradables). Dans les deux cas, le \zh{布袋} (sac en tissu) et le \zh{纸袋} (sac en papier) redeviennent des accessoires du quotidien.

\textbf{Vocabulaire bonus :}
\begin{center}
\begin{tabularx}{\linewidth}{|X|X|X|X|}
\hline
\rowcolor{popblueL} \textbf{Caractères} & \textbf{Pinyin} & \textbf{Nature} & \textbf{Traduction} \\
\hline
限塑令 & xiàn sù lìng & nom & Ordre de restriction du plastique (Chine) \\
\hline
可降解 & kě jiàngjiě & adj. & biodégradable \\
\hline
超市 & chāoshì & nom & supermarché \\
\hline
免费 & miǎnfèi & adj. & gratuit \\
\hline
收费 & shōufèi & verbe & faire payer, taxer \\
\hline
厚度 & hòudù & nom & épaisseur \\
\hline
市场 & shìchǎng & nom & marché \\
\hline
\end{tabularx}
\end{center}
\end{savaistu}

\courssub{Exercice d'application : À vous de jouer !}

\begin{exercice}[Rédaction guidée : Mon mini-débat écrit]
Rédigez un mini-débat complet (5 répliques : A-B-A-B-A) sur le sujet : \zh{保护树木：少用筷子，自带筷子} (Protéger les arbres : moins de baguettes jetables, apporter les siennes).
\textbf{Contraintes :}
\begin{enumerate}
    \item A propose : \zh{我认为我们应该自带筷子。}
    \item B s'oppose : \zh{一次性筷子很方便。} (Utilisez \zh{一次性} ).
    \item A objecte avec \zh{可是} : argument déforestation (\zh{砍树} ).
    \item B concède avec \zh{你说得也有道理} et propose solution : \zh{那我们可以用金属筷子 / 竹筷子循环使用。}
    \item A conclut : \zh{好主意！}
\end{enumerate}
Écrivez en \textbf{caractères}, \textbf{pinyin} (avec tons) et \textbf{traduction française}.
\end{exercice}

\begin{corrige}[Exercice : Mon mini-débat écrit]
\textbf{A :} 我认为我们应该自带筷子。\\
Wǒ rènwéi wǒmen yīnggāi zìdài kuàizi.\\
« Je pense que nous devrions apporter nos propres baguettes. »

\textbf{B :} 我不同意。一次性筷子很方便。\\
Wǒ bù tóngyì. Yīcìxìng kuàizi hěn fāngbiàn.\\
« Je ne suis pas d'accord. Les baguettes jetables sont pratiques. »

\textbf{A :} 可是制作一次性筷子导致砍树太多，污染环境。\\
Kěshì zhìzuò yīcìxìng kuàizi dǎozhì kǎn shù tài duō, wūrǎn huánjìng.\\
« Mais la fabrication de baguettes jetables entraîne l'abattage de trop d'arbres et pollue l'environnement. »

\textbf{B :} 你说得也有道理。那我们可以用金属筷子或竹筷子循环使用。\\
Nǐ shuō de yě yǒu dàolǐ. Nà wǒmen kěyǐ yòng jīnshǔ kuàizi huò zhú kuàizi xúnhuán shǐyòng.\\
« Ce que tu dis a du sens. Alors nous pouvons utiliser des baguettes en métal ou en bambou de façon réutilisable. »

\textbf{A :} 好主意！\\
Hǎo zhǔyi!\\
« Bonne idée ! »

\textbf{Notes de correction :}
\begin{itemize}
    \item \zh{一次性} (yīcìxìng) = « à usage unique, jetable ». Structure : \zh{一} + \zh{次} + \zh{性}.
    \item \zh{制作} (zhìzuò) = « fabriquer, faire ». \zh{导致} (dǎozhì) = « entraîner, causer (souvent négatif) ». Structure : \zh{制作…导致…} (La fabrication de… cause…). Plus précis que \zh{一次性筷子砍树}.
    \item \zh{砍树} (kǎn shù) = « couper des arbres » (verbe + objet). \zh{太多} (tài duō) = « trop ».
    \item \zh{金属} (jīnshǔ) = « métal » ; \zh{竹} (zhú) = « bambou » ; \zh{循环使用} (xúnhuán shǐyòng) = « utiliser en boucle / réutiliser ».
    \item \zh{或} (huò) = « ou » (écrit/formel), interchangeable avec \zh{或者} (huòzhě) à l'oral.
    \item Vérifiez les tons : \zh{筷子} (kuàizi, 4-neutre), \zh{竹} (zhú, 2), \zh{金属} (jīnshǔ, 1-3), \zh{导致} (dǎozhì, 3-4).
\end{itemize}
\end{corrige}

\coursec{Travail phonétique : les tons du thème}

Après avoir mis en place les arguments et les structures du mini-débat (section précédente), il est indispensable de sécuriser la prononciation du vocabulaire spécifique à l'environnement. En chinois, un ton modifié change radicalement le sens d'un mot : dire \zh{树 shù} (arbre) avec une montée de voix (2e ton) le transforme en \zh{数 shǔ} (nombre) ou \zh{书 shū} (livre) selon le contexte. Cette section ancre les contours toniques des mots-clés du thème et entraîne l'oreille à distinguer les paires minimales pièges.

\courssub{Observation et écoute active : les mots-clés du thème}

Écoutons d'abord les six mots « pivots » de cette leçon, classés par type de contour tonique. Répétez chaque mot en exagérant le mouvement de la voix.

\begin{center}
\begin{tabularx}{\linewidth}{|X|X|X|X|}
\hline
\rowcolor{popblueL} \textbf{Caractères} & \textbf{Pinyin} & \textbf{Nature} & \textbf{Traduction} \\
\hline
环境 & huánjìng & nom & environnement \\
\hline
保护 & bǎohù & verbe & protéger \\
\hline
污染 & wūrǎn & verbe/nom & polluer / pollution \\
\hline
垃圾 & lājī & nom & déchets, ordures \\
\hline
应该 & yīnggāi & verbe modal & devoir, devrait \\
\hline
树 & shù & nom & arbre \\
\hline
\end{tabularx}
\end{center}

\courssub{Analyse détaillée des contours toniques}

Observons chaque mot sous la loupe phonétique. Le chinois standard (putonghua) possède quatre tons de base (plus le ton neutre). Chaque ton est un \cle{contour mélodique} précis, non une simple hauteur.

\begin{propriete}[Le 1er ton (阴平 yīnpíng) : haut et plat]
Le 1er ton se maintient à une hauteur aiguë constante, comme une note de musique tenue. \textbf{Ne pas laisser descendre la voix.}
\begin{itemize}
    \item \zh{污 wū} (dans \zh{污染 wūrǎn}) : partie haute, plate.
    \item \zh{垃 lā} et \zh{圾 jī} (\zh{垃圾 lājī}) : deux 1ers tons consécutifs → la voix reste haute sur les deux syllabes. C'est fatigant pour une oreille francophone habituée à la variation.
    \item \zh{应 yīng} et \zh{该 gāi} (\zh{应该 yīnggāi}) : même configuration, deux syllabes hautes et plates.
\end{itemize}
\emph{Astuce :} Imaginez un fil tendu horizontalement au-dessus de votre tête ; votre voix glisse dessus.
\end{propriete}

\begin{propriete}[Le 2e ton (阳平 yángpíng) : montant]
Le 2e ton part du milieu vers le haut, comme une question surprise en français (« Hein ? », « Vraiment ? »).
\begin{itemize}
    \item \zh{环 huán} (dans \zh{环境 huánjìng}) : départ moyen → arrivée haute. Le 4e ton \zh{境 jìng} qui suit chute brutalement.
\end{itemize}
\emph{Enchaînement 2e + 4e :} La montée de \zh{huán} prépare la chute de \zh{jìng}. Visualisez une montagne russe : montée raide, descente verticale.
\end{propriete}

\begin{propriete}[Le 3e ton (上声 shǎngshēng) : bas, descendant-remontant (citation) / demi-3e ton (contexte)]
En isolation (citation), le 3e ton descend bas puis remonte (forme en V). \textbf{Mais en phrase, devant un 1er, 2e ou 4e ton, il ne fait que descendre (demi-3e ton).}
\begin{itemize}
    \item \zh{保 bǎo} (dans \zh{保护 bǎohù}) : devant \zh{护 hù} (4e ton), \zh{保} se prononce bas, court, sans remontée. On note souvent \zh{bǎo\kern-1.5pt\raisebox{1pt}{\tiny{}}hù} pour marquer l'absence de remontée.
    \item \zh{污 wū} est 1er ton, mais \zh{染 rǎn} est 3e ton. En fin de phrase ou isolé, \zh{rǎn} fait le V complet. Devant un autre mot (ex: \zh{污染严重 wūrǎn yánzhòng}), \zh{rǎn} devient demi-3e ton.
\end{itemize}
\end{propriete}

\begin{propriete}[Le 4e ton (去声 qùshēng) : descendant fort]
Le 4e ton part du haut et tombe violemment vers le bas, comme un ordre sec (« Arrête ! », « Non ! »).
\begin{itemize}
    \item \zh{境 jìng}, \zh{护 hù}, \zh{树 shù} : chute nette, courte, énergique.
    \item \zh{该 gāi} est 1er ton (attention, pas 4e !), mais \zh{树 shù} est bien 4e.
\end{itemize}
\end{propriete}

\courssub{Visualisation graphique : les courbes toniques}

Le schéma ci-dessous place nos mots-clés sur les quatre contours canoniques. Observez la position relative de \zh{环} (2e), \zh{保} (3e/demi), \zh{污} (1er), \zh{树} (4e).

\begin{popfigure}
\begin{tikzpicture}[
    scale=1.2,
    tone/.style={very thick, popblue},
    toneA/.style={very thick, poporange},
    toneB/.style={very thick, popgreen},
    toneC/.style={very thick, poppurple},
    axis/.style={thin, popmuted, ->},
    label/.style={font=\small, popdark},
    word/.style={font=\normalsize\bfseries, popink, anchor=west}
]
% Axes
\draw[axis] (0,0) -- (5,0) node[right, label] {Durée};
\draw[axis] (0,0) -- (0,4.5) node[above, label] {Hauteur (F0)};

% Grille de repère (5 niveaux Chao)
\foreach \y in {1,2,3,4,5} {
    \draw[popline, dashed] (0,\y) -- (5,\y);
}
\foreach \y/\ylab in {1/1, 2/2, 3/3, 4/4, 5/5} {
    \node[label, anchor=east] at (-0.2,\y) {\ylab};
}

% 1er Ton : Haut plat (niveau 5)
\draw[tone] (0.5,5) -- (2,5) node[word, pos=0.5, above=2mm] {\zh{污 wū} / \zh{垃 lā} / \zh{应 yīng}};
\node[label] at (1.25, 4.2) {1er ton (55)};

% 2e Ton : Montant (35)
\draw[toneA] (2.5,3) .. controls (3,4) and (3.5,4.5) .. (4,5) node[word, pos=1, above=2mm] {\zh{环 huán}};
\node[label] at (3.25, 3.5) {2e ton (35)};

% 3e Ton : Bas descendant (21) - Demi-ton en contexte
\draw[toneB] (0.5,2) -- (1.5,1) node[word, pos=1, below=2mm] {\zh{保 bǎo (demi)}};
\node[label] at (1, 2.5) {3e ton (21) / demi};

% 4e Ton : Descendant fort (51)
\draw[toneC] (3,5) -- (4,1) node[word, pos=1, below=2mm] {\zh{树 shù} / \zh{境 jìng} / \zh{护 hù}};
\node[label] at (3.5, 4.2) {4e ton (51)};

% Légende mots complets
\node[label, align=left, anchor=north west, popdark] at (0, -0.8) {
    \zh{环境 huánjìng} : 2e + 4e \quad
    \zh{保护 bǎohù} : 3e(demi) + 4e \quad
    \zh{垃圾 lājī} : 1er + 1er \quad
    \zh{应该 yīnggāi} : 1er + 1er
};
\end{tikzpicture}
\legende{Contours toniques des mots-clés de l'environnement (échelle de Chao 1-5). Le 3e ton \zh{保} est représenté en demi-ton (21) car suivi d'un 4e ton.}
\end{popfigure}

\courssub{Technique « main-ton » : incarner la mélodie}

\begin{methode}[Technique « main-ton » : dessine le ton avec la main]
Cette technique kinesthésique lie le geste à la prosodie. Elle est utilisée à l'oral pour ancrer le muscle laryngé.
\begin{enumerate}
    \item \textbf{1er ton (haut plat) :} Main horizontale, paume vers le bas, niveau menton/haut du front. Glisse la main vers l'avant sans bouger la hauteur. \emph{Ex: \zh{垃圾 lājī}.}
    \item \textbf{2e ton (montant) :} Main partie hauteur taille, paume vers l'intérieur. Monte la main vers le visage en prononçant. \emph{Ex: \zh{环 huán}.}
    \item \textbf{3e ton (demi-ton bas) :} Main partie hauteur taille, descends vers les genoux, \textbf{ne remonte pas}. Marque un temps d'arrêt bas. \emph{Ex: \zh{保 bǎo} dans \zh{保护}.}
    \item \textbf{4e ton (descendant) :} Main partie haut (front), frappe vers le bas (hanche) net et sec. \emph{Ex: \zh{树 shù}, \zh{护 hù}.}
    \item \textbf{Enchaînement :} Enchaîne les gestes pour \zh{环境} (montée \zh{环} → chute \zh{境}), \zh{保护} (bas \zh{保} → chute \zh{护}).
\end{enumerate}
\emph{Entraîne-toi devant un miroir : le geste doit être fluide, synchrone avec la voix.}
\end{methode}

\courssub{Exercices de discrimination auditive : paires minimales et mots proches}

L'oreille francophone confond souvent des tons voisins (1er/2e, 3e/4e) ou projette l'intonation française (montée finale) sur les tons chinois. Travaillons les paires critiques du thème.

\begin{exemplebox}[Discrimination : 树 shù (4e) vs 水 shuǐ (3e) vs 说 shuō (1er)]
Trois mots, même consonne \emph{sh}, même voyelle \emph{u} (ou \emph{uo}), \textbf{tons différents} = sens radicalement différents.
\begin{center}
\begin{tabularx}{\linewidth}{|X|X|X|X|}
\hline
\rowcolor{popblueL} \textbf{Caractère} & \textbf{Pinyin} & \textbf{Ton} & \textbf{Sens} \\
\hline
树 & shù & 4e (chute) & arbre \\
\hline
水 & shuǐ & 3e (bas / V) & eau \\
\hline
说 & shuō & 1er (haut plat) & parler, dire \\
\hline
\end{tabularx}
\end{center}
\emph{Exercice oral :} Le professeur prononce l'un des trois ; l'élève lève la carte « 1 », « 3 » ou « 4 » et donne le sens.
\begin{itemize}
    \item \zh{树很高。} (Shù hěn gāo.) — L'arbre est haut.
    \item \zh{水很干净。} (Shuǐ hěn gānjìng.) — L'eau est propre.
    \item \zh{他说中文。} (Tā shuō Zhōngwén.) — Il parle chinois.
\end{itemize}
\end{exemplebox}

\begin{exemplebox}[Paire minimale : 种 zhòng (4e) vs 中 zhōng (1er)]
Très fréquent dans le vocabulaire écologique : \zh{种树 zhòng shù} (planter des arbres) vs \zh{中国 Zhōngguó} (Chine) / \zh{中间 zhōngjiān} (milieu).
\begin{center}
\begin{tabularx}{\linewidth}{|X|X|X|X|}
\hline
\rowcolor{popblueL} \textbf{Caractère} & \textbf{Pinyin} & \textbf{Ton} & \textbf{Sens / Exemple} \\
\hline
种 & zhòng & 4e (chute) & planter (verbe) — \zh{种树 zhòng shù} \\
\hline
中 & zhōng & 1er (haut plat) & milieu, centre — \zh{中国 Zhōngguó}, \zh{中间 zhōngjiān} \\
\hline
\end{tabularx}
\end{center}
\emph{Attention :} \zh{种} se prononce aussi \zh{zhǒng} (3e ton) = « espèce, genre, graine » (nom). Ici, c'est le verbe « planter » → 4e ton.
\end{exemplebox}

\begin{attention}[Piège francophone : l'accent final montant (intonation interrogative)]
En français, une phrase affirmative se termine souvent par une légère montée (continuation) ou une chute faible. En chinois, \textbf{le ton lexical prime sur l'intonation syntaxique}.
\begin{itemize}
    \item \textbf{Erreur :} Dire \zh{保护环境 bǎohù huánjíng\kern-1pt\raisebox{1pt}{\textasciicircum}} avec une montée finale sur \zh{境} (comme une question « Tu protèges l'environnement ? »).
    \item \textbf{Correct :} \zh{保护环境 bǎohù huánjìng.} \zh{境} doit chuter (4e ton), même en fin de phrase affirmative. La chute du 4e ton \emph{est} le point final.
    \item \textbf{Erreur sur 树 :} \zh{种树 zhòng shú\kern-1pt\raisebox{1pt}{\textasciicircum}} (montée finale) → compris comme \zh{种熟} (mûrir ?) ou incompréhensible. \zh{树 shù} \textbf{doit descendre} violemment.
\end{itemize}
\cle{Règle d'or} : En chinois, on « atterrit » sur le ton du dernier mot. Pas de « queue » mélodique française.
\end{attention}

\begin{savaistu}[Un ton change le sens : planter vs milieu]
\zh{种树 zhòng shù} (planter des arbres) est une action écologique majeure. Si tu prononces \zh{zhōng shù} (1er ton + 4e ton), tu dis « arbre du milieu / arbre central » (\zh{中树}), qui n'est pas un mot courant et ne veut pas dire « planter ». Pire, \zh{zhǒng shù} (3e + 4e) signifie « espèce d'arbre / graine d'arbre ». Le verbe « planter » est \textbf{exclusivement} \zh{zhòng} (4e ton). C'est la différence entre \emph{agir pour la planète} et \emph{désigner un objet}.
\end{savaistu}

\begin{exoresolu}[Écoute et identification des tons]
\textbf{Consigne :} Écoute l'enseignant (ou l'audio) lire la liste suivante. Pour chaque item, note le numéro du ton (1, 2, 3, 4) de la syllabe \textbf{surlignée} et le sens du mot.

\begin{enumerate}
    \item \zh{\textbf{环}境} \hfill (Ton : \_\_\_ ; Sens : \_\_\_\_\_\_\_\_\_\_\_\_)
    \item \zh{保\textbf{护}} \hfill (Ton : \_\_\_ ; Sens : \_\_\_\_\_\_\_\_\_\_\_\_)
    \item \zh{\textbf{污}染} \hfill (Ton : \_\_\_ ; Sens : \_\_\_\_\_\_\_\_\_\_\_\_)
    \item \zh{垃\textbf{圾}} \hfill (Ton : \_\_\_ ; Sens : \_\_\_\_\_\_\_\_\_\_\_\_)
    \item \zh{应该} (premier caractère \zh{\textbf{应}}) \hfill (Ton : \_\_\_ ; Sens : \_\_\_\_\_\_\_\_\_\_\_\_)
    \item \zh{\textbf{树}} \hfill (Ton : \_\_\_ ; Sens : \_\_\_\_\_\_\_\_\_\_\_\_)
    \item \zh{\textbf{种}}树 \hfill (Ton : \_\_\_ ; Sens : \_\_\_\_\_\_\_\_\_\_\_\_)
    \item \zh{中\textbf{国}} \hfill (Ton : \_\_\_ ; Sens : \_\_\_\_\_\_\_\_\_\_\_\_)
\end{enumerate}

\tcblower
\textbf{Solution détaillée :}
\begin{enumerate}
    \item \zh{\textbf{环}境} : \textbf{2e ton} (montant). Sens : environnement.
    \item \zh{保\textbf{护}} : \textbf{4e ton} (chute). Sens : protéger (verbe). Note : \zh{保} est 3e ton (demi-ton bas ici).
    \item \zh{\textbf{污}染} : \textbf{1er ton} (haut plat). Sens : polluer / pollution.
    \item \zh{垃\textbf{圾}} : \textbf{1er ton} (haut plat). Sens : déchets (le mot complet \zh{垃圾} est 1er+1er).
    \item \zh{\textbf{应}}该 : \textbf{1er ton} (haut plat). Sens : devoir / devrait (verbe modal).
    \item \zh{\textbf{树}} : \textbf{4e ton} (chute forte). Sens : arbre.
    \item \zh{\textbf{种}}树 : \textbf{4e ton} (chute). Sens : planter (verbe). \emph{Piège évité : pas 1er ton (zhōng), pas 3e ton (zhǒng).}
    \item \zh{中\textbf{国}} : \textbf{2e ton} (montant) sur \zh{国}. \zh{中} est 1er ton. Sens : Chine.
\end{enumerate}
\emph{Auto-évaluation :} Si tu as 8/8, tes repères toniques sont solides. Si erreur sur 3e/4e ou 1er/2e, refais l'exercice « main-ton » sur les mots concernés.
\end{exoresolu}

\courssub{Entraînement en phrase complète : le ton dans le flux}

Les tons ne vivent pas isolés ; ils s'enchaînent (liaison, sable tonal). Répétons une phrase modèle intégrant plusieurs mots du thème.

\begin{exemplebox}[Phrase modèle : « L'eau est très propre, il faut la protéger. »]
\zh{水很干净，应该保护。}
\emph{Shuǐ hěn gānjìng, yīnggāi bǎohù.}

\textbf{Analyse ton par ton :}
\begin{center}
\begin{tabularx}{\linewidth}{|X|X|X|X|}
\hline
\rowcolor{popblueL} \textbf{Syllabe} & \textbf{Pinyin} & \textbf{Ton cité} & \textbf{Ton réalisé (contexte)} \\
\hline
水 & shuǐ & 3e (V complet) & 3e (V complet, fin de proposition) \\
\hline
很 & hěn & 3e & \textbf{Demi-3e ton (bas, 21)} (règle : 3e + 1er/2e/4e → demi-3e) \\
\hline
干 & gān & 1er & 1er (haut plat) \\
\hline
净 & jìng & 4e & 4e (chute) \\
\hline
应 & yīng & 1er & 1er (haut plat) \\
\hline
该 & gāi & 1er & 1er (haut plat) \\
\hline
保 & bǎo & 3e & \textbf{Demi-3e ton} (bas, pas de remontée car suivi de 4e) \\
\hline
护 & hù & 4e & 4e (chute finale) \\
\hline
\end{tabularx}
\end{center}

\textbf{Points de vigilance :}
\begin{itemize}
    \item \zh{很 hěn} : Devant \zh{干 gān} (1er ton), le 3e ton de \zh{很} devient \textbf{demi-3e ton (bas, sans remontée)}. \textbf{Attention :} La règle « 3e + 3e → 2e + 3e » ne s'applique \textbf{que} si la syllabe suivante est \emph{aussi} un 3e ton. Ici \zh{干} est 1er ton → \zh{很} reste bas.
    \item \zh{保 bǎo} : Devant \zh{护 hù} (4e ton), \zh{保} reste \textbf{bas} (demi-3e ton). Ne remontez pas !
    \item \zh{水 shuǐ} : En fin de premier segment (virgule), il fait le V complet (descend bas, remonte un peu). Marque bien la descente.
\end{itemize}
\emph{Exercice « main-ton » sur la phrase complète :} Trace le parcours de la main pour les 11 syllabes. La main danse : V (\zh{shuǐ}), \textbf{bas} (\zh{hěn}), plat (\zh{gān}), chute (\zh{jìng}), pause, plat (\zh{yīng}), plat (\zh{gāi}), bas (\zh{bǎo}), chute (\zh{hù}).
\end{exemplebox}

\begin{exercice}[Entraînement personnel : lecture à voix haute et enregistrement]
\textbf{Consigne :} Enregistre-toi (téléphone, ordinateur) en lisant les items ci-dessous. Réécoute-toi en comparant avec le modèle de l'enseignant / l'audio de référence. Coche les cases quand c'est stable.

\begin{enumerate}
    \item \textbf{Mots isolés (tons de citation) :} \zh{环境} — \zh{保护} — \zh{污染} — \zh{垃圾} — \zh{应该} — \zh{树} — \zh{种(zhòng)} — \zh{水} — \zh{说} — \zh{中}.
    \item \textbf{Enchaînements bisyllabiques / trisyllabiques :}
    \begin{itemize}
        \item \zh{保护环境} (3e+demi + 4e + 2e + 4e)
        \item \zh{污染水} (\zh{wū} 1er + \zh{rǎn} 3e→\textbf{2e} car devant \zh{shuǐ} 3e + \zh{shuǐ} 3e V final)
        \item \zh{种树} (4e + 4e)
        \item \zh{应该做} (1er + 1er + 4e)
    \end{itemize}
    \item \textbf{Phrases complètes :}
    \begin{enumerate}
        \item \zh{我们应该保护环境。} (Wǒmen yīnggāi bǎohù huánjìng.)
        \item \zh{垃圾要分类。} (Lājī yào fēnleì.) — \emph{注意 : 分类 fēnleì (1er + 4e).}
        \item \zh{水很干净，我们种树。} (Shuǐ hěn gānjìng, wǒmen zhòng shù.)
        \item \zh{中国保护大熊猫。} (Zhōngguó bǎohù dàxióngmāo.) — \emph{注意 : 大熊猫 dàxióngmāo (4e + 2e + 1er).}
    \end{enumerate}
\end{enumerate}
\textbf{Critères de réussite :} Aucune montée finale parasite sur les 4e tons ; \zh{很} devient demi-3e ton (bas) devant 1er ton ; \zh{保} reste bas devant \zh{护} ; distinction nette \zh{zhòng} (4e) / \zh{zhōng} (1er) ; \zh{染} devient 2e ton devant \zh{水} (3e).
\end{exercice}

\begin{corrige}[Exercice 1]
\textbf{Auto-correction guidée :}
\begin{itemize}
    \item \textbf{Mots isolés :} Vérifie au spectrogramme (appli type Praat, Pleco, ou simple réécoute) : \zh{环} monte (3→5), \zh{保} descend bas (2→1), \zh{污} plat haut (5), \zh{垃} \zh{圾} \zh{应} \zh{该} plats haut (5), \zh{树} \zh{种(zhòng)} \zh{护} \zh{境} chutent (5→1), \zh{水} V bas (2→1→3), \zh{说} plat haut (5), \zh{中} plat haut (5).
    \item \textbf{Enchaînements :}
        \begin{itemize}
            \item \zh{保护环境} : \zh{bǎo} (bas) \zh{hù} (chute) \zh{huán} (montée) \zh{jìng} (chute). Pas de remontée sur \zh{bǎo}.
            \item \zh{污染水} : \zh{wū} (plat) \zh{rǎn} (3e→\textbf{2e montant} car devant \zh{shuǐ} 3e) \zh{shuǐ} (3e V final). \textbf{Règle : 3e + 3e → 2e + 3e.}
            \item \zh{种树} : Deux 4e tons consécutifs → deux chutes sèches. \zh{zhòng shù}. Pas de liaison mélodique.
            \item \zh{应该做} : Trois tons plats/hauts puis chute. \zh{yīng gāi zuò}.
        \end{itemize}
    \item \textbf{Phrases :}
        \begin{itemize}
            \item \zh{我们应该保护环境} : \zh{wǒmen} (3e+ton neutre) \zh{yīnggāi} (1+1) \zh{bǎo} (bas) \zh{hù} (chute) \zh{huán} (montée) \zh{jìng} (chute).
            \item \zh{垃圾要分类} : \zh{lājī} (1+1) \zh{yào} (4e) \zh{fēn} (1er) \zh{lèi} (4e).
            \item \zh{水很干净，我们种树} : \zh{shuǐ} (V) \zh{hěn} (\textbf{demi-3e, bas}) \zh{gān} (1) \zh{jìng} (4) // \zh{wǒmen} (3+neutre) \zh{zhòng} (4) \zh{shù} (4).
            \item \zh{中国保护大熊猫} : \zh{Zhōng} (1) \zh{guó} (2) \zh{bǎo} (bas) \zh{hù} (4) \zh{dà} (4) \zh{xióng} (2) \zh{māo} (1).
        \end{itemize}
\end{itemize}
\emph{Si tu as identifié le changement de \zh{污染} \zh{rǎn} en 2e ton devant \zh{水} ET la stabilité du demi-3e ton sur \zh{很} devant \zh{干}, tu as maîtrisé le sable tonal du 3e ton !}
\end{corrige}

\coursec{Jeux de rôle et débat guidé}

Après avoir travaillé la prononciation des tons spécifiques au vocabulaire de l'environnement, nous passons maintenant à la mise en situation communicative. L'objectif est de mobiliser le lexique, les structures grammaticales (\cle{应该}, \cle{觉得}, \cle{因为}...\cle{所以}) et les formules fonctionnelles vues précédemment pour interagir de façon fluide et culturellement appropriée. Nous allons progresser du jeu de rôle guidé vers le débat semi-libre, pour finir par l'exposé individuel évalué.

\courssub{Jeu de rôle 1 : Journaliste et Maire — La propreté du quartier}

\begin{methode}[Consignes du jeu de rôle]
\begin{enumerate}
    \item \textbf{Rôles :} Élève A = Journaliste (记者 \zh{记者} \emph{jìzhě}) ; Élève B = Maire (市长 \zh{市长} \emph{shìzhǎng}) ou Chef de quartier.
    \item \textbf{Situation :} Une émission de radio locale. Le journaliste interroge le maire sur l'état de la propreté dans le quartier (Yaoundé, Douala, Buea...) et les mesures prévues.
    \item \textbf{Structures cibles :} \zh{你怎么看...？} (Qu'en pensez-vous ?) ; \zh{我们应该...} (Nous devrions...) ; \zh{政府打算...} (Le gouvernement compte...).
    \item \textbf{Durée :} 2 à 3 minutes par binôme. Inversez les rôles.
\end{enumerate}
\end{methode}

\begin{textebox}[Modèle de dialogue : Journaliste — Maire]
\textbf{记者}：市长你好！今天我们谈谈咱们社区的环境卫生。\textbf{你怎么看现在的垃圾问题？}\\
\textbf{市长}：你好。\textbf{我觉得}垃圾确实有点多，\textbf{因为}有些居民乱扔垃圾。\\
\textbf{记者}：\textbf{那我们应该怎么办？}\\
\textbf{市长}：\textbf{我们应该}增加垃圾桶，\textbf{还要}加强宣传教育。\textbf{政府打算}下个月组织一次大扫除。\\
\textbf{记者}：很好！谢谢市长接受采访。\\
\textbf{市长}：不客气，谢谢关注。
\end{textebox}

\begin{center}
\begin{tabularx}{\linewidth}{|X|X|X|X|}
\hline \rowcolor{popblueL}
\textbf{Caractères} & \textbf{Pinyin} & \textbf{Nature} & \textbf{Traduction} \\
\hline
社区 & shèqū & nom & quartier, communauté \\
环境卫生 & huánjìng wèishēng & nom & hygiène de l'environnement / salubrité \\
居民 & jūmín & nom & habitant, résident \\
乱扔 & luàn rēng & verbe & jeter n'importe où \\
增加 & zēngjiā & verbe & augmenter, ajouter \\
垃圾桶 & lājītǒng & nom & poubelle \\
宣传教育 & xuānchuán jiàoyù & nom & sensibilisation, éducation civique \\
组织 & zǔzhī & verbe & organiser \\
大扫除 & dàsǎochú & nom & grand nettoyage (collectif) \\
采访 & cǎifǎng & verbe/nom & interviewer / interview \\
\hline
\end{tabularx}
\end{center}

\begin{exemplebox}[Variantes utiles pour le journaliste]
\begin{itemize}
\item \zh{请问，市里有什么新措施？} \emph{Qǐngwèn, shì lǐ yǒu shénme xīn cuòshī?} (Madame le Maire, quelles nouvelles mesures la ville prévoit-elle ?)
\item \zh{市民对空气质量不满意，你怎么看？} \emph{Shìmín duì kōngqì zhìliàng bù mǎnyì, nǐ zěnme kàn?} (Les citoyens sont insatisfaits de la qualité de l'air, quel est votre avis ?)
\end{itemize}
\end{exemplebox}

\begin{experience}[Pratique : Interview croisée]
En binôme, préparez 3 questions supplémentaires sur : \zh{噪音污染} (pollution sonore), \zh{排水系统} (système d'évacuation), \zh{公园绿化} (végétalisation des parcs). Jouez la scène devant la classe. Les observateurs notent l'usage de \zh{你怎么看} et \zh{我们应该}.
\end{experience}

\courssub{Jeu de rôle 2 : Deux Voisins — Déchets dans la rue et solutions}

\begin{methode}[Consignes du jeu de rôle]
\begin{enumerate}
    \item \textbf{Rôles :} Élève A = Voisin 1 (inquiet, propose des solutions) ; Élève B = Voisin 2 (au début sceptique, puis d'accord).
    \item \textbf{Situation :} Devant leur immeuble ou dans la rue (quartier Mokolo, Nlongkak, etc.). Ils voient des déchets par terre.
    \item \textbf{Formules de politesse/relance :} \zh{哎，你看...} (Hé, regarde...) ; \zh{是啊，真是的。} (Oui, c'est vrai, c'est n'importe quoi.) ; \zh{要不我们...？} (Et si nous... ?).
\end{enumerate}
\end{methode}

\begin{textebox}[Modèle de dialogue : Deux Voisins]
\textbf{邻居甲}：哎，老王，\textbf{你看咱们楼下的垃圾}，真臭！\\
\textbf{邻居乙}：是啊，\textbf{真是的}。谁扔的这么没公德心？\\
\textbf{邻居甲}：\textbf{我觉得}大家环保意识不够。\textbf{要不我们}给物业打个电话？\\
\textbf{邻居乙}：好主意。\textbf{还得}大家自己注意，\textbf{不能乱扔}。\\
\textbf{邻居甲}：对，\textbf{以后}我们带头捡一下，带个好头。\\
\textbf{邻居乙}：行，\textbf{那就这么说定了}！
\end{textebox}

\begin{center}
\begin{tabularx}{\linewidth}{|X|X|X|X|}
\hline \rowcolor{popgreenL}
\textbf{Caractères} & \textbf{Pinyin} & \textbf{Nature} & \textbf{Traduction} \\
\hline
邻居 & línjū & nom & voisin \\
楼下 & lóuxià & nom & en bas de l'immeuble \\
臭 & chòu & adj & sentir mauvais, puant \\
公德心 & gōngdéxīn & nom & sens civique, conscience citoyenne \\
环保意识 & huánbǎo yìshi & nom & conscience écologique \\
物业 & wùyè & nom & gestion d'immeuble / syndic \\
打电话 & dǎ diànhuà & v + obj & passer un coup de fil \\
还得 & hái děi & adv + verbe & il faut aussi / encore faut-il que \\
带头 & dàitóu & verbe & donner l'exemple, montrer la voie \\
带个好头 & dài gè hǎo tóu & expr & donner le bon exemple \\
这么说定了 & zhème shuō dìng le & expr & c'est entendu / c'est dit \\
\hline
\end{tabularx}
\end{center}

\begin{attention}[Piège culturel : \zh{真是的} vs \zh{真的}]
\zh{真是的} (zhēn shì de) exprime l'agacement, le reproche (« C'est n'importe quoi ! », « Vraiment ! »). C'est une interjection courante à l'oral.
\zh{真的} (zhēn de) signifie « vraiment » (adverbe de degré) ou « c'est vrai ? » (question).
Ne confondez pas : \zh{真是的，你怎么又迟到了！} (C'est n'importe quoi, tu es encore en retard !) $\neq$ \zh{你真的来了？} (Tu es vraiment venu ?).
\end{attention}

\begin{experience}[Improvisation guidée]
Changez le sujet : \zh{楼道里堆旧家具} (vieux meubles dans le couloir), \zh{邻居晚上卡拉OK太吵} (karaoké trop fort le soir). Utilisez \zh{要不我们...} et \zh{还得...}.
\end{experience}

\courssub{Débat guidé en groupes de 4 : « Faut-il planter plus d'arbres à l'école ? »}

C'est l'activité centrale de cette leçon. Elle structure l'interaction complexe : tour de parole, argumentation, synthèse.

\begin{methode}[Organisation du débat — Rôles et Déroulement]
\begin{enumerate}
    \item \textbf{Groupes de 4 élèves} : 
    \begin{itemize}
        \item \textbf{Animateur (主持人 \zh{主持人} \emph{zhǔchírén})} : Lance le sujet, distribue la parole, reformule (\zh{你怎么看？}, \zh{你呢？}).
        \item \textbf{Pour (正方 \zh{正方} \emph{zhèngfāng})} : Défend la plantation (\zh{我觉得应该...}, \zh{因为...所以...}).
        \item \textbf{Contre (反方 \zh{反方} \emph{fǎnfāng})} : Soulève des obstacles (\zh{我不同意}, \zh{可是...}, \zh{不过...}).
        \item \textbf{Rapporteur (记录员/汇报人 \zh{记录员} \emph{jìlùyuán})} : Note les arguments, fait la synthèse finale (\zh{所以，我们认为...}).
    \end{itemize}
    \item \textbf{Sujet :} \zh{学校要不要多种树？} (L'école doit-elle planter plus d'arbres ?).
    \item \textbf{Temps :} 5 min préparation (notes en chinois), 8 min débat, 2 min rapport oral par le rapporteur.
    \item \textbf{Grille de soutien (affichée au tableau) :} Utilisez ces connecteurs pour structurer vos phrases.
\end{enumerate}
\end{methode}

\begin{center}
\begin{tabularx}{\linewidth}{|X|X|X|}
\hline \rowcolor{poporangeL}
\textbf{Fonction} & \textbf{Formule chinoise + Pinyin} & \textbf{Français} \\
\hline
Donner son avis & \zh{我觉得...} (Wǒ juéde...) & Je pense que... \\
Approuver & \zh{我同意。/ 赞成。} (Wǒ tóngyì. / Zànchéng.) & Je suis d'accord. / J'approuve. \\
Désapprouver & \zh{我不同意。/ 我反对。} (Wǒ bù tóngyì. / Wǒ fǎnduì.) & Je ne suis pas d'accord. / Je m'oppose. \\
Nuancer / Contraster & \zh{可是... / 不过...} (Kěshì... / Bùguò...) & Mais... / Cependant... \\
Relancer / Inclure & \zh{你呢？} (Nǐ ne?) & Et toi ? / Et vous ? \\
Conclure & \zh{所以...} (Suǒyǐ...) & Donc / C'est pourquoi... \\
\hline
\end{tabularx}
\end{center}

\begin{popfigure}
\begin{tikzpicture}[
    node distance=1.2cm and 1.5cm,
    every node/.style={align=center, font=\small},
    role/.style={rectangle, rounded corners, draw=popdark, thick, fill=popcream, minimum width=2.2cm, minimum height=1.5cm},
    arrow/.style={-Stealth, thick, popblue},
    labelarrow/.style={midway, fill=white, inner sep=1pt, font=\tiny, text=popdark}
]
% Nœuds
\node[role, align=center] (anim){Animateur\\ \zh{主持人}\\\zh{你怎么看？}};
\node[role, right=of anim, align=center] (pour){Pour\\ \zh{正方}\\\zh{我觉得应该...}};
\node[role, below=of pour, align=center] (contre){Contre\\ \zh{反方}\\\zh{我不同意，可是...}};
\node[role, left=of contre, align=center] (rap){Rapporteur\\ \zh{记录员}\\\zh{所以，我们...}};

% Flèches
\draw[arrow] (anim) -- node[labelarrow] {Lance le sujet} (pour);
\draw[arrow] (pour) -- node[labelarrow, right] {Argumente} (contre);
\draw[arrow] (contre) -- node[labelarrow] {Réplique / Nuance} (rap);
\draw[arrow] (rap) -- node[labelarrow] {Synthèse} (anim);
\draw[arrow, dashed] (anim) -- node[labelarrow, above] {Relance: \zh{你呢？}} (contre);
\draw[arrow, dashed] (pour) .. controls +(up:1cm) and +(up:1cm) .. node[labelarrow, above] {Échange libre} (rap);
\end{tikzpicture}
\legende{Organigramme du débat en groupe : flux de parole et formules clés}
\end{popfigure}

\begin{textebox}[Modèle de débat (transcription simplifiée)]
\textbf{主持人}：同学们好，今天我们讨论：\textbf{学校要不要多种树？} 小李，\textbf{你怎么看？}\\
\textbf{正方(小李)}：\textbf{我觉得应该多种树}。\textbf{因为}树能净化空气，\textbf{还能}遮阳，\textbf{所以}环境会更好。\\
\textbf{反方(小王)}：\textbf{我不同意}。\textbf{可是}学校场地小，种树没地方。\textbf{而且}落叶难扫，\textbf{反而}增加工作量。\\
\textbf{正方(小李)}：\textbf{不过}可以种小树苗，不占地。\\
\textbf{主持人}：小张，\textbf{你呢？}\\
\textbf{记录员(小张)}：\textbf{所以，我们认为}可以种树，\textbf{但是}要选合适的树种，\textbf{并且}安排值日生扫叶。
\end{textebox}

\begin{center}
\begin{tabularx}{\linewidth}{|X|X|X|X|}
\hline \rowcolor{poporangeL}
\textbf{Caractères} & \textbf{Pinyin} & \textbf{Nature} & \textbf{Traduction} \\
\hline
落叶 & luòyè & nom & feuilles mortes \\
值日生 & zhírìshēng & nom & élève de service / responsable du jour \\
树种 & shùzhǒng & nom & essence d'arbre, espèce d'arbre \\
合适 & héshì & adj & approprié, convenable \\
安排 & ānpái & verbe & arranger, planifier \\
扫叶 & sǎo yè & v + obj & balayer les feuilles \\
\hline
\end{tabularx}
\end{center}

\begin{exoresolu}[Structuration d'argumentation]
\textbf{Énoncé :} Réorganisez les segments suivants pour former un argument cohérent pour le rôle « Pour » (正方).
\begin{enumerate}
    \item \zh{所以，我赞成多种树。}
    \item \zh{因为树可以吸收二氧化碳，释放氧气。}
    \item \zh{我觉得学校应该多种树。}
    \item \zh{而且树木能美化校园环境。}
\end{enumerate}
\textbf{Solution détaillée :} L'ordre logique en chinois suit souvent le schéma : \textbf{Avis} $\rightarrow$ \textbf{Cause 1} $\rightarrow$ \textbf{Cause 2 (而且)} $\rightarrow$ \textbf{Conclusion (所以)}.
Ordre correct : \textbf{3 $\rightarrow$ 2 $\rightarrow$ 4 $\rightarrow$ 1}.
\zh{我觉得学校应该多种树。因为树可以吸收二氧化碳，释放氧气。而且树木能美化校园环境。所以，我赞成多种树。}
\end{exoresolu}

\courssub{Mini-exposé individuel : Problème environnemental local et solution}

Chaque élève prépare et présente un exposé court (4 phrases, ~45 secondes) devant la classe. C'est l'évaluation sommative de la séquence orale.

\begin{methode}[Plan en 4 phrases pour ton exposé]
Suis ce canevas rigoureux pour structurer ta prise de parole :
\begin{enumerate}
    \item \textbf{Sujet (Accroche) :} \zh{今天我要谈一谈...} (Aujourd'hui, je veux parler de...) + \textbf{Lieu/Problème}.
    \item \textbf{Problème (Constat) :} \zh{这里的...太多 / 很脏 / 很差。} (Ici, il y a trop de... / c'est très sale / mauvais.)
    \item \textbf{Avis (Opinion) :} \zh{我觉得... / 我认为...} (Je pense que... / J'estime que...).
    \item \textbf{Solution (Proposition) :} \zh{我们应该... / 建议...} (Nous devrions... / Je suggère...). Terminer par \zh{谢谢大家！}.
\end{enumerate}
\end{methode}

\begin{textebox}[Modèle d'exposé : Déchets à l'école]
\zh{今天我要谈一谈我们学校的环境。} \\
\zh{这里的垃圾太多，空气也不好。} \\
\zh{我觉得我们应该多种树，不要乱扔垃圾。} \\
\zh{谢谢大家！}
\end{textebox}

\begin{exemplebox}[Autres modèles selon le quartier (Cameroun)]
\begin{itemize}
\item \textbf{Problème : Érosion / Boues (Yaoundé, saison des pluies).} \\
\zh{今天我要谈一谈我们小区的路。这里的泥太多，路很滑。我觉得政府应该修好排水沟，铺柏油路。谢谢大家！}
\item \textbf{Problème : Fumée / Brûlage déchets (Douala, quartiers populaires).} \\
\zh{今天我要谈一谈空气污染。这里的烟太多，空气很脏。我觉得居民不应该烧垃圾，应该分类回收。谢谢大家！}
\item \textbf{Problème : Déforestation / Érosion (Buea, zones de colline).} \\
\zh{今天我要谈一谈我们山上的树。这里的树被砍太多，下雨土就流走。我觉得我们应该保护树，多种树。建议学校组织植树活动。谢谢大家！}
\end{itemize}
\end{exemplebox}

\begin{center}
\begin{tabularx}{\linewidth}{|X|X|X|X|}
\hline \rowcolor{poppinkL}
\textbf{Caractères} & \textbf{Pinyin} & \textbf{Nature} & \textbf{Traduction} \\
\hline
泥 & ní & nom & boue \\
滑 & huá & adj & glissant \\
排水沟 & páishuǐgōu & nom & fossé de drainage, caniveau \\
铺 & pū & verbe & poser, paver \\
柏油路 & bǎiyóu lù & nom & route goudronnée \\
烟 & yān & nom & fumée \\
烧 & shāo & verbe & brûler \\
分类回收 & fēnlèi huíshōu & v + obj & trier et recycler \\
砍 & kǎn & verbe & couper (arbre) \\
流走 & liúzǒu & verbe & s'écouler, partir (eau/boue) \\
保护 & bǎohù & verbe & protéger \\
植树 & zhíshù & v + obj / nom & planter des arbres / reboisement \\
水土流失 & shuǐtǔ liúshī & nom & érosion des sols \\
\hline
\end{tabularx}
\end{center}

\begin{experience}[Passage à l'oral — Évaluation par les pairs]
Chaque élève passe au tableau. La classe utilise une grille simplifiée (voir critères ci-dessous) pour noter sur 5 : Tons, Formules, Fluidité, Contenu. L'enseignant donne un feedback immédiat sur 1 point fort et 1 point à améliorer.
\end{experience}

\courssub{Critères d'évaluation orale (Grille enseignant / auto-évaluation)}

\begin{center}
\begin{tabularx}{\linewidth}{|X|X|X|}
\hline \rowcolor{poppurpleL}
\textbf{Critère} & \textbf{Indicateurs de réussite (Niveau HSK 3 / 1ère A)} & \textbf{Points d'attention} \\
\hline
\textbf{Exactitude des tons} & Les tons des mots-clés (垃圾 lājī, 环境 huánjìng, 应该 yīnggāi, 因为 yīnwèi, 所以 suǒyǐ) sont corrects. Les tons neutres (de, le, ma, ne) sont légers. & Confusion \zh{垃圾} (lājī) vs \zh{辣鸡} (làjī) : tons 1/1 vs 4/1. Sandhi ton 3 + ton 3 non appliqué : \zh{很好} $\rightarrow$ \emph{hén hǎo}. \\
\hline
\textbf{Usage des formules} & Utilisation spontanée de \zh{我觉得 / 我同意 / 我不同意 / 可是 / 所以 / 你呢？}. Au moins 3 formules différentes dans le débat. & Traduction mot-à-mot du français : \zh*{Je pense que} $\rightarrow$ \zh{我想} (incorrect pour l'opinion), \zh{Mais} $\rightarrow$ \zh{但是} (trop formel à l'oral vs \zh{可是}). \\
\hline
\textbf{Fluidité / Débit} & Phrases complètes, pauses aux bons endroits (groupes respiratoires). Pas de silence long $>$ 3s. Recherche de mots visible mais gérée (\zh{那个...}, \zh{怎么说...}). & Hésitations excessives, phrases télégraphiques (sujet + verbe seulement), lecture de notes. \\
\hline
\textbf{Participation au débat} & Réagit aux arguments des autres (pas de monologue). Utilise \zh{你呢？} pour inclure. Respecte le tour de parole. & Interrompt (voir \textbf{Attention} ci-dessous), ne regarde que ses notes, ne répond pas au sujet. \\
\hline
\textbf{Prononciation / Intelligibilité} & Voyelles (ü, ong, ang) et consonnes (j/q/x, zh/ch/sh/r, z/c/s) distinctes. \zh{ü} correct après j/q/x/y. & \zh{绿 lǜ} $\rightarrow$ lu, \zh{种 zhòng} $\rightarrow$ zong, \zh{树 shù} $\rightarrow$ su. \\
\hline
\end{tabularx}
\end{center}

\begin{attention}[Erreur culturelle majeure : Interrompre]
En classe de chinois, comme dans la société chinoise, \textbf{interrompre son interlocuteur (打断 \zh{打断} \emph{dǎduàn}) est perçu comme un manque de respect (不礼貌 \zh{不礼貌} \emph{bù lǐmào})}.
\begin{itemize}
    \item \textbf{Interdit :} Couper la parole pour dire \zh{我不同意} avant la fin de la phrase.
    \item \textbf{Correct :} Attendre la particule finale (\zh{。}, \zh{吗}, \zh{呢}) ou la pause. Signaler son intention par un hochement de tête, \zh{嗯} (èn), ou lever légèrement la main. L'animateur donne la parole : \zh{请说} (Qǐng shuō).
\end{itemize}
\end{attention}

\begin{savaistu}[Les exposés oraux (口头报告 \zh{口头报告} \emph{kǒutóu bàogào}) en Chine]
Dans le système éducatif chinois (de l'école primaire au lycée), l'exposé oral est une compétence évaluée aussi sérieusement que l'écrit. Dès le collège, les élèves préparent des \zh{演讲} (yǎnjiǎng, discours) de 3 à 5 minutes sur des sujets de société, d'histoire ou de sciences. Les concours d'éloquence (演讲比赛) sont très populaires. La structure standard enseignée est : \zh{开头} (intro) $\rightarrow$ \zh{主体} (corps : arguments + exemples) $\rightarrow$ \zh{结尾} (conclusion + appel à l'action). C'est exactement le plan en 4 phrases que nous venons de travailler !
\end{savaistu}

\begin{exercice}[Entraînement final : Prépare ton passage]
\begin{enumerate}
    \item Choisis un problème réel près de chez toi (marché, école, rue, rivière).
    \item Rédige ton exposé de 4 phrases en caractères + pinyin sur ton cahier.
    \item Entraîne-toi à le dire sans notes, en respectant les tons et le débit.
    \item Échange avec un camarade : il joue l'animateur, il te demande \zh{你怎么看？}, tu réponds par ton exposé.
\end{enumerate}
\end{exercice}

\begin{corrige}[Exercice 1]
\textbf{Sujet :} Déforestation / Érosion colline (Buea).
\zh{今天我要谈一谈我们山上的树。} (Aujourd'hui je parle des arbres sur la colline.)
\zh{这里的树被砍太多，下雨土就流走。} (Ici on coupe trop d'arbres, à la pluie la terre part.)
\zh{我觉得我们应该保护树，多种树。} (Je pense qu'il faut protéger les arbres, en planter plus.)
\zh{建议学校组织植树活动。谢谢大家！} (Je suggère que l'école organise une activité de plantation. Merci !)
\end{corrige}

\newpage

\coursec{Activité d'intégration}

\courssub{Objectifs de l'activité}
Cette activité d'intégration vise à mobiliser l'ensemble des savoirs et savoir-faire travaillés lors de la leçon 17 dans une situation de communication authentique et stimulante. À l'issue de ce débat, l'élève doit être capable de :
\begin{itemize}
    \item \textbf{Mobiliser le lexique} de l'environnement (\cle{环境}、\cle{垃圾}、\cle{回收}、\cle{污染}、\cle{保护}、\cle{绿化} etc.) de façon pertinente et spontanée.
    \item \textbf{Utiliser les formules fonctionnelles} pour structurer un échange oral : présenter un sujet (\cle{今天我们讨论……}), donner son avis (\cle{我觉得……}、\cle{我认为……}), exprimer accord/désaccord (\cle{我同意}、\cle{我不同意}、\cle{可是/但是}), relancer (\cle{你怎么看？}、\cle{还有什么想法？}) et conclure (\cle{所以，我们应该……}).
    \item \textbf{Maîtriser la prononciation} : tons corrects sur le vocabulaire cible, liaison naturelle, intonation interrogative/affirmative/impérative.
    \item \textbf{Interagir en temps réel} : écouter l'interlocuteur, réagir à ses arguments, construire une argumentation cohérente à plusieurs voix.
    \item \textbf{Produire une trace écrite} synthétique (extension) reprenant les arguments clés sous forme de 4 phrases structurées.
\end{itemize}

\courssub{Situation}
\begin{textebox}[Affiche de la Mairie de Yaoundé — Consultation Jeunesse]
\zh{雅温得市政府青年咨询通告} \\[4pt]
\zh{主题：共建整洁雅温得，青年献计献策} \\[4pt]
\zh{亲爱的同学们：} \\[4pt]
\zh{为了改善我们社区的环境卫生，市政府诚邀中学生代表参加“整洁社区”咨询会。} \\[4pt]
\zh{议题：1. 垃圾分类与回收利用；2. 减少塑料袋污染；3. 增加社区绿化面积。} \\[4pt]
\zh{形式：分组辩论，每组 3 分钟，陈述观点并提出建议。} \\[4pt]
\zh{时间：下周一上午 9 点，市政府大会议室。} \\[4pt]
\zh{欢迎积极发言，为雅温得出一份力！} \\[8pt]
\textit{Yǎwēndé shìzhèngfǔ qīngnián zīxún tōnggào. — Zhǔtí: Gòngjiàn zhěngjié Yǎwēndé, qīngnián xiànjì xiàncè. Qīn'ài de tóngxuémen: Wèile gǎishàn wǒmen shèqū de huánjìng wèishēng, shìzhèngfǔ chéngyāo zhōngxuéshēng dàibiǎo cānjiā “zhěngjié shèqū” zīxúnhuì. Yìtí: 1. Lājī fēnlèi yǔ huíshōu lìyòng; 2. Jiǎnshǎo sùliàodài wūrǎn; 3. Zēngjiā shèqū lǜhuà miànjí. Xíngshì: fēnzǔ biànlùn, měizǔ 3 fēnzhōng, chénshù guāndiǎn bìng tíchū jiànyì. Shíjiān: xià zhōuyī shàngwǔ 9 diǎn, shìzhèngfǔ dà huìyìshì. Huānyíng jījí fāyán, wèi Yǎwēndé chū yī fèn lì!} \\[4pt]
\textbf{Traduction :} Avis de consultation jeunesse de la Mairie de Yaoundé. Thème : « Construisons ensemble un Yaoundé propre, la jeunesse propose des idées ». Chers élèves, pour améliorer l'hygiène de notre quartier, la mairie invite des représentants lycéens à une consultation « Quartier propre ». Sujets : 1. Tri et recyclage des déchets ; 2. Réduction de la pollution par les sacs plastiques ; 3. Augmentation des espaces verts du quartier. Format : débat par groupes, 3 minutes par groupe, exposer points de vue et propositions. Horaire : lundi prochain 9h, grande salle de conférence de la mairie. Bienvenue à tous pour prendre la parole et contribuer pour Yaoundé !
\end{textebox}

\courssub{Consignes}
\begin{enumerate}
    \item \textbf{Formation des groupes (5 min) :} La classe est divisée en groupes de 4 élèves. Chaque groupe tire au sort un des trois sujets de l'affiche (Tri/Recyclage, Sacs plastiques, Espaces verts). Les rôles sont attribués : \textbf{Animateur} (lance le sujet, relance, gère le temps), \textbf{Orateur A} (défend le pour / solution 1), \textbf{Orateur B} (défend le contre / solution alternative / nuances), \textbf{Rapporteur} (prend notes, synthétise, conclut).
    \item \textbf{Préparation écrite (15 min) :} Chaque élève rédige \textbf{ses 4 phrases clés} (brouillon) en caractères + pinyin sur une fiche : 1 phrase d'introduction/argument principal, 1 phrase avec \cle{因为……所以……}, 1 phrase avec \cle{可是/但是} (concession/opposition), 1 phrase de conclusion avec \cle{应该}. L'animateur prépare 2 questions de relance (\cle{你怎么看？}、\cle{你赞成吗？}). Le rapporteur prépare la structure de conclusion (\cle{所以，我们建议……}).
    \item \textbf{Répétition phonétique (10 min) :} Lecture à voix basse dans le groupe. Focus sur les \textbf{tons} du vocabulaire technique (\cle{huánjìng} 2-4, \cle{lājī} 1-1, \cle{huíshōu} 2-1, \cle{wūrǎn} 1-3, \cle{lǜhuà} 4-4). L'animateur vérifie les tons de chaque membre ; correction par les pairs.
    \item \textbf{Passage devant la classe (3 min/groupe) :} Le groupe se place face à la classe. L'animateur ouvre : \zh{大家好，今天我们讨论……你怎么看？}. Déroulement libre mais structuré. Le professeur note sur la grille d'évaluation (voir Ressources).
    \item \textbf{Évaluation par les pairs (5 min) :} Les groupes spectateurs disposent d'une \textbf{grille d'écoute} (cases à cocher : \textit{Formules entendues : 我觉得 / 我同意 / 我不同意 / 可是 / 所以 / 你怎么看 ?} ; \textit{Tons : globalement corrects / hésitants / erreurs fréquentes} ; \textit{Interaction : naturelle / forcée}). Remise au groupe présentateur.
    \item \textbf{Retour professeur (5 min) :} Commentaire global sur la prononciation (tons sandhi, neutralisation), la richesse lexicale, la cohérence syntaxique (\cle{比}, \cle{比较}, \cle{最}), la fluidité de l'interaction.
    \item \textbf{Extension écrite (Devoir / 15 min en classe) :} Chaque élève rédige \textbf{proprement} son exposé personnel (4 phrases minimum) en caractères, pinyin et traduction française sur son cahier de chinois. Titre : \zh{我的环保建议：[Sujet choisi]}. Cela constitue la trace écrite évaluée.
\end{enumerate}

\courssub{Ressources à mobiliser}
\begin{center}
\begin{tabularx}{\linewidth}{|X|X|X|X|}
\hline
\rowcolor{popblueL} \textbf{Caractères} & \textbf{Pinyin} & \textbf{Nature} & \textbf{Traduction} \\
\hline
环境 & huánjìng & n. & environnement \\
垃圾 & lājī & n. & déchets, ordures \\
垃圾分类 & lājī fēnlèi & v./n. & tri des déchets \\
回收 & huíshōu & v. & recycler, récupérer \\
利用 & lìyòng & v. & utiliser, valoriser \\
塑料袋 & sùliàodài & n. & sac en plastique \\
污染 & wūrǎn & v./n. & polluer, pollution \\
减少 & jiǎnshǎo & v. & réduire \\
绿化 & lǜhuà & v./n. & verdir, espaces verts \\
植树 & zhíshù & v. & planter des arbres \\
空气 & kōngqì & n. & air \\
新鲜 & xīnxiān & adj. & frais, pur \\
脏 & zāng & adj. & sale \\
干净 & gānjìng & adj. & propre \\
保护 & bǎohù & v. & protéger \\
建议 & jiànyì & v./n. & suggérer, suggestion \\
应该 & yīnggāi & v. mod. & devoir, devrait \\
因为……所以…… & yīnwèi……suǒyǐ…… & conj. & parce que… donc… \\
可是 / 但是 & kěshì / dànshì & conj. & mais, cependant \\
我觉得 & wǒ juéde & expr. & je pense que, j'ai l'impression que \\
我认为 & wǒ rènwéi & expr. & j'estime que, je considère que \\
我同意 & wǒ tóngyì & expr. & je suis d'accord \\
我不同意 & wǒ bù tóngyì & expr. & je ne suis pas d'accord \\
你怎么看？ & nǐ zěnme kàn? & expr. & Qu'en penses-tu ? / Quel est ton avis ? \\
还有什么想法？ & hái yǒu shénme xiǎngfǎ? & expr. & D'autres idées ? \\
所以 & suǒyǐ & adv. & donc, par conséquent \\
\hline
\end{tabularx}
\end{center}

\vspace{0.5cm}
\noindent\textbf{Structures grammaticales clés (rappel) :}
\begin{itemize}
    \item \textbf{Comparaison :} \cle{A 比 B + Adj.} (\zh{植树比不植树好。} — Planter des arbres est mieux que de ne pas en planter.) / \cle{A 没有 B + Adj.} (moins... que).
    \item \textbf{Cause - Conséquence :} \cle{因为……所以……} (\zh{因为塑料袋污染大，所以我们应该少用。}).
    \item \textbf{Opposition / Concession :} \cle{虽然……，但是/可是……} (\zh{虽然分类麻烦，但是很有必要。}) ou simple \cle{可是/但是} en milieu de phrase.
    \item \textbf{Obligation / Conseil :} \cle{应该 / 应该要 + V.} (\zh{我们应该保护环境。}).
    \item \textbf{Particule \cle{了} (aspect accompli / changement d'état) :} \zh{现在环境变好了。} (\cle{变好} \textit{biànhǎo} s'améliorer).
\end{itemize}

\vspace{0.3cm}
\noindent\textbf{Repères phonétiques (Tons du thème) :}
\begin{itemize}
    \item Mots en \textbf{1er ton} (haut plat) : \cle{分} (fēn), \cle{新} (xīn), \cle{花} (huā), \cle{脏} (zāng), \cle{圾} (jī), \cle{收} (shōu), \cle{该} (gāi).
    \item Mots en \textbf{2e ton} (montant) : \cle{环} (huán), \cle{回} (huí), \cle{绿} (lǜ \rightarrow attention sandhi ! \cle{lǜ} devient \cle{lú} devant 4e ton, mais \cle{lǜ} devant 1/2/3/neutre. Ici \cle{lǜhuà} 4-4 \rightarrow \cle{lúhuà}).
    \item Mots en \textbf{3e ton} (bas descendant-remontant) : \cle{类} (lèi), \cle{污} (wū), \cle{减} (jiǎn), \cle{种} (zhǒng), \cle{染} (rǎn), \cle{少} (shǎo). \textbf{Sandhi 3e ton + 3e ton} : \cle{jiǎnshǎo} \rightarrow \cle{jiánshǎo} ; \cle{wūrǎn} reste 1-3 (1er + 3e) ; \cle{tǔrǎng} (土壤) 3-3 \rightarrow 2-3.
    \item Mots en \textbf{4e ton} (descendant) : \cle{境} (jìng), \cle{用} (yòng), \cle{化} (huà), \cle{护} (hù), \cle{议} (yì), \cle{看} (kàn).
    \item \textbf{Ton neutre} : \cle{了} (le), \cle{吗} (ma), \cle{呢} (ne), \cle{的} (de), \cle{么} (me).
\end{itemize}

\courssub{Proposition de production et corrigé commenté}
\begin{exoresolu}[Modèle de mini-débat — Groupe 1 : Sujet « Réduire les sacs plastiques »]
\textbf{Scénario :} Animateur (A), Orateur 1 — Pour interdiction stricte (B), Orateur 2 — Pour transition progressive / alternatives (C), Rapporteur (D). \\[6pt]
\tcblower
\textbf{Production orale modèle (Caractères | Pinyin | Traduction)} \\[4pt]
\zh{A: 大家好，今天我们讨论减少塑料袋污染。你怎么看？} \\
\textit{Dàjiā hǎo, jīntiān wǒmen tǎolùn jiǎnshǎo sùliàodài wūrǎn. Nǐ zěnme kàn?} \\
« Bonjour à tous, aujourd'hui nous discutons de la réduction de la pollution par les sacs plastiques. Qu'en penses-tu ? » \\[6pt]
\zh{B: 我觉得塑料袋污染很严重。因为塑料袋分解很慢，所以污染土壤和水。我建议超市禁止免费发塑料袋，顾客自己带布袋子。} \\
\textit{Wǒ juéde sùliàodài wūrǎn hěn yánzhòng. Yīnwèi sùliàodài fēnjiě hěn màn, suǒyǐ wūrǎn tǔrǎng hé shuǐ. Wǒ jiànyì chāoshì jìnzhǐ miǎnfēi fā sùliàodài, gùkè zìjǐ dài bù dàizi.} \\
« Je trouve que la pollution des sacs plastiques est très grave. Parce que les sacs plastiques se décomposent très lentement, ils polluent le sol et l'eau. Je suggère que les supermarchés interdisent la distribution gratuite de sacs plastiques, les clients apportent leurs propres sacs en tissu. » \\[6pt]
\zh{C: 我同意污染严重，可是突然禁止不现实。很多老人习惯用塑料袋，不方便。我觉得比较好的办法是：超市收费卖环保袋，慢慢改变习惯。} \\
\textit{Wǒ tóngyì wūrǎn yánzhòng, kěshì tūrán jìnzhǐ bù xiànshí. Hěnduō lǎorén xíguàn yòng sùliàodài, bù fāngbiàn. Wǒ juéde bǐjiào hǎo de bànfǎ shì: chāoshì shōufèi mài huánbǎo dài, mànmàn gǎibiàn xíguàn.} \\
« Je suis d'accord que la pollution est grave, mais une interdiction soudaine n'est pas réaliste. Beaucoup de personnes âgées ont l'habitude d'utiliser des sacs plastiques, c'est pas commode. Je trouve que la méthode relativement bonne est : les supermarchés vendent des sacs écologiques payants, changent les habitudes progressivement. » \\[6pt]
\zh{B: 可是收费也不行，穷人买不起。保护环境不能只靠收费。政府应该立法，强制用可降解袋子。} \\
\textit{Kěshì shōufèi yě bùxíng, qióngrén mǎi bu qǐ. Bǎohù huánjìng bùnéng zhǐ kào shōufèi. Zhèngfǔ yīnggāi lìfǎ, qiángzhì yòng kě jiàngjiě dàizi.} \\
« Mais le payant ne marche pas non plus, les pauvres ne peuvent pas acheter. Protéger l'environnement ne peut pas reposer seulement sur le payant. Le gouvernement devrait légiférer, imposer l'utilisation de sacs biodégradables. » \\[6pt]
\zh{A: 还有什么想法？我们时间不多了。} \\
\textit{Hái yǒu shénme xiǎngfǎ? Wǒmen shíjiān bù duō le.} \\
« D'autres idées ? Nous n'avons plus beaucoup de temps. » \\[6pt]
\zh{D: 所以，我们建议：第一，政府立法推广可降解袋子；第二，超市提供便宜的布袋子；第三，学校加强环保教育。谢谢大家！} \\
\textit{Suǒyǐ, wǒmen jiànyì: dì yī, zhèngfǔ lìfǎ tuīguǎng kě jiàngjiě dàizi; dì èr, chāoshì tígōng piányi de bù dàizi; dì sān, xuéxiào jiāqiáng huánbǎo jiàoyù. Xièxiè dàjiā!} \\
« Donc, nous proposons : premièrement, le gouvernement légifère pour promouvoir les sacs biodégradables ; deuxièmement, les supermarchés fournissent des sacs en tissu bon marché ; troisièmement, les écoles renforcent l'éducation environnementale. Merci à tous ! »
\end{exoresolu}

\noindent\textbf{Commentaire des choix de langue (pour l'enseignant) :}
\begin{enumerate}
    \item \textbf{Lexique ciblé :} \cle{严重} (grave), \cle{分解} (se décomposer), \cle{土壤} (sol), \cle{禁止} (interdire), \cle{免费} (gratuit), \cle{布袋子} (sac en tissu), \cle{现实} (réaliste), \cle{习惯} (habitude), \cle{环保袋} (sac écolo), \cle{穷人} (pauvres), \cle{立法} (légiférer), \cle{强制} (imposer), \cle{可降解} (biodégradable), \cle{推广} (promouvoir). Niveau HSK 3-4, adapté 1ère A.
    \item \textbf{Structures fonctionnelles respectées :}
    \begin{itemize}
        \item Ouverture/Relance (A) : \cle{今天我们讨论……} + \cle{你怎么看？} / \cle{还有什么想法？}.
        \item Opinion argumentée (B) : \cle{我觉得……} + \cle{因为……所以……} + \cle{我建议……}.
        \item Accord nuancé / Opposition (C) : \cle{我同意……，可是……} + \cle{我觉得比较好的办法是……}.
        \item Contre-argument (B) : \cle{可是……} + \cle{不能只靠……} + \cle{应该……}.
        \item Conclusion structurée (D) : \cle{所以，我们建议：第一……第二……第三……}.
    \end{itemize}
    \item \textbf{Points de vigilance phonétique :}
    \begin{itemize}
        \item \cle{jiǎnshǎo} $\rightarrow$ \cle{jiánshǎo} (sandhi 3+3).
        \item \cle{lǜhuà} (si utilisé) $\rightarrow$ \cle{lúhuà} (sandhi 4+4 pour \cle{lǜ}).
        \item \cle{yánzhòng} (2-4), \cle{fēnjiě} (1-3), \cle{tǔrǎng} (3-3 $\rightarrow$ 2-3), \cle{xiànshí} (4-2), \cle{jiàngjiě} (4-3), \cle{lìfǎ} (4-3).
        \item Intonation montante finale pour \cle{你怎么看？} \cle{还有什么想法？}.
    \end{itemize}
    \item \textbf{Cohérence interactionnelle :} L'animateur (A) cadre et gère le temps. Les orateurs (B, C) réagissent explicitement aux propos du précédent (\cle{我同意……可是}, \cle{可是收费也不行}). Le rapporteur (D) synthétise les deux positions (législation + alternatives + éducation) sans en ajouter de nouvelle.
\end{enumerate}

\courssub{Bilan de l'activité}
\begin{aretenir}[Bilan compétences — Leçon 17]
\begin{itemize}
    \item \textbf{Expression orale continue :} L'élève est capable de produire un énoncé court (3-4 phrases), structuré (introduction, argument, conclusion), sur un sujet sociétal familier, en utilisant le vocabulaire spécifique de l'environnement.
    \item \textbf{Expression orale en interaction :} L'élève sait prendre la parole à son tour, écouter pour rebondir (\cle{我同意/不同意}, \cle{可是}), utiliser des formules de relance (\cle{你怎么看？}) et gérer la fin de tour (\cle{所以……}).
    \item \textbf{Qualité de la langue :} Prononciation des tons globalement correcte sur le lexique cible (vérification sandhi 3e ton, \cle{lǜ}). Utilisation correcte des connecteurs logiques (\cle{因为……所以……}, \cle{虽然……但是……}, \cle{比较}, \cle{应该}).
    \item \textbf{Compétence socioculturelle :} L'élève ancre son propos dans la réalité camerounaise (Yaoundé, marchés, habitudes locales, rôle de la mairie) tout en mobilisant des références comparatives Chine (tri sélectif obligatoire à Shanghai, interdiction sacs plastiques, grande campagne de plantation d'arbres).
    \item \textbf{Trace écrite :} La rédaction des 4 phrases (personnelles) permet de fixer les structures, de vérifier l'écriture des caractères (ordre des traits : \cle{环} \textit{radical 阝 (droite), ordre : gauche puis droite} ; \cle{污} \textit{radical 氵 (gauche), ordre : haut → bas, gauche → droite} ; \cle{降} \textit{haut en bas, 土 au milieu}) et de servir de support à l'évaluation formative.
\end{itemize}
\end{aretenir}

\begin{attention}[Erreurs fréquentes observées lors du débat — À corriger collectivement]
\begin{itemize}
    \item \textbf{Tons :} \cle{环境} prononcé \textit{huánjǐng} (3e ton au lieu de 4e sur \cle{境}) ; \cle{垃圾} prononcé \textit{lājǐ} (3e ton sur \cle{圾}) ; \cle{绿化} prononcé \textit{lǜhuà} sans sandhi (doit être \textit{lúhuà}) ; \cle{减少} sans sandhi (\textit{jiǎnshǎo} au lieu de \textit{jiánshǎo}).
    \item \textbf{Ordre des mots (Complément de lieu/temps) :} \textit{*我们在超市应该带袋子} $\rightarrow$ \textbf{我们应该在超市带袋子} (Modaux avant locatif).
    \item \textbf{\cle{比} mal placé :} \textit{*塑料袋比布袋子脏} (sens inverse souvent voulu : « le sac plastique est plus sale ») $\rightarrow$ \textbf{布袋子比塑料袋干净} (le sac en tissu est plus propre) ou \textbf{塑料袋比布袋子脏} (vérifier le sens visé !).
    \item \textbf{Confusion \cle{觉得} / \cle{认为} :} \cle{觉得} plus subjectif/sentiment ; \cle{认为} plus rationnel/argumenté. En débat formel, \cle{我认为} est plus approprié.
    \item \textbf{Oubli de \cle{了} pour le changement d'état :} \textit{*现在环境变好} $\rightarrow$ \textbf{现在环境变好\cle{了}}.
    \item \textbf{Traduction littérale du français :} \textit{*我同意 avec vous} (pas de \cle{avec} en chinois) $\rightarrow$ \textbf{我同意你的看法} / \textbf{我跟你一样想}.
\end{itemize}
\end{attention}

\begin{experience}[Prolongement possible — Projet interdisciplinaire SVT / Géographie]
Organiser une \textbf{sortie terrain} (ou enquête de quartier) : « Audit propreté de mon quartier ». Les élèves, par binômes, photographient 3 points noirs (décharge sauvage, caniveaux bouchés, absence poubelles) et 3 points positifs (arbres, poubelles de tri, panneau sensibilisation). En classe, ils présentent en chinois (5 phrases) leur diagnostic et une proposition concrète. Production d'une affiche bilingue (Chinois/Français) pour le CDI : \zh{保护环境，从我做起} — « Protégeons l'environnement, commençons par nous-mêmes ».
\end{experience}

\newpage

\coursec{Méthodes et exercices résolus}

\coursec{Leçon 17 — Expression orale : protection de l’environnement}

\courssub{Vocabulaire clé : l'environnement}

\begin{definition}[Lexique thématique — Environnement et protection (HSK 2-3)]
Le tableau ci-dessous regroupe le vocabulaire essentiel pour comprendre, parler et débattre sur l'environnement. Mémorise les tons, le(s) classificateur(s) associé(s) et la nature grammaticale.
\begin{center}
\begin{tabularx}{\linewidth}{|X|X|X|X|}
\hline
\rowcolor{popblueL} \textbf{Caractères} & \textbf{Pinyin} & \textbf{Nature} & \textbf{Traduction} \\
\hline
环境 & huánjìng & nom & environnement \\
保护 & bǎohù & verbe & protéger \\
污染 & wūrǎn & verbe / nom & polluer / pollution \\
垃圾 & lājī & nom & déchets, ordures \\
塑料袋 & sùliàodài & nom & sac en plastique (classif. \zh{个} / \zh{只}) \\
空气 & kōngqì & nom & air \\
严重 & yánzhòng & adjectif & grave, sérieux \\
骑自行车 & qí zìxíngchē & v + obj & faire du vélo \\
公共汽车 & gōnggòng qìchē & nom & bus, autobus (classif. \zh{辆}) \\
种树 & zhòng shù & v + obj & planter des arbres (classif. \zh{棵} pour l'arbre) \\
捡 & jiǎn & verbe & ramasser (par terre) \\
河 & hé & nom & rivière, fleuve \\
比如 & bǐrú & conjonction & par exemple \\
同意 & tóngyì & verbe & être d'accord (avec) \\
看法 & kànfǎ & nom & opinion, point de vue \\
责任 & zérèn & nom & responsabilité \\
雅温得 & Yǎwēndé & n. propre & Yaoundé (capitale du Cameroun) \\
应该 & yīnggāi & modal & devoir, il faut \\
可以 & kěyǐ & modal & pouvoir (permission/capacité) \\
觉得 & juéde & verbe & penser, trouver (avis personnel) \\
认为 & rènwéi & verbe & considérer que (plus formel) \\
因为 & yīnwèi & conjonction & parce que \\
所以 & suǒyǐ & conjonction & donc, par conséquent \\
在 & zài & préposition & à, dans (lieu d'action) \\
里 & lǐ & locatif & à l'intérieur (après nom de lieu) \\
棵 & kē & classif. & arbres, plantes à tige ligneuse \\
个 & gè & classif. & classificateur général (secours) \\
\hline
\end{tabularx}
\end{center}
\textbf{Astuce :} Associe chaque nom à son classificateur dès l'apprentissage : \zh{一棵树} (yì kē shù), \zh{一个塑料袋} (yí ge sùliàodài), \zh{一辆公共汽车} (yí liàng gōnggòng qìchē).
\end{definition}

\courssub{Dialogue modèle 1 : Donner son avis sur la pollution à Yaoundé}

\begin{textebox}[Dialogue : A et B parlent de la qualité de l'air]
\textbf{A：}\zh{你好！你觉得雅温得的空气怎么样？} (Nǐ hǎo! Nǐ juéde Yǎwēndé de kōngqì zěnme yàng ?)\\
\textbf{B：}\zh{我觉得雅温得的空气不太好。因为汽车太多了，所以空气污染很严重。} (Wǒ juéde Yǎwēndé de kōngqì bù tài hǎo. Yīnwèi qìchē tài duō le, suǒyǐ kōngqì wūrǎn hěn yánzhòng.)\\
\textbf{A：}\zh{我同意你的看法。我们应该多坐公共汽车，少开车。} (Wǒ tóngyì nǐ de kànfǎ. Wǒmen yīnggāi duō zuò gōnggòng qìchē, shǎo kāichē.)\\
\textbf{B：}\zh{对。保护环境，人人有责。} (Duì. Bǎohù huánjìng, rénrén yǒuzé.)
\end{textebox}

\begin{propriete}[Analyse linguistique du dialogue]
\begin{enumerate}
\item \textbf{Question d'opinion :} \zh{你觉得……怎么样？} (Nǐ juéde…… zěnme yàng ?) structure standard pour demander un avis.
\item \textbf{Réponse structurée :} \zh{我觉得……} (avis) + \zh{因为……所以……} (justification cause-effet). En chinois, les deux connecteurs coexistent souvent dans la même phrase complexe.
\item \textbf{Accord :} \zh{我同意你的看法。} (Wǒ tóngyì nǐ de kànfǎ.) « Je suis d'accord avec ton avis. » \zh{同意} est transitif direct.
\item \textbf{Conseil / Obligation morale :} \zh{应该} + V. La négation est \zh{不应该} (bù yīnggāi). Attention au sandhi de \zh{不} : devant \zh{应该} (4\textsuperscript{e} ton), \zh{不} se prononce \emph{bú} (2\textsuperscript{e} ton) $\rightarrow$ \emph{bú yīnggāi}.
\item \textbf{Chengyu (idiome) oral :} \zh{保护环境，人人有责。} (Bǎohù huánjìng, rénrén yǒuzé.) « Protéger l'environnement, chacun a sa part de responsabilité. » Phrase toute faite, très valorisée à l'oral.
\end{enumerate}
\end{propriete}

\courssub{Méthode 1 : Présenter un exposé simple en chinois}

\begin{methode}[Structurer un mini-exposé sur l'environnement (5-8 phrases)]
Pour présenter un exposé oral clair et noté, suis l'architecture en trois temps. Note les formules sur une fiche bristol (mots-clés seulement, pas le texte intégral).
\begin{enumerate}
\item \textbf{Saluer et annoncer le sujet (1 phrase).} \\
Formule : \zh{大家好！今天我要谈一谈……} (Dàjiā hǎo! Jīntiān wǒ yào tántán……) « Bonjour à tous ! Aujourd'hui, je vais parler de…… » \\
\emph{Note :} \zh{谈一谈} (tán yi tán) = reduplication du verbe « parler un peu / discuter de », adoucit le ton.
\item \textbf{Développer en 3 idées distinctes (3 phrases).} Une idée = une phrase courte (SVO). \\
\zh{问题} (Problème) : \zh{雅温得的塑料垃圾太多了。} (Yǎwēndé de sùliào lājī tài duō le.) \\
\zh{原因} (Cause) : \zh{因为大家乱扔垃圾。} (Yīnwèi dàjiā luàn rēng lājī.) \\
\zh{解决} (Solution) : \zh{我们应该少用塑料袋，多用布袋子。} (Wǒmen yīnggāi shǎo yòng sùliàodài, duō yòng bù dàizi.)
\item \textbf{Conclure et remercier (1-2 phrases).} \\
Appel à l'action : \zh{保护环境，从我做起。} (Bǎohù huánjìng, cóng wǒ zuò qǐ.) « Protéger l'environnement, commençons par moi. » \\
Remerciement : \zh{我的发言完了，谢谢大家！} (Wǒ de fāyán wán le, xièxie dàjiā !)
\end{enumerate}
\textbf{Réflexes oraux :} Parle lentement, marque les tons (main qui monte/descend), contact visuel avec le jury/camarades. Si tu bloques : \zh{那个……} (nà ge…… « euh… ») pour gagner du temps sans couper la prosodie.
\end{methode}

\begin{exoresolu}[Rédiger le plan d'un exposé : « Protéger l'environnement à Yaoundé »]
\textbf{Énoncé.} Rédige en chinois (caractères + pinyin + traduction) les 4 blocs de l'exposé.
\tcblower
\textbf{Solution détaillée.}
\begin{enumerate}
\item \textbf{Annonce :} \zh{今天我要谈一谈雅温得的环境保护。} (Jīntiān wǒ yào tántán Yǎwēndé de huánjìng bǎohù.) « Aujourd'hui, je vais parler de la protection de l'environnement à Yaoundé. »
\item \textbf{Problème :} \zh{雅温得的塑料垃圾太多了。} (Yǎwēndé de sùliào lājī tài duō le.) « Il y a trop de déchets plastiques à Yaoundé. » (\zh{太……了} = excès).
\item \textbf{Cause :} \zh{因为很多人乱扔塑料袋。} (Yīnwèi hěn duō rén luàn rēng sùliàodài.) « Parce que beaucoup de gens jettent des sacs plastiques n'importe où. »
\item \textbf{Solution :} \zh{我们应该自带购物袋。} (Wǒmen yīnggāi zì dài gòuwù dài.) « Nous devrions apporter nos propres sacs de courses. » (\zh{自带} = apporter soi-même).
\item \textbf{Conclusion :} \zh{保护环境，从我做起。谢谢大家！} (Bǎohù huánjìng, cóng wǒ zuò qǐ. Xièxie dàjiā !) « Protégeons l'environnement, commençons par nous-mêmes. Merci à tous ! »
\end{enumerate}
\textbf{Conclusion.} Un exposé réussi = 5 phrases courtes, correctes, reliées par la logique (Problème $\rightarrow$ Cause $\rightarrow$ Solution).
\end{exoresolu}

\courssub{Méthode 2 : Exprimer et justifier son opinion}

\begin{methode}[Formules fonctionnelles : Opiner, justifier, (dés)accorder]
\begin{enumerate}
\item \textbf{Introduire son avis (du plus familier au plus formel).} \\
\zh{我觉得……} (Wǒ juéde……) « Je trouve que… / Je pense que… » (spontané, oral). \\
\zh{我认为……} (Wǒ rènwéi……) « Je considère que… / J'estime que… » (formel, débat noté, écrit). \\
\zh{在我看来，……} (Zài wǒ kànlái……) « À mon avis… / Selon moi… » (topique, met le locuteur en thème).
\item \textbf{Justifier : la structure \zh{因为……，所以……}.} \\
Schéma : \zh{Sujet + 因为 + Cause , 所以 + Conséquence.} \\
Ex. : \zh{我觉得空气不好，因为工厂太多，所以污染严重。} \\
\emph{Attention :} En français, on évite « Parce que…, donc… » (redondant). En chinois, c'est la norme standard à l'oral et à l'écrit.
\item \textbf{Exprimer l'accord.} \\
\zh{我同意你的看法。} (Wǒ tóngyì nǐ de kànfǎ.) Standard. \\
\zh{对，我也是这么想的。} (Duì, wǒ yě shì zhème xiǎng de.) « Oui, je pense pareil. » (Naturel). \\
\zh{你说得很有道理。} (Nǐ shuō de hěn yǒu dàolǐ.) « Ton argument est très valable. » (Valorisant).
\item \textbf{Exprimer le désaccord (politesse chinoise : indirectivité).} \\
\zh{我不太同意……} (Wǒ bù tài tóngyì……) « Je ne suis pas tout à fait d'accord… » (\zh{不太} adoucit). \\
\zh{我觉得可能不太一样。} (Wǒ juéde kěnéng bù tài yíyàng.) « Je pense que c'est peut-être un peu différent. » \\
\zh{有道理，但是……} (Yǒu dàolǐ, dànshì……) « C'est vrai, mais… » (Concession + pivot).
\end{enumerate}
\end{methode}

\begin{exoresolu}[Compléter et traduire un dialogue d'opinion]
\textbf{Énoncé.} Complète la réplique de B avec une formule d'opinion + justification. Traduis ensuite la réplique complète de B.
A：\zh{你觉得我们学校的环境怎么样？}\\
B：\_\_\_\_\_\_\_\_\_\_\_\_，因为\_\_\_\_\_\_\_\_\_\_\_\_，所以\_\_\_\_\_\_\_\_\_\_\_\_。\\
A：\zh{我同意。我们应该一起捡垃圾。}
\tcblower
\textbf{Solution détaillée.}
\begin{enumerate}
\item \textbf{Complétion :} \zh{我觉得我们学校的环境很好，因为树很多，所以空气很新鲜。} (Wǒ juéde wǒmen xuéxiào de huánjìng hěn hǎo, yīnwèi shù hěn duō, suǒyǐ kōngqì hěn xīnxiān.)
\item \textbf{Traduction :} « Je trouve que l'environnement de notre école est très bien, parce qu'il y a beaucoup d'arbres, donc l'air est très frais. »
\item \textbf{Points de grammaire :}
\begin{itemize}
\item \zh{很好} : s'écrit \emph{hěn hǎo}, se prononce \emph{hén hǎo} (sandhi 3+3 $\rightarrow$ 2+3).
\item \zh{新鲜} (xīnxiān) : « frais » (air, nourriture), 1\textsuperscript{er} + 1\textsuperscript{er} ton.
\item \zh{树很多} : \zh{很} (degré) + adjectif verbe \zh{多}. Pas de verbe \zh{是} devant l'adjectif en chinois.
\end{itemize}
\end{enumerate}
\end{exoresolu}

\courssub{Dialogue modèle 2 : Mini-débat — Faut-il interdire les sacs plastiques ?}

\begin{textebox}[Débat guidé : 4 répliques fonctionnelles]
\textbf{A：}\zh{我认为我们不应该用塑料袋。} (Wǒ rènwéi wǒmen bù yīnggāi yòng sùliàodài.)\\
\textbf{B：}\zh{为什么？你能举个例子吗？} (Wèishénme? Nǐ néng jǔ ge lìzi ma?)\\
\textbf{A：}\zh{因为塑料袋会污染环境。比如，河里的塑料垃圾很多。} (Yīnwèi sùliàodài huì wūrǎn huánjìng. Bǐrú, hé lǐ de sùliào lājī hěn duō.)\\
\textbf{B：}\zh{我同意你的看法。总之，我们应该少用塑料袋。同学们，你们怎么看？} (Wǒ tóngyì nǐ de kànfǎ. Zǒngzhī, wǒmen yīnggāi shǎo yòng sùliàodài. Tóngxuémen, nǐmen zěnme kàn?)
\end{textebox}

\begin{propriete}[Analyse des formules de débat]
\begin{itemize}
\item \zh{不应该} (bù yīnggāi) : Négation du modal \zh{应该}. Prononciation \emph{bú yīnggāi} (sandhi \zh{不}).
\item \zh{举个例子} (jǔ ge lìzi) : « Donner un exemple ». \zh{个} est le classificateur de \zh{例子}.
\item \zh{比如} (bǐrú) : « Par exemple », introduit une illustration concrète. Très oral.
\item \zh{总之} (zǒngzhī) : « En résumé / En conclusion », marque la synthèse.
\item \zh{同学们，你们怎么看？} : Relance collective obligatoire pour montrer la compétence « interaction » au jury.
\end{itemize}
\end{propriete}

\courssub{Méthode 3 : Relancer la parole dans un débat (Interaction)}

\begin{methode}[Gérer son tour de parole et relancer l'échange]
\begin{enumerate}
\item \textbf{Demander l'avis / passer le relais :} \zh{你怎么看？} (Nǐ zěnme kàn?), \zh{你呢？} (Nǐ ne?), \zh{你同意吗？} (Nǐ tóngyì ma?).
\item \textbf{Demander une précision / une preuve :} \zh{为什么？} (Wèishénme?), \zh{你能举个例子吗？} (Nǐ néng jǔ ge lìzi ma?), \zh{具体怎么说？} (Jùtǐ zěnme shuō? « Concrètement, comment le dire ? »).
\item \textbf{Prendre la parole / Ajouter une idée :} \zh{我想补充一点。} (Wǒ xiǎng bǔchōng yì diǎn.), \zh{我还有另一个想法。} (Wǒ hái yǒu lìng yí ge xiǎngfǎ.).
\item \textbf{Conclure collectivement :} \zh{总之，我们都应该……} (Zǒngzhī, wǒmen dōu yīnggāi……), \zh{大家都同意吗？} (Dàjiā dōu tóngyì ma?).
\end{enumerate}
\textbf{Critère d'évaluation (Bac) :} L'interaction compte pour 30-40\% de la note. Un candidat qui ne fait que réciter son texte sans regarder ni questionner son partenaire perd des points. \textbf{À chaque fin de réplique, pose une question.}
\end{methode}

\begin{exoresolu}[Mini-débat : « Planter des arbres à l'école »]
\textbf{Énoncé.} Rédige un échange de 4 répliques. Sujet : \zh{我们学校应该种树吗？} A propose, B demande pourquoi, A donne 2 arguments (air, ombre), B conclut et relance la classe.
\tcblower
\textbf{Solution détaillée.}
\begin{enumerate}
\item \textbf{A (Proposition) :} \zh{我认为我们学校应该种树。} (Wǒ rènwéi wǒmen xuéxiào yīnggāi zhòng shù.)
\item \textbf{B (Relance cause) :} \zh{为什么？具体有什么好处？} (Wèishénme? Jùtǐ yǒu shénme hǎochù?)
\item \textbf{A (Arguments) :} \zh{因为树可以让空气变好，而且夏天可以乘凉。} (Yīnwèi shù kěyǐ ràng kōngqì biàn hǎo, érqiě xiàtiān kěyǐ chéngliáng.) \emph{Note :} \zh{而且} (érqiě) « de plus / en plus » pour lister arguments. \zh{乘凉} (chéngliáng) « profiter de la fraîcheur/ombre ».
\item \textbf{B (Synthèse + Relance) :} \zh{我同意。种树很好。总之，我们下周一起去种树吧？同学们，你们愿意吗？} (Wǒ tóngyì. Zhòng shù hěn hǎo. Zǒngzhī, wǒmen xià zhōu yìqǐ qù zhòng shù ba? Tóngxuémen, nǐmen yuànyì ma?) \emph{Note :} \zh{吧} (ba) suggestion douce ; \zh{愿意} (yuànyì) « vouloir bien / être volontaire ».
\end{enumerate}
\end{exoresolu}

\courssub{Méthode 4 : Prononcer correctement les tons du thème (Phonétique)}

\begin{methode}[Tons, sandhi et pièges phonétiques du vocabulaire environnement]
\begin{enumerate}
\item \textbf{Tons de base des mots-clés (à surligner sur ta fiche).} \\
\zh{环境} (huán\textcolor{poporange}{jíng} : 2+4), \zh{保护} (bǎo\textcolor{poporange}{hù} : 3+4), \zh{污染} (wū\textcolor{poporange}{rǎn} : 1+3), \zh{垃圾} (lā\textcolor{poporange}{jī} : 1+1), \zh{空气} (kōng\textcolor{poporange}{qì} : 1+4), \zh{严重} (yán\textcolor{poporange}{zhòng} : 2+4), \zh{应该} (yīng\textcolor{poporange}{gāi} : 1+1), \zh{可以} (kě\textcolor{poporange}{yǐ} : 3+3 $\rightarrow$ \emph{ké} yǐ).
\item \textbf{Règle du sandhi (changement de ton) : 3\textsuperscript{e} ton + 3\textsuperscript{e} ton.} \\
Le premier 3\textsuperscript{e} ton devient un 2\textsuperscript{e} ton (montant). \\
Ex. : \zh{很好} (hěn hǎo) $\rightarrow$ \emph{hé\textcolor{poporange}{n} hǎo} ; \zh{可以} (kě yǐ) $\rightarrow$ \emph{ké yǐ} ; \zh{保护} (bǎo hù) \textbf{ne change pas} (3+4).
\item \textbf{Sandhi de \zh{不} (bù) et \zh{一} (yī) — Crucial pour l'oral !} \\
\zh{不} (4\textsuperscript{e} ton) $\rightarrow$ \emph{bú} (2\textsuperscript{e} ton) \textbf{devant un 4\textsuperscript{e} ton} : \zh{不应该} \emph{bú yīnggāi}, \zh{不是} \emph{bú shì}, \zh{不好} \emph{bú hǎo}. \\
\zh{一} (1\textsuperscript{e} ton) $\rightarrow$ \emph{yí} (2\textsuperscript{e} ton) devant 4\textsuperscript{e} ton ; \emph{yì} (4\textsuperscript{e} ton) devant 1/2/3\textsuperscript{e} ton. Ex. : \zh{一个} \emph{yí ge}, \zh{一样} \emph{yí yàng}, \zh{一棵} \emph{yì kē}.
\item \textbf{Le 4\textsuperscript{e} ton « franc » :} \zh{护} (hù), \zh{气} (qì), \zh{污} (wū - 1\textsuperscript{er} mais fort), \zh{严} (yán - 2\textsuperscript{e} montant). Un 4\textsuperscript{e} ton « mou » (qui ne descend pas bas) sonne comme une hésitation ou une question.
\item \textbf{Entraînement en chaîne (chaîne tonale) :} Répète 3 fois de suite en tapant le rythme sur la table : \\
\zh{保护环境} (3-4 | 2-4) $\rightarrow$ \zh{人人有责} (2-2 | 3-2) $\rightarrow$ \zh{从我做起} (2-3 | 4-3).
\end{enumerate}
\end{methode}

\begin{exoresolu}[Phonétique : Identification et correction]
\textbf{Énoncé.} (a) Donne les tons (chiffres) et la prononciation réelle (sandhi appliqué) pour : \zh{环境}, \zh{保护}, \zh{垃圾}, \zh{应该}, \zh{不应该}, \zh{一棵树}. (b) Pourquoi \zh{空气} est-il un piège pour les francophones ? (c) Corrige la prononciation erronée : \emph{huán jīng}, \emph{báo hù}, \emph{là jǐ}.
\tcblower
\textbf{Solution détaillée.}
\begin{enumerate}
\item \textbf{Tableau des tons :}
\begin{center}
\begin{tabularx}{\linewidth}{|X|X|X|X|}
\hline
\rowcolor{popblueL} \textbf{Mot} & \textbf{Tons citation} & \textbf{Prononciation réelle (sandhi)} & \textbf{Note} \\
\hline
\zh{环境} & 2 + 4 & huán jìng & Stable \\
\zh{保护} & 3 + 4 & bǎo hù & Stable (3+4) \\
\zh{垃圾} & 1 + 1 & lā jī & Stable (hauts plats) \\
\zh{应该} & 1 + 1 & yīng gāi & Stable \\
\zh{不应该} & 4 + 1 + 1 & \textbf{bú} yīng gāi & \zh{不} $\rightarrow$ 2\textsuperscript{e} ton devant 4\textsuperscript{e} ton (\zh{应}) \\
\zh{一棵树} & 1 + 1 + 4 & \textbf{yì} kē shù & \zh{一} $\rightarrow$ 4\textsuperscript{e} ton devant 1\textsuperscript{er} ton (\zh{棵}) \\
\hline
\end{tabularx}
\end{center}
\item \textbf{\zh{空气} (kōng qì) :} Le francophone prononce souvent \emph{kōng qí} (montée finale à la française sur le « i »). Or \zh{气} est un 4\textsuperscript{e} ton \textbf{descendant brutal} (comme un « Non ! » sec). Il faut « tomber » sur le \emph{qì}.
\item \textbf{Corrections :} \emph{huán jìng} (pas \emph{jīng} 1\textsuperscript{er} ton = « cristal/miroir » \zh{镜}) ; \emph{bǎo hù} (pas \emph{báo} 2\textsuperscript{e} ton = « mince/faible » \zh{薄} ni \emph{bào} 4\textsuperscript{e} = « journal » \zh{报}) ; \emph{lā jī} (pas \emph{là} 4\textsuperscript{e} = « cire/épicer » \zh{蜡} ni \emph{jǐ} 3\textsuperscript{e} = « quelques » \zh{几}).
\end{enumerate}
\end{exoresolu}

\courssub{Méthode 5 : Construire des phrases orales correctes (Grammaire / Ordre des mots)}

\begin{methode}[Ordre des mots : Temps — Lieu — Manière — Verbe]
En chinois, les compléments circonstanciels se placent \textbf{avant} le verbe, dans un ordre fixe : \textbf{Sujet + Temps + Lieu + Manière/Modalité + Verbe + Complément d'objet}.
\begin{enumerate}
\item \textbf{Complément de temps :} Après le sujet, avant le lieu. \zh{我们明天去公园。} (Wǒmen míngtiān qù gōngyuán.)
\item \textbf{Complément de lieu :} Avec \zh{在} (zài), avant le verbe. \zh{我们在学校种树。} (Wǒmen zài xuéxiào zhòng shù.) \textbf{Jamais} \zh{我们种树在学校。} (Calque français).
\item \textbf{Locatif \zh{里} (lǐ) / \zh{上} (shàng) :} Après le nom de lieu si précision « dans / sur ». \zh{在公园里} (zài gōngyuán lǐ), \zh{在桌子上} (zài zhuōzi shàng).
\item \textbf{Modalités (\zh{应该}, \zh{可以}, \zh{要}, \zh{想}) :} Avant le verbe principal, après le lieu. \zh{我们明天应该在学校种树。} (Wǒmen míngtiān yīnggāi zài xuéxiào zhòng shù.)
\item \textbf{Classificateurs obligatoires :} \zh{一棵树} (yì kē shù), \zh{一个塑料袋} (yí ge sùliàodài), \zh{一辆公共汽车} (yí liàng gōnggòng qìchē). Oublier le classif. = erreur majeure (niveau HSK1/2).
\end{enumerate}
\end{methode}

\begin{exoresolu}[Remise en ordre, traduction, correction d'erreurs]
\textbf{Énoncé.} (a) Ordre : \zh{种树 / 我们 / 在 / 学校 / 里 / 明天 / 应该}. (b) Traduis : « Demain, nous devrions ramasser des déchets dans le parc. » (c) Corrige : \zh{我们捡垃圾在公园。} \zh{我有一个树。}
\tcblower
\textbf{Solution détaillée.}
\begin{enumerate}
\item \textbf{Ordre correct :} \zh{我们明天应该在学校里种树。} (Wǒmen míngtiān yīnggāi zài xuéxiào lǐ zhòng shù.) Ordre : Sujet (\zh{我们}) + Temps (\zh{明天}) + Modal (\zh{应该}) + Lieu (\zh{在学校里}) + Verbe (\zh{种树}).
\item \textbf{Traduction :} \zh{我们明天应该在公园捡垃圾。} (Wǒmen míngtiān yīnggāi zài gōngyuán jiǎn lājī.) \zh{捡垃圾} (jiǎn lājī) verbe + objet.
\item \textbf{Corrections :} \\
1. \zh{我们捡垃圾在公园。} $\rightarrow$ \zh{我们在公园捡垃圾。} (Lieu \zh{在公园} avant verbe). \\
2. \zh{我有一个树。} $\rightarrow$ \zh{我有一棵树。} (Classif. \zh{棵} pour arbres). Ou \zh{我有一棵树要种。} (J'ai un arbre à planter).
\end{enumerate}
\end{exoresolu}

\courssub{Écriture des caractères clés : Analyse et ordre des traits}

\begin{definition}[Caractères à maîtriser pour l'écrit et la reconnaissance (HSK 3)]
\begin{center}
\begin{tabularx}{\linewidth}{|X|X|X|X|X|}
\hline
\rowcolor{popblueL} \textbf{Caractère} & \textbf{Pinyin} & \textbf{Clé (Radical)} & \textbf{Nb traits} & \textbf{Ordre des traits (règles clés)} \\
\hline
环 & huán & \zh{王} (jade/roi) & 8 & Haut→bas : \zh{王} puis \zh{亏} (enfermement). Clé = jade/anneau $\rightarrow$ environnement (ce qui entoure). \\
\hline
境 & jìng & \zh{土} (terre) & 14 & Gauche→droite : \zh{土} + \zh{立} + \zh{一} + \zh{小}. Idée : limite de terre $\rightarrow$ frontière, milieu. \\
\hline
保 & bǎo & \zh{亻} (homme) & 9 & Clé \zh{亻} + \zh{呆} (phonétique). Homme qui veille $\rightarrow$ protéger, garder. \\
\hline
护 & hù & \zh{扌} (main) & 7 & Clé \zh{扌} + \zh{戌} (phonétique). Action de la main qui protège. \\
\hline
污 & wū & \zh{氵} (eau) & 6 & Clé \zh{氵} (3 traits) + \zh{于}. Eau sale $\rightarrow$ polluer, souiller. \\
\hline
染 & rǎn & \zh{木} (bois) & 9 & Haut→bas : \zh{木} + \zh{氵} (eau/teinture) + \zh{田}. Teindre le bois $\rightarrow$ teindre, contaminer. \\
\hline
垃 & lā & \zh{土} (terre) & 10 & \zh{土} + \zh{立} + \zh{刂}. Déchets sur la terre. \\
\hline
圾 & jī & \zh{土} (terre) & 11 & \zh{土} + \zh{及}. Deuxième partie de \zh{垃圾}. \\
\hline
\end{tabularx}
\end{center}
\textbf{Exercice d'écriture :} Écris chaque caractère 5 fois en respectant l'ordre des traits. Dicte-toi le pinyin + tons en écrivant.
\end{definition}

\courssub{Élément socioculturel : Cameroun — Chine, politiques vertes}

\begin{savaistu}[Gestion des déchets et reforestation : regards croisés]
\begin{itemize}
\item \textbf{Chine :} Depuis 2019, \zh{垃圾分类} (lājī fēnlèi « tri sélectif ») est obligatoire à Shanghai, puis étendu aux grandes villes (Pékin, Shenzhen). 4 poubelles : \zh{可回收物} (recyclables), \zh{有害垃圾} (dangereux), \zh{厨余垃圾} (organiques), \zh{其他垃圾} (autres). Grande campagne \zh{光盘行动} (Guāngpán xíngdòng « Assiette vide ») contre le gaspillage alimentaire. Reforestation massive : \zh{三北防护林} (Sānběi fánghùlín « Ceinture verte des Trois-Nord », « Grande Muraille Verte ») depuis 1978.
\item \textbf{Cameroun :} Interdiction des sacs plastiques non biodégradables (\zh{不可降解塑料袋} bù kě jiàngjiě sùliàodài) depuis 2014 (application inégale). Projet \zh{Sahel Vert} / \zh{Grande Muraille Verte} panafricaine (reboisement nord). Défis : gestion des ordures ménagères à Yaoundé/Douala (\zh{HYSACAM}), pollution plastique des cours d'eau (Wouri, Mfoundi).
\item \textbf{Point commun :} Les deux pays font face à l'urbanisation rapide et à la pollution plastique. La Chine mise sur la technologie (IA pour le tri, incinération) et la contrainte administrative ; le Cameroun sur la sensibilisation (\zh{École verte}, \zh{Journée de l'arbre} le 20 mars) et la coopération internationale.
\item \textbf{Pour ton oral :} Citer un fait précis (ex. \zh{上海垃圾分类很严格} « Le tri à Shanghai est très strict ») montre une culture générale appréciée du jury.
\end{itemize}
\end{savaistu}

\courssub{Activité de classe : Jeu de rôle et débat guidé (20 min)}

\begin{experience}[Simulation : « Comité Éco-Lycée »]
\textbf{Contexte :} Vous êtes le bureau du Club Environnement du lycée. Vous devez proposer 3 actions concrètes pour le trimestre.
\textbf{Rôles (groupes de 4) :}
\begin{enumerate}
\item \textbf{Président(e)} : Ouvre la séance, donne la parole, synthétise, vote.
\item \textbf{Secrétaire} : Note les propositions en chinois (mots-clés), lit le compte-rendu final.
\item \textbf{Deux membres} : Font des propositions, argumentent, s'accordent/désaccordent.
\end{enumerate}
\textbf{Déroulement (chronométré) :}
\begin{enumerate}
\item \textbf{2 min — Préparation individuelle :} Chacun écrit 1 proposition complète sur un papier : \zh{我认为我们应该……，因为……。} (Utilise vocabulaire fiche).
\item \textbf{8 min — Débat :} Tour de table. Président lance : \zh{同学们，我们要保护学校环境。谁有建议？} Membres répondent avec formules (Methode 2/3). Secrétaire note. Président relance : \zh{你同意吗？ / 还有别的想法吗？}
\item \textbf{5 min — Synthèse et vote :} Président : \zh{总之，我们决定……} (3 actions). Vote : \zh{同意的举手。}
\item \textbf{5 min — Restitution orale :} Chaque groupe désigne un rapporteur pour présenter les 3 actions décidées en 30 secondes devant la classe. \textbf{Critères :} Tons justes, phrases complètes, connecteurs \zh{因为/所以}, \zh{而且}, \zh{总之}, interaction visible.
\end{enumerate}
\textbf{Variante (niveau +) :} Imposer un désaccord simulé : un membre \emph{doit} contester une proposition avec \zh{我不太同意……因为……}.
\end{experience}

\courssub{Exercices d'entraînement (Devoir / Travail personnel)}

\begin{exercice}[Entraînement écrit et oral — Niveau Bac Blanc]
\textbf{Exercice 1 — Vocabulaire \& Classificateurs.} Complète avec le bon classificateur : \zh{个, 棵, 只, 辆, 个}.
(a) \zh{一\_\_\_树} \quad (b) \zh{一\_\_\_塑料袋} \quad (c) \zh{一\_\_\_公共汽车} \quad (d) \zh{一\_\_\_问题} \quad (e) \zh{一\_\_\_河}.

\textbf{Exercice 2 — Ordre des mots (Grammaire).} Remets dans l'ordre (caractères + pinyin).
(a) \zh{应该 / 我们 / 少用 / 塑料袋} \\
(b) \zh{在 / 公园 / 我们 / 捡 / 垃圾 / 明天} \\
(c) \zh{因为 / 空气 / 污染 / 严重 / 所以 / 我们 / 骑 / 自行车}

\textbf{Exercice 3 — Traduction (Version / Thème).}
(a) \zh{保护环境，人人有责。}\\
(b) \zh{我觉得雅温得的河里垃圾太多了。}\\
(c) « Nous devrions planter des arbres à l'école demain. »\\
(d) « À mon avis, l'air n'est pas bon parce qu'il y a trop de voitures. »\\
(e) « Peux-tu donner un exemple ? Je ne suis pas tout à fait d'accord. »

\textbf{Exercice 4 — Expression orale guidée (Préparation 10 min).}
Prépare un exposé de 1 minute sur l'un des sujets (fiche bristol autorisée, pas de texte lu) :
\begin{enumerate}
\item \zh{我的学校如何保护环境} (Comment mon école protège l'environnement).
\item \zh{少用塑料袋，从我做起} (Moins de sacs plastiques, commençons par moi).
\item \zh{骑自行车上学的好处} (Les avantages d'aller à l'école à vélo).
\end{enumerate}
Structure imposée : Salutation + Annonce + 2 arguments (\zh{因为……所以……}) + Conclusion slogan + Remerciement.
\end{exercice}

\courssub{Corrigés des exercices d'entraînement}

\begin{corrige}[Exercice 1 — Vocabulaire \& Classificateurs]
\begin{enumerate}
\item \zh{一棵树} (yì kē shù) — \zh{棵} pour arbres.
\item \zh{一个塑料袋} (yí ge sùliàodài) — \zh{个} (ou \zh{只} zhǐ pour sacs/conteneurs).
\item \zh{一辆公共汽车} (yí liàng gōnggòng qìchē) — \zh{辆} pour véhicules à roues.
\item \zh{一个问题} (yí ge wèntí) — \zh{个} général.
\item \zh{一条河} (yì tiáo hé) — \zh{条} (tiáo) pour choses longues (rivières, routes, pantalons). \emph{Piège : pas \zh{个} pour les fleuves !}
\end{enumerate}
\end{corrige}

\begin{corrige}[Exercice 2 — Ordre des mots]
\begin{enumerate}
\item \zh{我们应该少用塑料袋。} (Wǒmen yīnggāi shǎo yòng sùliàodài.) Sujet + Modal + Verbe + Objet.
\item \zh{我们明天在公园捡垃圾。} (Wǒmen míngtiān zài gōngyuán jiǎn lājī.) Sujet + Temps + Lieu + Verbe + Objet.
\item \zh{因为空气污染严重，所以我们骑自行车。} (Yīnwèi kōngqì wūrǎn yánzhòng, suǒyǐ wǒmen qí zìxíngchē.) Structure cause-effet complète. \zh{骑自行车} (verbe + objet).
\end{enumerate}
\end{corrige}

\begin{corrige}[Exercice 3 — Traduction]
\begin{enumerate}
\item \zh{保护环境，人人有责。} (Bǎohù huánjìng, rénrén yǒuzé.) « Protéger l'environnement, c'est la responsabilité de tous. »
\item \zh{我觉得雅温得的河里垃圾太多了。} (Wǒ juéde Yǎwēndé de hé lǐ lājī tài duō le.) « Je trouve qu'il y a trop de déchets dans les rivières de Yaoundé. »
\item \zh{我们明天应该在学校种树。} (Wǒmen míngtiān yīnggāi zài xuéxiào zhòng shù.) Ordre : Temps + Modal + Lieu + Verbe.
\item \zh{在我看来，空气不好，因为汽车太多了。} (Zài wǒ kànlái, kōngqì bù hǎo, yīnwèi qìchē tài duō le.) \zh{在我看来} en tête ; \zh{不太好} ou \zh{不好} ; \zh{太多了} excès.
\item \zh{你能举个例子吗？我不太同意。} (Nǐ néng jǔ ge lìzi ma? Wǒ bù tài tóngyì.) Formules de relance et désaccord poli.
\end{enumerate}
\end{corrige}

\begin{corrige}[Exercice 4 — Grille d'auto-évaluation pour l'exposé oral]
Utilise cette grille pour t'enregistrer (téléphone) et te corriger avant le passage en classe.
\begin{center}
\begin{tabularx}{\linewidth}{|X|X|X|}
\hline
\rowcolor{popblueL} \textbf{Critère} & \textbf{Objectif (Oui/Non)} & \textbf{Commentaire / Note} \\
\hline
Structure & 5 blocs respectés (Salut, Annonce, Arg1, Arg2, Conclusion+Merci) & \\
Grammaire & \zh{因为……所以……} utilisé correctement au moins 1 fois & \\
Vocabulaire & 5+ mots du thème employés (environnement, protéger, pollution, sac, arbre, air, responsabilité) & \\
Tons & Mots-clés (\zh{环境, 保护, 污染, 垃圾, 空气, 应该}) prononcés tons justes (vérifie sandhi \zh{不应该, 很好, 可以}) & \\
Ordre mots & Aucun complément de lieu/temps après le verbe & \\
Classificateurs & \zh{棵} pour arbre, \zh{辆} pour bus, \zh{条} pour rivière utilisés si nombres & \\
Interaction & Regard vers public, débit fluide, pas de lecture & \\
\hline
\end{tabularx}
\end{center}
\textbf{Objectif :} 7/7 « Oui » = Prêt pour le Bac. 5-6 = Réviser points faibles. $<$5 = Refaire la fiche méthode + s'enregistrer 3 fois.
\end{corrige}

\begin{attention}[Les 6 erreurs qui coûtent des points au Bac (Révisées)]
\begin{enumerate}
\item \textbf{Tons aplatis, inversés ou sandhi oubliés.} Dire \emph{huán jīng} (1\textsuperscript{er} ton) au lieu de \emph{huán jìng} (4\textsuperscript{e} ton) change le sens (\zh{镜} = miroir). Oublier le sandhi de \zh{很好} (\emph{hén hǎo}), \zh{可以} (\emph{ké yǐ}), \zh{不应该} (\emph{bú yīnggāi}), \zh{一棵} (\emph{yì kē}). \textbf{Action :} Surligne les tons sur ta fiche de prep ; répète les mots-clés en chaîne 5 min avant l'examen.
\item \textbf{Complément de lieu APRÈS le verbe (Calque français).} \zh{*我们种树在学校。} $\rightarrow$ \textbf{Faux.} Correction : \zh{我们在学校种树。} Règle : \zh{在} + Lieu \textbf{avant} Verbe. Toujours.
\item \textbf{Classificateur absent ou incorrect.} \zh{*一个树}, \zh{*一辆树}, \zh{*一个河}. \textbf{Action :} Apprends les triplets : \zh{一棵树, 一辆车, 一条河, 一个袋子}. \zh{个} = secours seulement.
\item \textbf{Confusion caractères proches (Clé radicale ignorée).} \zh{境} (environnement, clé \zh{土} terre) vs \zh{镜} (miroir, clé \zh{钅} métal). \zh{保} (protéger, clé \zh{亻} homme) vs \zh{报} (journal/rapport, clé \zh{扌} main). \zh{污} (sale, clé \zh{氵} eau) vs \zh{乌} (corbeau/noir). \textbf{Vérifie la clé : elle donne le sens.}
\item \textbf{Opinion brute sans formule fonctionnelle.} Répondre \zh{*不好！} ou \zh{*不同意！} sans \zh{我觉得/我认为/我不太同意}. \textbf{Jury attend :} Formule + Contenu + Justification (\zh{因为}).
\item \textbf{Absence d'interaction (Monologue).} Ne pas regarder le partenaire, ne pas poser de question (\zh{你怎么看？}, \zh{为什么？}), ne pas réagir (\zh{我同意}, \zh{有道理}). \textbf{Note interaction = 0.} \textbf{Action :} Termine CHAQUE réplique par une question ou une ouverture vers l'autre.
\end{enumerate}
\end{attention}

\newpage

\coursec{Exercices}

\courssub{Je vérifie mes connaissances}

\begin{exercice}[QCM : Vocabulaire de l'environnement]
Choisissez la bonne réponse parmi les trois propositions. Une seule réponse est correcte.
\begin{enumerate}
    \item Comment dit-on « protéger l'environnement » en chinois ?
    \begin{itemize}
        \item A. 破坏环境 (pòhuài huánjìng)
        \item B. \zh{保护环境} (bǎohù huánjìng)
        \item C. 污染环境 (wūrǎn huánjìng)
    \end{itemize}
    \item Quel mot désigne « les déchets / les ordures » ?
    \begin{itemize}
        \item A. \zh{垃圾} (lājī)
        \item B. 树木 (shùmù)
        \item C. 空气 (kōngqì)
    \end{itemize}
    \item Pour dire « économiser l'eau », on utilise le verbe :
    \begin{itemize}
        \item A. 浪费 (làngfèi)
        \item B. \zh{节约} (jiéyuē)
        \item C. 喝 (hē)
    \end{itemize}
    \item La phrase « L'air est très pollué » se traduit par :
    \begin{itemize}
        \item A. 空气很干净 (Kōngqì hěn gānjìng)
        \item B. \zh{空气很污染} (Kōngqì hěn wūrǎn)
        \item C. 空气很好 (Kōngqì hěn hǎo)
    \end{itemize}
\end{enumerate}
\end{exercice}

\begin{exercice}[Vrai ou Faux ? Justifiez en français]
Lisez les affirmations suivantes. Indiquez VRAI ou FAUX et justifiez votre réponse en vous basant sur le vocabulaire de la leçon.
\begin{enumerate}
    \item \zh{种树} (zhòng shù) signifie « couper les arbres ».
    \item \zh{塑料袋} (sùliào dài) signifie « sac en plastique ».
    \item \zh{回收} (huíshōu) signifie « jeter n'importe où ».
    \item \zh{噪音污染} (zàoyīn wūrǎn) désigne la pollution de l'eau.
\end{enumerate}
\end{exercice}

\begin{exercice}[Vocabulaire : Associez les mots aux définitions]
Reliez chaque mot chinois (caractères) à sa définition française correspondante.
\begin{center}
\begin{tabularx}{\linewidth}{|X|X|}
\hline
\rowcolor{popblueL} \textbf{Mot chinois (Caractères + Pinyin)} & \textbf{Définition} \\ \hline
1. \zh{环境} (huánjìng) & A. Action de trier les déchets pour les réutiliser \\ \hline
2. \zh{污染} (wūrǎn) & B. Milieu naturel ou cadre de vie \\ \hline
3. \zh{回收} (huíshōu) & C. Rendre quelque chose sale, impur, nocif \\ \hline
4. \zh{节约} (jiéyuē) & D. Utiliser avec parcimonie, ne pas gaspiller \\ \hline
5. \zh{垃圾分类} (lājī fēnlèi) & E. Trier les ordures par catégories \\ \hline
\end{tabularx}
\end{center}
\end{exercice}

\begin{exercice}[Phonétique : Associez le pinyin au caractère]
Écoutez (ou lisez) les pinyin suivants et écrivez le caractère chinois correspondant (choisi dans la liste ci-dessous).
\textbf{Liste :} \zh{树, 水, 说, 中, 重, 种}
\begin{enumerate}
    \item shuǐ $\rightarrow$ \underline{\hspace{3cm}}
    \item shuō $\rightarrow$ \underline{\hspace{3cm}}
    \item shù $\rightarrow$ \underline{\hspace{3cm}}
    \item zhòng $\rightarrow$ \underline{\hspace{3cm}}
    \item zhōng $\rightarrow$ \underline{\hspace{3cm}}
    \item zhǒng $\rightarrow$ \underline{\hspace{3cm}}
\end{enumerate}
\emph{Note : Attention aux tons ! Lisez ensuite les six mots à voix haute en respectant les tons.}
\end{exercice}

\courssub{Je m'entraîne}

\begin{exercice}[Traduction : Du français vers le chinois]
Traduisez les phrases suivantes en chinois (caractères + pinyin). Utilisez le vocabulaire de la leçon.
\begin{enumerate}
    \item Nous devons protéger l'environnement.
    \item Ne jetez pas les déchets par terre ! (Interdiction avec \zh{别} ou \zh{不要})
    \item L'air à Yaoundé est un peu pollué aujourd'hui.
    \item Je pense qu'il faut économiser l'eau et l'électricité.
    \item Est-ce que tu tries les ordures à la maison ?
\end{enumerate}
\end{exercice}

\begin{exercice}[Transformation : Exprimez l'avis et relancez]
À partir de la phrase de base, produisez les deux phrases demandées en utilisant les formules fonctionnelles vues en cours (\zh{我觉得}, \zh{你觉得呢？}, \zh{我同意}, \zh{我不同意}, \zh{所以}).
\begin{enumerate}
    \item \textbf{Base :} \zh{我们应该种树。} (Wǒmen yīnggāi zhòng shù.)
    \begin{itemize}
        \item A. Exprimez votre avis personnel : \underline{\hspace{8cm}}
        \item B. Demandez l'avis de l'interlocuteur : \underline{\hspace{8cm}}
    \end{itemize}
    \item \textbf{Base :} \zh{塑料袋污染环境。} (Sùliào dài wūrǎn huánjìng.)
    \begin{itemize}
        \item A. Exprimez votre accord : \underline{\hspace{8cm}}
        \item B. Donnez une conséquence avec \zh{所以} : \underline{\hspace{8cm}}
    \end{itemize}
    \item \textbf{Base :} \zh{开车上班很方便。} (Kāichē shàngbān hěn fāngbiàn.)
    \begin{itemize}
        \item A. Exprimez un désaccord poli (pollution) : \underline{\hspace{8cm}}
        \item B. Proposez une alternative avec \zh{可以} : \underline{\hspace{8cm}}
    \end{itemize}
\end{enumerate}
\end{exercice}

\begin{exercice}[Texte à trous : Complétez le dialogue]
Complétez le dialogue ci-dessous avec les mots ou formules manquants choisis dans la boîte.
\textbf{Mots à disposition :} \zh{我觉得, 我同意, 你呢？, 所以, 因为, 应该, 但是}
\begin{textebox}[Dialogue : Protéger l'environnement]
A : \zh{现在的空气污染很严重，\_\_\_\_\_\_我们要少开车。} (Xiànzài de kōngqì wūrǎn hěn yánzhòng, \_\_\_\_\_\_ wǒmen yào shǎo kāichē.)
B : \zh{\_\_\_\_\_\_。坐公共汽车比较好。} (\_\_\_\_\_\_. Zuò gōnggòng qìchē bǐjiào hǎo.)
A : \zh{\_\_\_\_\_\_，公共汽车有时候太挤了。} (\_\_\_\_\_\_, gōnggòng qìchē yǒushíhou tài jǐ le.)
B : \zh{\_\_\_\_\_\_我们可以骑自行车，\_\_\_\_\_\_？} (\_\_\_\_\_\_ wǒmen kěyǐ qí zìxíngchē, \_\_\_\_\_\_?)
A : \zh{好主意！\_\_\_\_\_\_既环保又健康。} (Hǎo zhǔyì! \_\_\_\_\_\_ jì huánbǎo yòu jiànkāng.)
\end{textebox}
\end{exercice}

\begin{exercice}[Remise en ordre : Reconstituez les phrases]
Remettez les mots dans l'ordre correct pour former des phrases grammaticales. Écrivez la phrase complète en caractères, pinyin et traduction.
\begin{enumerate}
    \item \zh{环境 / 保护 / 我们 / 应该 / 。}
    \item \zh{不 / 随地 / 扔 / 垃圾 / 要 / 。}
    \item \zh{你 / 怎么 / 看 / 这个 / 问题 / ？}
    \item \zh{虽然 / 累 / 但是 / 有意义 / 植树 / 。}
\end{enumerate}
\end{exercice}

\courssub{Je me prépare au Bac}

\begin{exercice}[Compréhension de l'oral / écrit : Texte original sur l'environnement au Cameroun]
Lisez attentivement le texte suivant (caractères + pinyin) et répondez aux questions en français.
\begin{textebox}[Mon quartier, mon environnement]
\zh{我叫保罗，住在喀麦隆的首都雅温得。我的家乡很美，但是现在环境有问题。} \\
\zh{第一，空气污染很严重，因为汽车太多了，工厂也排放废气。} \\
\zh{第二，塑料袋到处都是，河里也有很多垃圾，鱼都死掉了。} \\
\zh{我觉得我们每个人都应该保护环境。第一，我们应该少开车，多坐公交车或骑自行车。} \\
\zh{第二，出门购物要带自己的布袋子，不要用塑料袋。第三，要把垃圾分类，回收利用。} \\
\zh{如果大家都行动起来，我们的雅温得会变得更美、更干净。加油！}
\end{textebox}
\textbf{Questions de compréhension :}
\begin{enumerate}
    \item Où habite Paul ? Quelle est la situation environnementale générale de sa ville ?
    \item Citez les deux causes principales de la pollution de l'air mentionnées dans le texte.
    \item Quel est l'impact des déchets plastiques sur la rivière selon le texte ?
    \item Quelles sont les trois solutions concrètes que Paul propose aux citoyens ?
    \item Trouvez dans le texte le mot qui signifie « usine / fabrique » et celui qui signifie « émettre / rejeter (des gaz) ».
\end{enumerate}
\textbf{Questions de langue :}
\begin{enumerate}
    \setcounter{enumi}{5}
    \item Analysez la structure de la phrase : \zh{如果大家都行动起来，我们的雅温得会变得更美、更干净。} (Identifiez la conjonction de condition, le sujet, l'adverbe de futur, le verbe principal et les adjectifs).
    \item Traduisez en français : \zh{出门购物要带自己的布袋子，不要用塑料袋。}
    \item Transformez la phrase \zh{我们应该少开车} en une phrase donnant un avis personnel avec \zh{我觉得} et en une question relançant l'interlocuteur avec \zh{你觉得呢？}.
\end{enumerate}
\end{exercice}

\begin{exercice}[Thème et Version : Traduction spécialisée]
\textbf{Partie A : Thème (Français $\rightarrow$ Chinois)}
Traduisez le paragraphe suivant en chinois. Utilisez les structures \zh{应该 / 可以 / 因为...所以...} et le vocabulaire de l'environnement. (Environ 80-100 caractères).
\begin{quote}
« La protection de l'environnement est l'affaire de tous. À Douala, la pollution de l'air et des déchets plastiques est grave. Je pense que le gouvernement devrait interdire les sacs plastiques. Les citoyens peuvent utiliser des sacs en tissu. Si nous agissons ensemble, notre ville deviendra plus propre. »
\end{quote}

\textbf{Partie B : Version (Chinois $\rightarrow$ Français)}
Traduisez le texte suivant en français naturel et correct.
\begin{textebox}[Version]
\zh{现在，全世界都在关注气候变化。中国种了很多树，建设了很多“森林城市”。喀麦隆也有大森林，保护森林就是保护我们的未来。作为中学生，我们虽然力量小，但可以从小事做起：节约一张纸，少用一个塑料袋。保护环境，从我做起，从现在做起！}
\end{textebox}
\end{exercice}

\begin{exercice}[Expression écrite / orale : Sujet type Baccalauréat]
\textbf{Sujet :} \emph{« Dans votre ville (Yaoundé, Douala, Bafoussam, etc.), un problème environnemental vous préoccupe particulièrement (pollution de l'air, déchets plastiques, bruit, eau sale, coupe d'arbres...). »}

\textbf{Consignes :}
\begin{enumerate}
    \item Présentez ce problème en \textbf{5 phrases} minimum en chinois (caractères + pinyin).
    \item Utilisez obligatoirement : \zh{因为...所以...} (cause-conséquence), \zh{应该} (obligation morale), \zh{可以} (possibilité/solution).
    \item Proposez \textbf{deux solutions concrètes} adressées aux citoyens ou aux autorités.
    \item Soignez la prononciation (tons) et la cohérence logique. Nombre de caractères visé : 100-150.
\end{enumerate}
\textbf{Aide au plan :}
\begin{itemize}
    \item Phrase 1 : Présentation du problème (Où ? Quoi ? Gravité).
    \item Phrase 2 : Cause principale (\zh{因为...}).
    \item Phrase 3 : Conséquence (\zh{所以...}).
    \item Phrase 4 : Avis / Obligation (\zh{我觉得...应该...}).
    \item Phrase 5 : Solution 1 (\zh{可以...}).
    \item Phrase 6 : Solution 2 (\zh{也可以... / 还应该...}).
\end{itemize}
\end{exercice}

\newpage

\coursec{Corrigés des exercices}

\begin{corrige}[Exercice 1 — QCM : Vocabulaire de l'environnement]
\begin{enumerate}
    \item \textbf{Réponse B} : \zh{保护环境} (bǎohù huánjìng) = « protéger l'environnement ». La proposition A (\zh{破坏环境}) signifie « détruire l'environnement » et la C (\zh{污染环境}) « polluer l'environnement ». Piège classique : les trois verbes se construisent avec le même complément \zh{环境}, seul le verbe change le sens.
    \item \textbf{Réponse A} : \zh{垃圾} (lājī) = « déchets, ordures ». B \zh{树木} = « les arbres » ; C \zh{空气} = « l'air ».
    \item \textbf{Réponse B} : \zh{节约} (jiéyuē) = « économiser ». A \zh{浪费} = « gaspiller » (c'est l'antonyme, piège fréquent) ; C \zh{喝} = « boire ».
    \item \textbf{Réponse B} : \zh{空气很污染} (Kōngqì hěn wūrǎn) = « L'air est très pollué ». A signifie « L'air est très propre » et C « L'air est très bon ».
\end{enumerate}
\end{corrige}

\begin{attention}[Point de vigilance : l'adjectif-verbe chinois]
En chinois, \zh{污染} peut fonctionner comme verbe (« polluer ») ou comme adjectif (« pollué ») : \zh{空气很污染} est correct sans ajouter \zh{是}. Ne dites jamais \zh{空气是很污染} : devant un adjectif qualificatif, \zh{很} suffit et \zh{是} est de trop. Règle générale : Sujet + \zh{很} + adjectif (sans \zh{是}).
\end{attention}

\begin{corrige}[Exercice 2 — Vrai ou Faux ?]
\begin{enumerate}
    \item \textbf{FAUX.} \zh{种树} (zhòng shù) signifie « planter des arbres ». Le verbe \zh{种} (zhòng, 4\up{e} ton) signifie « planter, cultiver ». « Couper les arbres » se dit \zh{砍树} (kǎn shù). Attention à ne pas confondre avec \zh{种} (zhǒng, 3\up{e} ton) = « graine, espèce ».
    \item \textbf{VRAI.} \zh{塑料袋} (sùliào dài) = « sac en plastique » : \zh{塑料} = « plastique », \zh{袋} = « sac, sachet ».
    \item \textbf{FAUX.} \zh{回收} (huíshōu) signifie « recycler, récupérer » (\zh{回} = « retour » + \zh{收} = « recevoir, collecter »). « Jeter n'importe où » se dit \zh{随地扔垃圾} (suídì rēng lājī).
    \item \textbf{FAUX.} \zh{噪音污染} (zàoyīn wūrǎn) désigne la \emph{pollution sonore} (\zh{噪音} = « bruit, vacarme »). La pollution de l'eau se dit \zh{水污染} (shuǐ wūrǎn).
\end{enumerate}
\end{corrige}

\begin{corrige}[Exercice 3 — Associez les mots aux définitions]
\begin{center}
\begin{tabularx}{\linewidth}{|X|X|X|}
\hline
\rowcolor{popblueL} \textbf{Mot chinois} & \textbf{Réponse} & \textbf{Justification} \\ \hline
1. \zh{环境} (huánjìng) & B & \zh{环} = « entourer », \zh{境} = « territoire, milieu » : le milieu qui nous entoure. \\ \hline
2. \zh{污染} (wūrǎn) & C & \zh{污} = « sale », \zh{染} = « teindre, souiller » : rendre impur. \\ \hline
3. \zh{回收} (huíshōu) & A & \zh{回} = « retour » + \zh{收} = « collecter » : récupérer pour réutiliser. \\ \hline
4. \zh{节约} (jiéyuē) & D & \zh{节} = « économiser » + \zh{约} = « restreindre » : ne pas gaspiller. \\ \hline
5. \zh{垃圾分类} (lājī fēnlèi) & E & \zh{分类} = « trier par catégories » (\zh{分} = « diviser », \zh{类} = « catégorie »). \\ \hline
\end{tabularx}
\end{center}
\textbf{Correspondances :} 1-B, 2-C, 3-A, 4-D, 5-E.

\textbf{Méthode :} pour retenir le vocabulaire composé, décomposez toujours le mot en ses caractères et cherchez le sens de chacun : le chinois est une langue très « logique » où les mots composés racontent leur propre définition.
\end{corrige}

\begin{corrige}[Exercice 4 — Phonétique : Associez le pinyin au caractère]
\begin{enumerate}
    \item shuǐ $\rightarrow$ \zh{水} (l'eau) — 3\up{e} ton.
    \item shuō $\rightarrow$ \zh{说} (parler, dire) — 1\up{er} ton.
    \item shù $\rightarrow$ \zh{树} (arbre) — 4\up{e} ton.
    \item zhòng $\rightarrow$ \zh{种} (planter) — 4\up{e} ton.
    \item zhōng $\rightarrow$ \zh{中} (milieu, centre ; Chine) — 1\up{er} ton.
    \item zhǒng $\rightarrow$ \zh{种} (graine, espèce) — 3\up{e} ton.
\end{enumerate}
\end{corrige}

\begin{attention}[Piège majeur : les tons changent le sens]
Ce sont les mêmes syllabes mais des tons différents : \zh{水} shuǐ / \zh{说} shuō / \zh{树} shù, et \zh{中} zhōng / \zh{种} zhǒng / \zh{种} zhòng. Le caractère \zh{种} est un polyphone : il se lit \emph{zhǒng} (graine) ou \emph{zhòng} (planter). Dans \zh{种树} « planter des arbres », on dit obligatoirement zhòng shù. À l'oral, une erreur de ton peut transformer « planter des arbres » en « graines d'arbre » ! Entraînez-vous à voix haute en frappant le rythme des tons avec la main.
\end{attention}

\begin{corrige}[Exercice 5 — Traduction : Du français vers le chinois]
\begin{enumerate}
    \item \zh{我们应该保护环境。} \\
    \emph{Wǒmen yīnggāi bǎohù huánjìng.} \\
    Structure : Sujet + \zh{应该} (devoir) + Verbe + Complément. « Devoir » moral = \zh{应该}, pas \zh{必须} (trop fort, « il faut absolument »).
    \item \zh{别随地扔垃圾！} ou \zh{不要随地扔垃圾！} \\
    \emph{Bié suídì rēng lājī!} / \emph{Bú yào suídì rēng lājī!} \\
    L'interdiction se forme avec \zh{别} ou \zh{不要} + verbe. \zh{随地} = « n'importe où, par terre ». Ne pas utiliser \zh{不} seul pour une interdiction.
    \item \zh{今天雅温得的空气有点污染。} \\
    \emph{Jīntiān Yǎwēndé de kōngqì yǒudiǎn wūrǎn.} \\
    Le complément de temps \zh{今天} se place après le sujet (ou en tête de phrase). « Un peu » = \zh{有点儿} / \zh{有点}, placé avant l'adjectif. \zh{雅温得的空气} : possession avec \zh{的}.
    \item \zh{我觉得我们应该节约用水和电。} \\
    \emph{Wǒ juéde wǒmen yīnggāi jiéyuē yòng shuǐ hé diàn.} \\
    Avis personnel : \zh{我觉得} + phrase complète. « Économiser l'eau » = \zh{节约用水} (litt. « économiser dans l'usage de l'eau ») ; on peut aussi dire \zh{节约水电}.
    \item \zh{你在家垃圾分类吗？} ou \zh{你在家把垃圾分类吗？} \\
    \emph{Nǐ zài jiā lājī fēnlèi ma?} / \emph{Nǐ zài jiā bǎ lājī fēnlèi ma?} \\
    Question oui/non avec la particule finale \zh{吗}. Le lieu \zh{在家} se place avant le verbe. La construction avec \zh{把} est plus idiomatique : Sujet + \zh{把} + objet + verbe + complément.
\end{enumerate}
\end{corrige}

\begin{corrige}[Exercice 6 — Transformation : Exprimez l'avis et relancez]
\begin{enumerate}
    \item Base : \zh{我们应该种树。}
    \begin{itemize}
        \item A. \zh{我觉得我们应该种树。} \emph{Wǒ juéde wǒmen yīnggāi zhòng shù.} « Je pense que nous devrions planter des arbres. »
        \item B. \zh{我们应该种树，你觉得呢？} \emph{Wǒmen yīnggāi zhòng shù, nǐ juéde ne?} « Nous devrions planter des arbres, qu'en penses-tu ? » La relance \zh{你觉得呢？} se place en fin de phrase.
    \end{itemize}
    \item Base : \zh{塑料袋污染环境。}
    \begin{itemize}
        \item A. \zh{我同意，塑料袋污染环境。} \emph{Wǒ tóngyì, sùliào dài wūrǎn huánjìng.} « Je suis d'accord, les sacs plastiques polluent l'environnement. »
        \item B. \zh{塑料袋污染环境，所以我们应该少用。} \emph{Sùliào dài wūrǎn huánjìng, suǒyǐ wǒmen yīnggāi shǎo yòng.} « Les sacs plastiques polluent, donc nous devrions moins les utiliser. » \zh{所以} introduit la conséquence.
    \end{itemize}
    \item Base : \zh{开车上班很方便。}
    \begin{itemize}
        \item A. \zh{我不同意，开车会污染空气。} \emph{Wǒ bù tóngyì, kāichē huì wūrǎn kōngqì.} « Je ne suis pas d'accord, conduire pollue l'air. » Négation de \zh{同意} : \zh{不同意}.
        \item B. \zh{我们可以坐公共汽车或骑自行车上班。} \emph{Wǒmen kěyǐ zuò gōnggòng qìchē huò qí zìxíngchē shàngbān.} « Nous pouvons prendre le bus ou aller au travail à vélo. » \zh{可以} + verbe exprime la possibilité/la solution.
    \end{itemize}
\end{enumerate}
\end{corrige}

\begin{attention}[Point de vigilance : relancer et conclure]
\zh{呢} seul ne suffit pas pour relancer : il faut \zh{你觉得呢？} ou \zh{你呢？}. Et \zh{所以} ne s'emploie jamais en tête de texte sans cause exprimée avant : il faut une première proposition (éventuellement introduite par \zh{因为}). La paire correcte est \zh{因为...所以...}, mais en chinois on peut aussi n'utiliser que \zh{所以} si la cause vient d'être mentionnée.
\end{attention}

\begin{textebox}[Dialogue complété : Protéger l'environnement]
A : \zh{现在的空气污染很严重，所以我们要少开车。} (Xiànzài de kōngqì wūrǎn hěn yánzhòng, suǒyǐ wǒmen yào shǎo kāichē.) \\
B : \zh{我同意。坐公共汽车比较好。} (Wǒ tóngyì. Zuò gōnggòng qìchē bǐjiào hǎo.) \\
A : \zh{但是，公共汽车有时候太挤了。} (Dànshì, gōnggòng qìchē yǒushíhou tài jǐ le.) \\
B : \zh{我觉得我们可以骑自行车，你呢？} (Wǒ juéde wǒmen kěyǐ qí zìxíngchē, nǐ ne?) \\
A : \zh{好主意！因为既环保又健康。} (Hǎo zhǔyì! Yīnwèi jì huánbǎo yòu jiànkāng.)
\end{textebox}

\begin{corrige}[Exercice 7 — Texte à trous : Complétez le dialogue]
\textbf{Traduction du dialogue complété :} A : « La pollution de l'air est grave en ce moment, donc nous devons moins conduire. » B : « Je suis d'accord. Prendre le bus, c'est mieux. » A : « Mais le bus est parfois trop bondé. » B : « Je pense qu'on peut faire du vélo, et toi ? » A : « Bonne idée ! Parce que c'est à la fois écologique et bon pour la santé. »

\textbf{Justifications :}
\begin{itemize}
    \item Trou 1 : \zh{所以} — la première partie donne une cause (pollution grave), la seconde une conséquence (moins conduire).
    \item Trou 2 : \zh{我同意} — B approuve et ajoute une proposition (le bus).
    \item Trou 3 : \zh{但是} — A introduit une objection (le bus est bondé) : opposition.
    \item Trous 4 et 5 : \zh{我觉得} + \zh{你呢？} — B donne son avis puis relance l'interlocuteur.
    \item Trou 6 : \zh{因为} — A justifie son enthousiasme ; la structure \zh{既...又...} (jì... yòu...) signifie « à la fois... et... ».
\end{itemize}
\textbf{Méthode :} pour réussir un texte à trous, lisez d'abord tout le dialogue en global, repérez la logique de chaque échange (cause, accord, objection, proposition, conclusion), puis choisissez le mot-outil qui correspond à cette logique.
\end{corrige}

\begin{corrige}[Exercice 8 — Remise en ordre : Reconstituez les phrases]
\begin{enumerate}
    \item \zh{我们应该保护环境。} \\
    \emph{Wǒmen yīnggāi bǎohù huánjìng.} « Nous devons protéger l'environnement. » \\
    Ordre : Sujet (\zh{我们}) + modal (\zh{应该}) + verbe (\zh{保护}) + complément (\zh{环境}).
    \item \zh{不要随地扔垃圾。} \\
    \emph{Bú yào suídì rēng lājī.} « Ne jetez pas les ordures n'importe où. » \\
    Interdiction : \zh{不要} + complément de manière/lieu (\zh{随地}) + verbe (\zh{扔}) + objet (\zh{垃圾}).
    \item \zh{你怎么看这个问题？} \\
    \emph{Nǐ zěnme kàn zhège wèntí?} « Que penses-tu de ce problème ? » \\
    Structure : Sujet + \zh{怎么} (comment) + verbe (\zh{看} = « envisager, juger » ici) + complément (\zh{这个问题}). C'est une formule de relance très utile au débat, équivalente de \zh{你觉得呢？}.
    \item \zh{植树虽然累，但是有意义。} \\
    \emph{Zhíshù suīrán lèi, dànshì yǒu yìyì.} « Planter des arbres, bien que fatigant, a du sens. » \\
    Concession : \zh{虽然} + adjectif, \zh{但是} + proposition principale. En chinois, \zh{虽然} et \zh{但是} s'emploient \emph{ensemble} (contrairement au français où « bien que » et « mais » sont incompatibles).
\end{enumerate}
\end{corrige}

\begin{attention}[Erreur fréquente des francophones]
Ne traduisez pas « bien que... mais... » en choisissant l'un des deux : en chinois, la paire \zh{虽然...但是...} est obligatoire et correcte. À l'inverse, n'utilisez jamais \zh{因为...但是...} ensemble : cause et opposition ne se combinent pas.
\end{attention}

\begin{corrige}[Exercice 9 — Compréhension : Mon quartier, mon environnement]
\textbf{Questions de compréhension :}
\begin{enumerate}
    \item Paul habite à Yaoundé (\zh{雅温得}), la capitale du Cameroun. Sa ville est belle, mais l'environnement y a maintenant des problèmes : pollution de l'air grave et déchets (sacs plastiques) partout.
    \item Les deux causes de la pollution de l'air sont : le nombre excessif de voitures (\zh{汽车太多了}) et les usines qui rejettent des gaz d'échappement (\zh{工厂也排放废气}).
    \item Les sacs plastiques sont partout et la rivière est pleine de déchets ; les poissons en meurent (\zh{鱼都死掉了}).
    \item Les trois solutions proposées : (1) moins conduire, prendre le bus ou le vélo ; (2) apporter son propre sac en tissu pour faire les courses, ne pas utiliser de sacs plastiques ; (3) trier les déchets et les recycler.
    \item « Usine » = \zh{工厂} (gōngchǎng) ; « émettre, rejeter » = \zh{排放} (páifàng). \zh{废气} (fèiqì) signifie « gaz d'échappement, gaz résiduaires ».
\end{enumerate}
\textbf{Questions de langue :}
\begin{enumerate}
    \setcounter{enumi}{5}
    \item Analyse de \zh{如果大家都行动起来，我们的雅温得会变得更美、更干净。} : conjonction de condition \zh{如果} (« si ») ; sujet de la subordonnée \zh{大家} (« tout le monde ») ; sujet de la principale \zh{我们的雅温得} ; adverbe de futur \zh{会} ; verbe principal \zh{变得} (« devenir ») suivi du degré \zh{更} (« plus ») ; adjectifs \zh{更美、更干净} (« plus belle, plus propre »). Traduction : « Si tout le monde se met en action, notre Yaoundé deviendra plus belle et plus propre. »
    \item \zh{出门购物要带自己的布袋子，不要用塑料袋。} : « Quand on sort faire les courses, il faut apporter son propre sac en tissu, et ne pas utiliser de sacs en plastique. » (\zh{要} = « il faut » ; \zh{自己的} = « son propre ».)
    \item Avis personnel : \zh{我觉得我们应该少开车。} \emph{Wǒ juéde wǒmen yīnggāi shǎo kāichē.} « Je pense que nous devrions moins conduire. » Relance : \zh{我们应该少开车，你觉得呢？} \emph{Wǒmen yīnggāi shǎo kāichē, nǐ juéde ne?} « Nous devrions moins conduire, qu'en penses-tu ? »
\end{enumerate}
\textbf{Barème indicatif (sur 20) :} compréhension 10 pts (2 pts par question) ; questions de langue 10 pts (analyse 4 pts, traduction 3 pts, transformation 3 pts).
\end{corrige}

\begin{corrige}[Exercice 10 — Thème et Version]
\textbf{Partie A : Thème — proposition de traduction}

\zh{保护环境是每个人的事情。在杜阿拉，空气污染和塑料垃圾问题很严重。我觉得政府应该禁止塑料袋。市民可以用布袋。如果我们大家一起行动，我们的城市会变得更干净。}

\emph{Bǎohù huánjìng shì měi ge rén de shìqing. Zài Dù'ālā, kōngqì wūrǎn hé sùliào lājī wèntí hěn yánzhòng. Wǒ juéde zhèngfǔ yīnggāi jìnzhǐ sùliào dài. Shìmín kěyǐ yòng bùdài. Rúguǒ wǒmen dàjiā yìqǐ xíngdòng, wǒmen de chéngshì huì biàn de gèng gānjìng.}

Points de méthode : « l'affaire de tous » = \zh{每个人的事情} ; « interdire » = \zh{禁止} (jìnzhǐ) ; « agir ensemble » = \zh{一起行动} ; la condition \zh{如果...} appelle \zh{会} dans la principale ; « devenir + adjectif » = \zh{变得} + adjectif.

\textbf{Partie B : Version — proposition de traduction}

« Aujourd'hui, le monde entier s'intéresse au changement climatique. La Chine a planté beaucoup d'arbres et a construit de nombreuses "villes-forêts". Le Cameroun possède aussi de grandes forêts : protéger les forêts, c'est protéger notre avenir. En tant que lycéens, même si nos forces sont modestes, nous pouvons commencer par de petits gestes : économiser une feuille de papier, utiliser un sac plastique de moins. Protéger l'environnement commence par moi, et commence maintenant ! »

Points de méthode : \zh{关注} = « s'intéresser à, suivre de près » ; \zh{虽然力量小，但可以...} = concession \zh{虽然...但...} ; \zh{从小事做起} = « commencer par les petites choses » ; \zh{从我做起，从现在做起} = slogan « à partir de moi, à partir de maintenant ».

\textbf{Barème indicatif (sur 20) :} thème 10 pts (exactitude grammaticale 4, vocabulaire 4, caractères 2) ; version 10 pts (fidélité 5, qualité du français 3, vocabulaire 2).
\end{corrige}

\begin{corrige}[Exercice 11 — Expression écrite / orale : Sujet type Baccalauréat]
\textbf{Proposition de rédaction modèle (environ 120 caractères) :}

\zh{我住在杜阿拉。我们城市的塑料垃圾问题很严重。因为很多人随地扔塑料袋，所以街上和河里到处都是垃圾，很不卫生。我觉得每个市民都应该保护环境。第一，我们出门购物可以带自己的布袋，不用塑料袋。第二，政府也可以多放垃圾桶，还应该禁止乱扔垃圾。如果大家一起努力，杜阿拉会变得更干净、更美！}

\emph{Wǒ zhù zài Dù'ālā. Wǒmen chéngshì de sùliào lājī wèntí hěn yánzhòng. Yīnwèi hěn duō rén suídì rēng sùliào dài, suǒyǐ jiēshang hé hé lǐ dàochù dōu shì lājī, hěn bù wèishēng. Wǒ juéde měi ge shìmín dōu yīnggāi bǎohù huánjìng. Dì-yī, wǒmen chūmén gòuwù kěyǐ dài zìjǐ de bùdài, bú yòng sùliào dài. Dì-èr, zhèngfǔ yě kěyǐ duō fàng lājītǒng, hái yīnggāi jìnzhǐ luàn rēng lājī. Rúguǒ dàjiā yìqǐ nǔlì, Dù'ālā huì biàn de gèng gānjìng, gèng měi!}

\textbf{Traduction :} « J'habite à Douala. Le problème des déchets plastiques de notre ville est très grave. Parce que beaucoup de gens jettent leurs sacs plastiques n'importe où, les rues et les rivières sont pleines d'ordures, ce qui est très insalubre. Je pense que chaque citoyen doit protéger l'environnement. Premièrement, quand nous sortons faire les courses, nous pouvons apporter notre propre sac en tissu au lieu d'utiliser des sacs plastiques. Deuxièmement, le gouvernement peut aussi installer plus de poubelles et devrait interdire de jeter des déchets n'importe où. Si tout le monde fait des efforts ensemble, Douala deviendra plus propre et plus belle ! »

\textbf{Vérification des contraintes :}
\begin{itemize}
    \item \zh{因为...所以...} : présent (phrase 3).
    \item \zh{应该} : présent (\zh{每个市民都应该保护环境} ; \zh{还应该禁止}).
    \item \zh{可以} : présent (\zh{可以带自己的布袋} ; \zh{也可以多放垃圾桶}).
    \item Deux solutions concrètes : sac en tissu (citoyens) + poubelles et interdiction (autorités).
    \item Plan respecté : problème $\rightarrow$ cause $\rightarrow$ conséquence $\rightarrow$ avis $\rightarrow$ solutions $\rightarrow$ conclusion avec \zh{如果...会...}.
\end{itemize}

\textbf{Barème indicatif (sur 20) :} respect des consignes et structures obligatoires 5 pts ; correction grammaticale (ordre des mots, modaux, tons à l'oral) 5 pts ; richesse du vocabulaire de l'environnement 4 pts ; cohérence logique du plan 3 pts ; caractères corrects / prononciation soignée 3 pts.
\end{corrige}

\begin{attention}[Points de vigilance pour l'oral]
\begin{itemize}
    \item Prononcez bien les tons : \zh{垃圾} lājī (1-1), \zh{环境} huánjìng (2-4), \zh{应该} yīnggāi (1-1), \zh{塑料袋} sùliào dài (4-4-4).
    \item \zh{不} devient bú devant un 4\up{e} ton : \zh{不用} bú yòng, \zh{不要} bú yào.
    \item N'oubliez pas \zh{的} dans \zh{自己的布袋} et \zh{我们的城市}.
    \item Évitez le calque du français « très grave » avec \zh{非常} systématique : \zh{很严重} est naturel et suffisant.
    \item À l'oral, articulez les modaux : \zh{应该} (devoir moral), \zh{可以} (possibilité), \zh{要} (nécessité) ne sont pas interchangeables.
\end{itemize}
\end{attention}

\newpage

\coursec{Fiche bilan}

\courssub{Fiche bilan — Leçon 17 : Expression orale — Protection de l’environnement}

\begin{popfigure}
\begin{center}\tcbox[colback=popink,colframe=popink,coltext=white,arc=9pt,boxsep=4pt,left=6pt,right=6pt]{\bfseries\small Protection de l’environnement \zh{环境保护} (Huánjìng bǎohù)}\end{center}
\vspace{-2pt}
\begin{tcbraster}[raster columns=3,raster equal height=rows,raster column skip=6pt,raster row skip=6pt,enhanced,blankest]
\begin{tcolorbox}[enhanced,colback=popgreen,colframe=popgreen,coltext=white,boxrule=0.9pt,arc=3pt,left=4pt,right=4pt,top=3pt,bottom=3pt,boxsep=1pt,halign=left]\bfseries\scriptsize Vocabulaire clé \zh{重点词汇}\end{tcolorbox}
\begin{tcolorbox}[enhanced,colback=poporange,colframe=poporange,coltext=white,boxrule=0.9pt,arc=3pt,left=4pt,right=4pt,top=3pt,bottom=3pt,boxsep=1pt,halign=left]\bfseries\scriptsize Dialogue 1 : Donner son avis \zh{发表观点}\end{tcolorbox}
\begin{tcolorbox}[enhanced,colback=poppurple,colframe=poppurple,coltext=white,boxrule=0.9pt,arc=3pt,left=4pt,right=4pt,top=3pt,bottom=3pt,boxsep=1pt,halign=left]\bfseries\scriptsize Formules fonctionnelles \zh{功能句式}\end{tcolorbox}
\begin{tcolorbox}[enhanced,colback=poppink,colframe=poppink,coltext=white,boxrule=0.9pt,arc=3pt,left=4pt,right=4pt,top=3pt,bottom=3pt,boxsep=1pt,halign=left]\bfseries\scriptsize Dialogue 2 : Mini-débat \zh{小型辩论}\end{tcolorbox}
\begin{tcolorbox}[enhanced,colback=popgold,colframe=popgold,coltext=white,boxrule=0.9pt,arc=3pt,left=4pt,right=4pt,top=3pt,bottom=3pt,boxsep=1pt,halign=left]\bfseries\scriptsize Phonétique : Tons du thème \zh{重点声调}\end{tcolorbox}
\begin{tcolorbox}[enhanced,colback=popdark,colframe=popdark,coltext=white,boxrule=0.9pt,arc=3pt,left=4pt,right=4pt,top=3pt,bottom=3pt,boxsep=1pt,halign=left]\bfseries\scriptsize Jeux de rôle \& Débat \zh{角色扮演与辩论}\end{tcolorbox}
\begin{tcolorbox}[enhanced,colback=popmuted,colframe=popmuted,coltext=white,boxrule=0.9pt,arc=3pt,left=4pt,right=4pt,top=3pt,bottom=3pt,boxsep=1pt,halign=left]\bfseries\scriptsize Culture Cameroun / Chine \zh{中喀文化对比}\end{tcolorbox}
\end{tcbraster}

\legende{Carte mentale de la leçon 17 — Expression orale : protection de l’environnement}
\end{popfigure}

\begin{bilan}

\textbf{1. Lexique clé (HSK 2–3) — \zh{重点词汇}}\par
\begin{center}
\begin{tabularx}{\linewidth}{|X|X|X|X|}
\hline
\rowcolor{popblueL} \textbf{Caractères} & \textbf{Pinyin} & \textbf{Nature} & \textbf{Traduction} \\
\hline
环境 & huánjìng & n. & environnement \\
保护 & bǎohù & v./n. & protéger / protection \\
污染 & wūrǎn & v./n. & polluer / pollution \\
垃圾 & lājī & n. & déchets, ordures \\
分类 & fēnlèi & v./n. & trier / classification (tri sélectif) \\
回收 & huíshōu & v. & recycler, récupérer \\
节约 & jiéyuē & v. & économiser (ressources) \\
浪费 & làngfèi & v. & gaspiller \\
气候变化 & qìhòu biànhuà & n. & changement climatique \\
植树 & zhíshù & v.o. & planter des arbres \\
可持续发展 & kěchíxù fāzhǎn & n. & développement durable \\
低碳生活 & dītàn shēnghuó & n. & vie bas carbone \\
\hline
\end{tabularx}
\end{center}

\textbf{2. Structures grammaticales \& Formules fonctionnelles — \zh{语法与功能句式}}\par
\begin{center}
\begin{tabularx}{\linewidth}{|X|X|X|}
\hline
\rowcolor{popblueL} \textbf{Structure / Fonction} & \textbf{Exemple (caractères + pinyin)} & \textbf{Traduction} \\
\hline
\zh{我觉得……} (Wǒ juéde…) — Exprimer son avis & \zh{我觉得保护环境很重要。} (Wǒ juéde bǎohù huánjìng hěn zhòngyào.) & Je pense que protéger l'environnement est important. \\
\zh{我认为……} (Wǒ rènwéi…) — Opinion plus formelle & \zh{我认为大家应该垃圾分类。} (Wǒ rènwéi dàjiā yīnggāi lājī fēnlèi.) & Je pense que tout le monde devrait trier les déchets. \\
\zh{一方面……，另一方面……} — Argumentation équilibrée & \zh{一方面要发展经济，另一方面要保护环境。} (Yīfāngmiàn yào fāzhǎn jīngjì, lìngyīfāngmiàn yào bǎohù huánjìng.) & D'un côté il faut développer l'économie, de l'autre protéger l'environnement. \\
\zh{不仅……，而且……} — Progression & \zh{他不仅自己分类，而且还宣传环保。} (Tā bùjǐn zìjǐ fēnlèi, érqiě hái xuānchuán huánbǎo.) & Non seulement il trie lui-même, mais en plus il fait de la pub pour l'écolo. \\
\zh{如果……，就……} (Rúguǒ… jiù…) — Hypothèse & \zh{如果人人都节约用水，水资源就不会短缺。} (Rúguǒ rénrén dōu jiéyuē yòngshuǐ, shuǐ zīyuán jiù bù huì duǎnquē.) & Si tout le monde économise l'eau, les ressources en eau ne manqueront pas. \\
\zh{建议 / 呼吁} (Jiànyì / Hūyù) — Proposition / Appel & \zh{我建议大家少用塑料袋。} (Wǒ jiànyì dàjiā shǎo yòng sùliào dài.) & Je suggère que tout le monde utilise moins de sacs plastiques. \\
\zh{同意 / 不同意} (Tóngyì / Bù tóngyì) — Accord / Désaccord & \zh{我同意你的看法。 / 我不同意，因为……} (Wǒ tóngyì nǐ de kànfǎ. / Wǒ bù tóngyì, yīnwèi…) & Je suis d'accord avec ton avis. / Je ne suis pas d'accord, car… \\
\zh{请问……} (Qǐngwèn…) — Relancer poliment & \zh{请问你怎么看待气候变化？} (Qǐngwèn nǐ zěnme kàndài qìhòu biànhuà?) & Comment voyez-vous le changement climatique ? \\
\hline
\end{tabularx}
\end{center}

\textbf{3. Méthodes express — \zh{口语表达技巧}}\par
\begin{itemize}
    \item \textbf{Préparer un exposé (2 min) :} 1. Accroche (\zh{大家好，今天我谈谈……}) $\rightarrow$ 2. Problème (\zh{现在……很严重}) $\rightarrow$ 3. Causes / Exemples (\zh{因为……比如……}) $\rightarrow$ 4. Solutions (\zh{我们应该……}) $\rightarrow$ 5. Conclusion / Appel (\zh{让我们一起行动吧！}).
    \item \textbf{Participer à un débat :} Écouter $\rightarrow$ Noter mots-clés $\rightarrow$ Réagir (\zh{我赞成 / 我反对，理由是……}) $\rightarrow$ Rebondir (\zh{此外，我想补充……}).
    \item \textbf{Tons à surveiller :} \zh{环境} (2-3), \zh{污染} (1-3), \zh{分类} (1-4), \zh{回收} (2-1), \zh{节约} (2-1), \zh{发展} (1-3). Attention aux tons 3 consécutifs (sandhi) : \zh{保护环境} $\rightarrow$ \zh{báohù huánjìng}.
\end{itemize}

\end{bilan}

\begin{aretenir}[Checklist avant le Bac — Leçon 17]
$\square$ Je connais le vocabulaire de base de l'environnement (\zh{环境、污染、垃圾分类、回收、节约、气候变化}) et je les écris en caractères.\par
$\square$ Je sais présenter un exposé oral structuré de 1 à 2 minutes en chinois (introduction, arguments, conclusion).\par
$\square$ J'utilise correctement \zh{我觉得 / 我认为} pour donner mon avis et \zh{建议 / 呼吁} pour proposer.\par
$\square$ Je maîtrise la structure \zh{一方面……，另一方面……} pour nuancer mon propos.\par
$\square$ Je sais exprimer l'accord (\zh{我同意}) et le désaccord (\zh{我不同意，因为……}) de façon polie.\par
$\square$ Je prononce correctement les tons des mots-clés (attention aux sandhi des tons 3).\par
$\square$ Je sais relancer mon interlocuteur avec \zh{请问……} ou \zh{你怎么看？}.\par
$\square$ Je connais au moins 2 différences concrètes entre le Cameroun et la Chine sur la gestion des déchets / énergie verte.\par
$\square$ Je suis capable de participer à un jeu de rôle « débat en classe » sans lire mes notes.\par
$\square$ J'écris correctement les caractères \zh{护、污、类、收、约、废} (ordre des traits, radicaux).\par
$\square$ Je traduis sans erreur : « Protégeons l'environnement, c'est protéger notre avenir » $\rightarrow$ \zh{保护环境就是保护我们的未来。}
\end{aretenir}

\findecours

\end{document}
